\quellseite{120}
tröstliche Wort, das ihn aufrichtet und erhält, sonst würde er durch des Gesetzes Wort haben müssen vergehen. Das sind die Worte, darinne die zwei Predigten begriffen sind; eine, die ihn zu einem Sünder macht und tödtet; die andere, dadurch er fromm und lebendig wird. Nun lasset uns hören, wie er sich darein schickt und dem Wort folgt.

9. Zum ersten, siehe, wie der Glaube des heiligen Patriarchen gestanden ist. Also haben wir gesagt: daß der Glaube eine gewisse Zuversicht ist der Dinge, die man nicht siehet [Hebr. 11, 1.], aber doch verheißen sind, dadurch man gewartet deß, das da wohl scheint, als werde es nimmermehr kommen. Darum ist es ein blindes und doch ein helles Erkenntniß, ein Licht mitten im Finstern, also, daß er auch solch Gut hat, das er nicht hat, das ist, daß er nicht sieht, und jedermann sagen muß, er habe nichts. Als, daß ich es mit einem groben Exempel erkläre: Wenn sich ein armer Mann in ehelichen Stand geben soll, und der Glaube nicht da ist, steht das Herz also: Was thue ich? soll ich ein Weib nehmen, und habe nichts, damit ich mich ernähre? Weil er nun nicht sieht, wo das Gut liegt, daß er sich ernähre, will er nicht hinan. Ist aber Glaube da, so spricht er: Ich will ehelich werden in GOttes Namen; wo das Gut ist, das ich haben soll, oder wo ich es nehmen soll, weiß ich nicht, da sehe ich nicht nach, lasse mir an dem genügen, daß Christus [Matth. 6, 33.] sagt: „Trachtet am ersten nach dem Reich GOttes und nach seiner Gerechtigkeit, so wird euch solches alles zufallen“; das fasse ich, und tröste mich deß, eben als hätte ich, was ich haben sollte,\footnote{Erlanger: soll.} und habe doch nichts.

10. Siehe, das ist des Glaubens Art, daß er mit den Gütern handelt, die er nicht sieht noch fühlt, und geht gerade damit um, als habe er sie in Händen; hat keinen andern Trost, denn daß er weiß, daß GOtt nicht leugt noch trügt. Solches thut er in allerlei andern Stücken auch; als, wenn ich sterben soll, und der Tod mir unter Augen tritt, da muß ich davon, und weiß nicht, wo ich den Fuß im ersten Tritte hinsetzen soll. Ist nun der Unglaube da, der zappelt, zagt, und spricht: Wo will ich nun hin? Wer weiß, wo ich bleibe? Er will immerdar sehen und fühlen, wo er bleiben soll; es wird aber nichts daraus, darum muß er verzweifeln. Der Glaube aber denkt also: Ich weiß nicht, wohin ich fahre, aus muß ich, sehe und fühle nichts, ich will mich aber dem befehlen, der da gesagt hat im 55. Psalm, V. 23.: „Wirf deine Sorge und Anliegen auf den HErrn, der wird dich versorgen“; darauf fahre ich dahin, denn ich weiß, er wird nicht lügen. So hat er das Leben, und sieht es nicht, ja, er sieht nichts denn das Widerspiel. Wovon weiß er es denn? Allein davon, daß GOtt zugesagt hat, er wolle die, so ihm trauen, nicht lassen aus seiner Hand fallen. Solcher Exempel sollte man überall in der Schrift fleißig wahrnehmen.

11. Also ging es dem Patriarchen Noah (wie [Cap. 6, 13.] gehört), daß GOtt sagte, er wolle die Welt ersäufen und vertilgen. Und es gingen hundert Jahre hinweg, er sahe nichts, ja, es schien, als sollte nichts daraus werden; darum hielten es die Leute für Lügentheidung; er aber sahe es hundert Jahre zuvor, ehe es kam, und handelte so gewiß darauf, als sollte es morgen des Tages kommen. So scharfe Augen hat der Glaube, daß er im Dunkeln sehen kann, da doch nichts überall scheinet, sieht, da nichts zu sehen ist, fühlt, da nichts zu fühlen ist. Also glauben wir auch an den HErrn Christum, daß er droben sitzt zur rechten Hand des allmächtigen Vaters, und regiert also, daß er alle Creaturen in Händen hat, und alles in uns wirkt. Das sehen wir nicht, fühlen es auch nicht; noch sieht das Herz durch den Glauben so gewiß, als wenn es mit Augen sähe. Wenn nun der Mensch sterben soll, so ist eben derselbe Glaube da, und weiß gewiß, daß er ihm in den Schooß fährt.

12. Solche Art des Glaubens preist hier die Schrift in dem heiligen Altvater Abram. Er war in dem Lande, und bei seinen Freunden bekannt, sollte da erben sein väterlich Gut; da war noch kein Glaube, denn es war alles vorhanden, was er haben wollte, fühlte keinen Mangel oder Noth. Da schickt GOtt zu ihm das Wort, reißt ihn aus dem allen, was er vor Augen hat, und weß er sich trösten möchte, spricht: „Gehe heraus.“ Wovon? „Aus deinem Vaterlande, und von deiner Freundschaft, und aus deines Vaters Hause“, das ist, laß fahren, und verzeihe dich alles, das du hier siehst und genießen kannst. Wohin? „In ein Land, das ich dir zeigen will.“ Da wird nichts genannt; [es] ist eine hohe Versuchung und Kampf des
\quellseite{121}
Glaubens, daß er aus dem Lande soll fahren, und weiß nicht wohin, geht dahin, wo ihn der Wind hinwehet; item, weiß nicht, ob er zu Freunden oder Feinden kommt, und muß sich erwägen, er komme am ersten zu Feinden. Meinst du, daß das ein kleiner Streit gewesen sei?

13. Denn das Herz hat es müssen also fühlen: Hier reißest du mich von allen Bekannten, und führest mich dahin, vielleicht mitten unter die ärgsten Feinde, da man mich gar nicht wird können leiden; machst mich gar zu einem armen Bettler, daß ich keine Stätte behalte, da ich sicher bleiben könne. Ja, er hat sich müssen erwägen mit seinem Weibe, und wer mit ihm gezogen ist, daß sie müßten sterben, und in die Fahrt geben Ehre, Gut, Leib und Leben; läßt fahren gewiß Gut, Freude und Sicherheit, und gibt sich in gewisse Feinde, Armuth und Tod. Was thut er? Nimmt nicht mehr denn das Wort, daß Gott sagt: „Ich will dich zum großen Volke machen und will dich segnen.“ Ja, hätte der Unglaube gesagt, du sagst viel von segnen, und ich sehe nur das Widerpiel; hier habe ich ein unfruchtbar Weib [1 Mos. 11, 30.], so wirfst du mich dazu in ein fremd Land und Volk; heißt das gesegnet? Ja, vielmehr in alles Unglück gejaget. Noch folget er denn bloßen Worte; als sollte er sagen: Du hast geredet, du wollest mich segnen; wenn mich gleich jedermann verflucht, so soll mir’s nicht schaden; darum will ich es frisch darauf wagen. So geht er in solchem Glauben in das Finstere hinein, weiß nicht, wo er hingehet, noch findet er das rechte Land; wie wir hören werden.

14. Solches preiset nun die Epistel an die Hebräer [Cap. 11, 8.]: „Durch den Glauben ward gehorsam“, der da genannt ist „Abraham, auszugehen in das Land, das er sollte zum Erbe empfangen, und ging aus, und wußte nicht, wo er hinkäme.“ Das ist nun um unserwillen geschrieben, daß wir auch lernen an seinem Worte hangen, so er uns zusagt, daß er Leib, Leben, und auch die Seele genugsam versorgen und erhalten will, ob es gleich nicht vor Augen, noch vorhanden ist. Hebe nur an, und nimm dich sein an, als hättest du es bereits. Glaubst du, so hast du es. Der Glaube kann dir nicht fehlen noch trügen, ob es schon sich so ansehen läßt, als wollten alle Dinge zu Boden gehen. Komm es nicht so bald, wie du siehest, daß GOtt Abram aufgezogen hat, so laß nicht ab, denn, daß er verzeucht, thut er, den Glauben zu stärken. Folgt nun weiter:

15. Da siehe, wie GOtt den Glauben versucht hat. Denn es ist nicht ein schlecht Exempel, weil es die Schrift fleißig anzeigt, welche diesem Mann allein die Ehre gibt, daß sie ihn heißt einen Vater des Glaubens [Röm. 4, 11. Gal. 3, 9.], und Christus selbst [Luc. 16, 22.] Abrahams Glauben den Schooß Abrahä nennt; daß auch GOtt selbst viel davon hält. Er ist fünf und siebenzig Jahre alt, sagt der Text, als er auszog; von denen hat er gezählt, bis daß er hundert Jahre alt ward, da überkam er erst einen Sohn. Da werden wir sehen, wie er in dem Glauben so lange versucht wird, daß er hundertmal hätte mögen verzweifeln. GOtt hat ihm verheißen, seinen Samen zu mehren, da wartet und hofft er alle Jahre. Es währte fünf, zehn, fünfzehn, zwanzig Jahre; [es] wollte nichts daraus werden. O! wie hat er diese Zeitlang müssen fechten mit dem Unglauben. Da er nun hundert Jahre alt ward, war er fast verzweifelt, denn es der Natur nach unmöglich war; dennoch ließ er nicht von dem Worte und Verheißung GOttes, daß wohl St. Paulus [Röm. 4, 18.] sagt: „Er glaubte“ wider die Hoffnung „auf die Hoffnung.“

16. So lerne nun, daß nicht genug ist, daß man anhebt zu glauben, sondern muß auch fortfahren und beharren, und sich nicht davon reißen lassen. Denn daß GOtt den Glauben versucht, thut er nur darum, wie gesagt, daß der Glaube kräftig und fest werde, und der alte Adam und Unglaube gewürgt werde, der nur am Gegenwärtigen hangt, daß er sieht und fühlt, auf daß er uns von dem Gegenwärtigen abreiße, daß wir allein an dem Worte haften, so haben wir genug an Leib und Seele.

17. Also hat er gelernt und geübt den Spruch
\quellseite{122}
im fünften Buche Mosis [Cap. 8, 3.]: „Der Mensch lebet nicht allein von dem Brode, sondern von einem jeglichen Worte, das durch den Mund Gottes gehet“ [Matth. 4, 4.], so lebt er von dem bloßen Worte, daß er daran hangt, und darinne bleibt; so muß er gesegnet werden mit allen Gütern. Das sei das Exempel, daran wir lernen, wie ferne wir noch vom Glauben sind.

18. Daß aber GOtt zuvor in der Zusagung anzeucht, und spricht: „Ich will segnen, die dich segnen, und verfluchen, die dich verfluchen“, da legt er ihm das heilige Kreuz auf. Das ist das andere Stück, das wir zu lernen haben aus der Väter Legenden. Denn das sind die Hauptstücke des christlichen Lebens, Glaube, Hoffnung, Liebe, und das Kreuz. Denn wo GOttes Wort, das heilige Evangelium, ist, da soll ja das heilige Kreuz nicht außen bleiben. Denn Paulus gibt ihm [1 Cor. 1, 18.] den Namen, daß er es heißt ein Wort des Kreuzes, das ist, es kostet das Kreuz, und muß das Kreuz tragen. Denn die Welt und das Fleisch ist so gesinnt, daß sie das Evangelium nicht mag leiden, denn sie will nur an dem hangen, das sie vor Augen sieht, will Ehre und Guts genug haben und der Sache gewiß sein. Wenn nun das Evangelium predigt, es sei nicht recht, so kann sie es nicht hören; denn es predigt nur wider sie. Darum sagt es GOtt dem Abram zuvor, daß er verflucht wird werden, auf daß er sich solches erwäge, und wisse, daß es sein muß, lasse sich solches nicht anfechten; und tröstet ihn daneben gleich mit denselben Worten, als sollte er sagen: Sei getrost, und halte fest, du sollst sehen, wer dich verflucht, den will ich verfluchen, und wer dich segnet, den will ich segnen.

19. Aber das geht auch im Glauben daher. Es ist nicht alles beschrieben, wie es gegangen ist, aber hernach, an einem andern Orte, wohl angezeigt, und hier auch, daß er von vielen verflucht ist worden, verdammt, verhöhnt und verlästert, er war ein Fremdling im Lande unter seinen Feinden, und des Lebens nicht sicher, mußte seines Weibes auch fürchten [Vers 10—12.]. So tröstet ihn nun der Heilige Geist; schweigt aber stille, und lässet ihn\footnote{Erlanger: ihnen.} gleichwohl lästern, thut, als sähe er es nicht. Darum ist das ein Segen und Fluchen GOttes im Geiste, als sollte er sagen: Die dich verfluchen, sind wohl vor der Welt gesegnet, und wiederum; aber schweig stille, und friß in dich: siehe auf mich, wie ich dich führe im Geiste, daß sie vor mir verflucht sind.

20. Also weist er ihn immer auf das Wort. Da mußte er nun leben in der Feinde Augen, Leib und Leben wagen, und was er hatte, nehmen lassen; daß GOtt dazu sich stellt, als höre und sehe er es nicht, wie man ihn lästert und schändet. Das ist groß; aber noch größer ist es, daß er ziehen muß mit Weib und großem Gesinde; denn wo er allein wäre gewesen, hätte er viel besser fortkommen mögen. [Es] ist wohl zu denken, sie sind Herren gewesen im Lande, wie oft man ihm Unrecht und Gewalt gethan hat, und das Seine entzogen.

21. Das ist uns alles zu einem hohen Exempel vorgestellt, dergleichen man nicht viel findet unter den Heiligen. Denn GOtt hat selbst geordnet, seine Legende mit allem Fleiß zu beschreiben, daß ihn die ganze Welt sollte heißen einen Vater des Glaubens, davon jedermann lerne, wie der Glaube geschickt muß sein, und was er für Art habe, nämlich, daß er hindurch geht zwischen Leben und Tod, Ehre und Schande, Gut und Armuth; daß er könne zu GOtt sagen: Du bist mächtig, kannst mich im Tode lebendig behalten; weil du genug hast und lebst, so muß ich auch leben und genug haben. Denn wiewohl er das Widerspiel sieht, richtet er sich doch immer nach der Verheißung.

22. Auf daß wir aber diese Verheißung ein wenig baß ansehen und erklären, habe ich also gesagt, daß in diesen Worten: „In dir sollen gesegnet werden alle Geschlechte auf Erden“, eigentlich Christus verheißen ist. Wiewohl die Juden, welchen eigentlich Christus verheißen ist,\footnote{„welchen eigentlich Christus verheißen ist“ fehlt in der Wittenberger.} den Spruch dahin ziehen, wie sie andere Schrift auch verdrehen und schänden, daß es so viel geredet sei: „In dir sollen gesegnet werden alle Geschlechte auf Erden“, das ist, durch dich sollen alle Geschlechte gelobt werden; also, daß, wenn sich die Heiden zum Glauben der Juden bekehren, oder zu ihrem Stande, und lassen sich beschneiden, so könnten sie sich rühmen, daß sie auch Abraham angehörten. So weit haben sie es gedeutet.

23. Aber wir sollen es weiter ziehen. Denn
\quellseite{123}
es geht auf den Samen Jesum Christum, wie es auch hernach baß ausgerichtet ist. Und den Segen soll man lassen bleiben auf dem ganzen Menschen, mit Leib und Seele. Denn der ein Christ ist, der ist gesegnet durch und durch, alles, daß er sieht, hört und fühlt; also, daß sich das Wort mächtig weit streckt.

24. So schließet nun dieser Spruch so viel, daß wir alle verflucht und verdammt sind in Adam, welchem das Urtheil aufgelegt ist [1 Mos. 3, 16. 17.]: „Verflucht sei der Acker um deinetwillen“; und dem Weibe: „Ich will dir viel Kummers schaffen, wenn du schwanger bist“ etc. Die Sprüche und Flüche gehen über Mann und Weib und uns alle, um der Sünde willen, damit die Sünde gestraft und dazu gewehrt würde. Aber solches alles hört nun auf, wenn Christus kommt, der für den Fluch eitel Segen bringt; nicht sofern, daß er den Schweiß und Arbeit, und den Schmerz der Geburt wegnimmt, sondern viel höher. Denn der Segen ist ganz geistlich, also, daß er am ersten die Sünde hinwegnimmt und würgt, um welcher willen der Fluch über uns gegangen ist; welche auch der rechte Fluch ist, dadurch wir im Schlunde des ewigen Todes stecken; wie oben [Cap. 2, 17.] gehört ist: „Welche Stunde du von dem Baume essen wirst, wirst du des Todes sterben.“

25. Wie er aber des Mannes Arbeit, und des Weibes Angst hinweg nimmt, sieht man nicht. Denn er nimmt es also hinweg wie den Tod, alsofern, daß der Jammer und Tod bleibt, aber gibt uns solche Arznei, macht uns so trunken, daß uns dies alles nicht irret, und leiden es mit fröhlichem Gewissen. Denn wer da glaubt, der geht fröhlich hinan,\footnote{Erlanger: hinein.} nicht mit Unlust und Bitterkeit wie zuvor; und ob es wohl Fleisch und Blut wehe thut, so bleibt doch das Gewissen fröhlich. Das Wegnehmen ist viel herrlicher, denn wenn er es gar aufhöbe; denn da erzeiget er seine Gewalt und Stärke, wie er so mächtig ist, daß mitten im Tode das Leben kommt und den Tod ersäuft. Desgleichen nimmt er auch nicht die Arbeit noch den Schweiß hinweg, sondern macht, daß es nicht wehe thut, und das Herz voller Freude, daß es meint, es könnte ihm nicht wehe thun. Also auch, daß mit dem Weibe macht er, daß sie alles gerne leidet, wenn es auch noch mehr wäre; [es] thut ihr wohl wehe, aber zuvor noch viel mehr; jetzt wird es doch dem Gewissen süß und leicht.

26. Also führt er seinen Segen so stark, daß er alle Menschen so segnet, daß der Fluch muß unschädlich werden; wird es auch endlich gar hinweg nehmen, es muß aber zuvor der ganze Adam sterben. Wenn wir da hindurch sind, so ist es alles hinweg. Also muß man den Segen verstehen, daß er in das Herz wirkt, tröstet und stärkt, daß ihm kein Unglück zu schwer wird. Wo das nicht ist, so ist der Schmerz zu groß und unträglich, ist eitel Unlust und Herzeleid da, also, daß mans immer gerne wollte überhoben sein. Der Unglaube kann nichts willig leiden, fleucht, weil er fliehen kann, daß er ja nicht arbeite, will vorhin genug haben und wohl versorget sein.

27. Also sieht man allenthalben, wie der Unglaube ein verflucht Ding ist, macht alles bitter und unträglich; wiederum, daß der Glaube ein selig Ding ist, bringt alles Gute und Freude. Das ist nun hier versprochen, daß ein solcher Segen kommen sollte und gehen über alle Geschlechte, weil die Welt steht, also, daß man bis auf den letzten Tag sagt und predigt: Abram ist todt, Christus ist im Himmel; noch ist der Segen da. Wo man das Evangelium predigt, verkündigt es, daß, wer da an Christum, Abrahams Fleisch und Blut, glaubt, dem sollen die Sünden vergeben sein, der Tod hinweg genommen und ein ewiges Leben gegeben werden. Dies ist je höchlich gesegnet, wird jedermann angeboten. Willst du gesegnet sein, so glaube an Christum. Darum ist es nichts Anderes, denn das Evangelium predigen, in welchem der Segen täglich gesprochen wird.

28. Das ist nun gewesen das Evangelium, das sie zu der Zeit gehabt haben, das eben dasselbige gesagt, und nichts Anderes ist, denn unser Evangelium, durch welches Christus in aller Welt verkündigt ist. Wer es annimmt und glaubt, der soll errettet werden von allem Unglück. Dasselbe Evangelium ist Adam auch gegeben, aber nicht so klar, als hier dem Abram. Von dem ist es immer fort getrieben und klarer ausgedrückt, bis so lange Christus gekommen ist.

29. So haben wir gesehen das heilige Kreuz, das der heilige Patriarch getragen hat. Er ist nun in das Land gekommen, weiß aber nicht, daß dies sei, das ihm GOtt geredet hatte. Darum zog er auf und nieder, wie ein Pilgrim, kam bis gen Sichem, fast mitten im Lande Canaan.
\quellseite{1001}

Da er nun also gehorchte seinem GOtt und ausgezogen ist, läßt er ihn nicht endlich ohne Trost. Darum sagt Moses:

\begin{schriftzitat}
V. 7–9. Da erschien der HErr Abram, und sprach: Deinem Samen will ich dies Land geben. Und er bauete daselbst dem HErrn einen Altar, der ihm erschienen war. Darnach brach er auf von dannen an einen Berg, der lag gegen dem Morgen der Stadt Bethel, und richtete seine Hütte auf, daß er Bethel gegen Abend, und Ai gegen dem Morgen hatte. Und bauete daselbst dem HErrn einen Altar, und rief den Namen des HErrn an. Darnach wich Abram ferner, und zog aus gegen dem Mittag.
\end{schriftzitat}

30. Dieser Text läßt sich ansehen, als sei es ein schlechter Text, als er auch ist; [es] sind aber drei Reisen darin angezeigt. Erstlich ist Abram das Land durchzogen hin und wieder, und hat sich einmal gelagert bei Sichem. Von da hat er ausgebrochen und fort gerückt, sich nieder begeben zwischen Bethel. Zum dritten ist er von dannen gewichen und ausgezogen gegen Mittag, bis in das Land Egypten, und hat an zweien Orten Altäre gebauet, da ihm GOtt erschien, und zu Bethel.

31. Der Prophet David hat diesen Text recht angesehen, da er [Ps. 39, 13.] spricht: „Ich bin ein Fremdling bei dir, und ein Gast, wie alle meine Väter.“ Item, die Epistel an die Hebräer [Cap. 11, 9.] sagt auch:\footnote{„sagt auch“ fehlt in der Erlanger.} „Durch den Glauben ist er ein Fremdling gewesen im Lande, das ihm verheißen war, als einem Fremden.“ Und Stephanus [Apost. 7, 4. 5.] spricht auch: „GOtt brachte unsern Vater Abraham herüber in dies Land, da ihr nun inne wohnt. Und gab ihm kein Erbtheil darinnen, auch nicht eines Fußes breit, und verhieß ihm, er wollte es\footnote{In den Ausgaben: sie.} geben ihm zu besitzen, und seinem Samen nach ihm, da er noch kein Kind hatte.“ So hat er müssen sein Lebenlang fremd sein, daß er auch sein Begräbniß mußte kaufen [1 Mos. 23, 16. 17.]. Darum hat er seine Hoffnung nicht gesetzt auf zeitlich Gut.

32. Es ist eine wunderliche Historie und wunderlicher Glaube, daß GOtt den Mann so führt von dem Zeitlichen zu dem Ewigen, nennt und deutet ihm zeitlich Gut, nämlich dies Land, und gibt es ihm nicht; aber der Glaube, der am Worte hangt, wiewohl es vom Zeitlichen redet, gibt ihm die ewige Seligkeit. Derhalben liegt es gar an GOttes Wort. Wenn GOtt redet, auch von einem Strohhalm, ist es dennoch ein ewig Wort, daß, wer daran glaubt, wird rechtfertig und fromm, daß er GOtt, und genug hat in Ewigkeit. Darum muß man nicht allein ansehen, wovon er redet; denn es gilt gleich, ist eines wie das andere; sondern du mußt auf das allermeiste seines Wortes wahrnehmen. Wenn du das gewiß hast, so thue die Augen zu, frage nicht, wovon es rede. Ist das Wort da, so hast du genug, denn es ist die ewige Wahrheit und GOtt selbst [Psalm 119, 160. Joh. 17, 17.]. Also ist der Glaube auch ein ewiger Schatz.

33. Das rede ich darum, daß man sich feindlich gebrochen hat, wie die Verheißungen des Alten und Neuen Testaments zu scheiden seien, und gemeiniglich so scheidet, daß im Alten Testamente zeitlich, im Neuen ewig Gut verheißen wird. Solcher Unterschied ist nicht genugsam, und schließt nicht. Denn dem Abram ist hier nicht ewig Gut zugesagt, da er sagt: „Dies Land will ich deinem Samen geben“, denn da redet er von leiblichen Kindern und Land; noch ist das Wort ewig, und der Glaube, der daran hangt, erlangt ewiges Leben und Seligkeit. Darum soll man das Wort GOttes in gleicher Acht bleiben lassen, er rede wovon er wolle, und immer daran hangen. Es ist alles GOttes Kraft und höchstes Gut. Sonst, wo man nach dem Unterschied sollte richten, so wäre Abram nicht durch den Glauben an dies Wort, von dem leiblichen Gute, ein Christ worden.

34. Dies ist nun der Trost, den uns GOtt gibt im Leiden. Denn die Art hat er, daß er den Glauben versucht, zeucht ihn auf, und macht es lange, daß wir ja die alte Haut ausziehen, und an dem bloßen Worte lernen hangen, und darauf pochen. Doch, wenn es zu lange will währen, und zu hart kommt, läßt er ihn nicht ohne Trost. Die Schrift ist aber vernünftig, sagt nicht heraus, wie man ihm mitgefahren hat. Warum ist er nicht zu Sichem geblieben? Er hat es freilich nicht von Lust wegen gethan. Moses zeucht es auch an im 5. Buche, Cap. 26, 5.: „Die Syrer wollten meinen Vater umbringen; der zog hinab in Egypten, und war daselbst ein Fremdling.“ Die Leute im Lande haben ihn nicht können leiden.
\quellseite{125}

35. Denn das ist gewiß, daß Abram hat GOttes Wort gepredigt, und auch also gelebt, wie er glaubte. Hat er denn sollen also predigen und leben, ist es unmöglich, daß er den Leuten gefallen habe. Denn er hat sie müssen strafen, und aus christlicher Liebe sagen: Euer Wesen ist nicht recht, und wider GOtt. Wenn er das nicht gethan hätte, wäre er nicht ein rechter frommer, heiliger Mann, denn die christliche Liebe kann nicht schweigen noch dulden, daß der Nächste irrt und sündigt, muß strafen und bessern, wo sie kann. Das hat er auch gethan; derhalben, wohin er kommt, bleibt er so lange, als er nicht die Wahrheit sagt. Wenn er das Maul aufthut, so muß er davon.

36. Es ist zu der Zeit noch eine feine Welt gewesen, viel weise, verständige Leute, dazu fromm und heilig. Aber er hat ihrer viele gefunden, sonderlich daß heuchlerische und vermessene Leute waren, welche vornehmlich die rechten Heiligen verfolgten, wie Moses rührt im fünften Buche, wie wir auch hernach hören werden. Darum hat er es nicht aus Leichtfertigkeit gethan, daß er so aufgebrochen ist; es hat ihn niemand leiden wollen. Und weil er so im heiligen Kreuze steht, hat niemand unter den Leuten, davon er getröstet würde, tröstet ihn GOtt selbst, setzt aber solche Worte, daß [es] wider allen Schein ist.\footnote{Erlanger: das Widerschein ist.} Alles sollte er sagen: Du bist ein Fremdling unter den Leuten, die dir feind sind, meinen, sie haben das Land inne, du müssest ihrer Gnade leben; ich will es aber umkehren, daß du Herr seiest, und sie nicht. Des hat sich Abram getröstet. Also verfährt GOtt mit allen seinen Gläubigen. Das ist ein Stück.

37. Das andere, daß Abram einen Altar bauet zu Sichem, darnach aber einen, als er gen Bethel kommt, ist nichts Anderes, denn daß der fromme Patriarch für sich und sein Gesinde einen Ort zugerichtet, dazu, daß sie zusammen kämen, das Evangelium zu hören, zu beten und opfern. Denn von Anfang war das der äußerliche Gottesdienst, daß sie Thiere schlachteten und opferten, wie Noah, Kain und Abel; wie wir haben die Messe oder Sacrament, dabei wir predigen und beten. Sie haben auch gepredigt und gebetet, wie wir, haben aber andere Zeichen gehabt, nämlich das Opfer. Zuweilen hat sich’s also begeben, wenn sie es auf den Altar legen, daß es das Feuer vom Himmel angezündet und verzehrt, als ich von Abels Opfer [Cap. 4, 22.] gesagt habe: Aber ich will glauben, daß [es] hier nicht geschehen sei; oder nicht allezeit. Man findet wohl hernach klarlich, daß [es] geschehen sei; als, von Gideon und Manoah im Buche der Richter [Cap. 6, 21. 13, 20.]. Man soll aber nicht denken, daß sie solchen Gottesdienst aufgerichtet haben, daß sie GOtt ein gut Werk damit thun wollten; sondern, weil man muß einen Ort und Stätte haben, da man zusammen komme, und GOttes Wort lehre und handele, zwingt solches die Noth aufzurichten, also daß eigentlich Abram solches um des Predigens willen gethan hat.

38. Was hat er denn gepredigt. Eben das, das wir predigen im Evangelio, nämlich also: Hier ist GOttes Wort, das mir zusagt, daß er mir wolle einen Samen geben, durch welchen alle Welt solle gesegnet werden, und mir will dies Land geben. Wenn wir dem Worte glauben, so sind wir fromm. [Er] hat es also ausgerichtet, wie man GOtt mit dem Glauben dienen solle, und wie in der Verheißung das ganze Evangelium begriffen ist, daß niemand GOtt ohne den Glauben gefallen kann [Hebr. 11, 6.], niemand aber könnte glauben ohne GOttes Wort, das ist, die Predigt.

39. Darnach hat er auch gebeten, daß die Verheißung erfüllt würde, für sich und sein Gesinde, und für seine Feinde, die ihm Leide thaten. Da ist er Hauswirth und Pfarrherr, hat beide Regimente in der Hand. Darum ist diese Historie nicht zu verwerfen. Er ist heiliger gewesen, denn wohl je ein Papst oder Bischof gewesen ist; hat auch das geistliche Amt besser geführt, denn nie keiner der Allerbesten. Das Predigen ist das höchste und vornehmste Amt. Wo das nicht ist, kann vor GOtt nichts einen Priester machen.

40. Also ist das Opfer nicht so zu verstehen, als habe er es gethan, daß es ein gut Werk wäre, sondern als ein Zeichen, damit die Predigt zu bestätigen, wie wir neben dem Worte, den Glauben zu stärken, auch das Sacrament nehmen. Denn, so du das empfäbest, und daneben das Wort fassest, und glaubst, daß sein Leib für dich gegeben, sein Blut für dich vergossen sei, so bist du selig. So hat er auch gesagt: Glauben wir, daß geschehen wird, was GOtt gesagt hat, so wollen wir zu Wahrzeichen
\quellseite{124}
ein Opfer thun, den Glauben zu stärken, auf daß wir desto sicherer werden. Solches hat die Schrift nicht umsonst gesetzt. Denn je weniger nun Gottesdienst, je besser es ist, wenn man nur bei dem Worte und Zeichen bleibt.

41. So hat nun Abram, als er weiter aufgebrochen ist und gen Bethel kommen, abermal einen Altar aufgerichtet, und so geopfert. Bethel heißt ein Gotteshaus. Denn Beth ist ein Haus, El GOtt; also daß, da der Mann gekommen ist, und eben den Ort getroffen, daß er da ein recht Gotteshaus aufrichtete, und GOttes Namen anriefe, spricht der Text. Das ist alles gesagt vom öffentlichen Amte; also daß, wo man findet, daß die Heiligen einen Altar oder desgleichen gebauet haben, daß man lerne, daß es nicht vom heimlichen Gebete und Gottesdienste gesagt sei. Denn dasselbe hat er gethan, daß [es] niemand gesehen hat; dies aber trifft das öffentliche Werk, das er gethan hat vor den Leuten, da man zusammenkommen ist.

42. Denn „GOttes Namen anrufen“ heißt in der Schrift „öffentlich anrufen“. Das ist nichts anders denn der äußerliche Gottesdienst, dadurch das Evangelium mit dem Munde bekannt und vor jedermann gepredigt wird. So wird er gethan haben, daß sein Volk dabei gewesen ist und zugehört hat; denen hat er gepredigt, und Sprüche in [den] Mund gegeben, wie sie GOtt lobeten und danketen um den zukünftigen Christum, und alle Güter, die ihm verheißen waren [Röm. 10, 9. 10.].

43. Das sind die drei Reisen, daß er muß ein Pilgrim sein, und fliehen vor dem heiligen Kreuz. Er wäre wohl gerne geblieben, da er am ersten den Altar bauete, aber er mußte immer fort, daß er wohl geübt würde im rechtschaffenen Glauben; hat immer gedacht, Gott würde ihm einen Ort geben, da er bliebe, mußte aber schweben zwischen Himmel und Erde.

44. Möchte nun jemand sagen: Ist doch verboten in Mose, daß nicht ein jeglicher ihm einen Gottesdienst aufrichte. Item, daß die Schrift sagt, GOtt will nicht leiden die Werke, die wir selbst erwählen, daß sie sollen vor GOtt wohl gethan sein [Matth. 15, 9. Marc. 7, 7.]. Denn, wie ich oft gesagt habe, was wir thun, muß so gethan sein, daß wir können sagen, das gefalle GOtt, und gewisse Sprüche haben, die da sagen, daß [es] GOtt haben will und gebeut. Wo das nicht [ist], soll ich nichts thun vor GOtt. Als, daß ein Christ soll glauben und lieben, ist GOttes Wort. Was ich nun thue, daß ich immer kann sagen: Das thue ich darum, daß es GOtt geheißen hat; so fahre ich recht vor GOtt. Sonst sind alle Werke nichts und verdammt, wie groß und schön sie sein mögen, wie aller Pfaffen [und] Mönche Wesen und Werke sind.

45. Denn kein Werk soll gut heißen und sein, denn daher, daß es GOttes Wort an sich hat; nicht, daß es mich oder dich gut dünkt, vor der Welt und Vernunft köstlich scheint, groß, lang und schwer\footnote{Erlanger: schwerer.} ist: oder daß man große Andacht und gute Meinung darinnen hat; es soll allein seine Güte und Adel davon haben, daß es in GOttes Gebot hergeht. Wenn nun alle Pfaffen und geistlicher Stand in der Welt Einen Spruch könnten aufbringen, daß GOtt ihr Ding befohlen und geboten hätte, so wollten wir es auch auf den Händen tragen. Es gehört mehr dazu denn gute Meinung, und was die Welt machen kann, daß es gut sei. GOttes Wort muß es machen.

46. Wie wollen wir nun den heiligen Vater retten, daß er zufährt, und bauet Altäre, so er doch kein Wort GOttes davon hat? So lange er ihn das nicht heißt, soll er ja keinen Altar bauen, wie heilig er ist, und wie gut er es meint. Antwort also: Daß er freilich Befehl von GOtt hat gehabt, ob es wohl nicht geschrieben ist, daß der Heilige Geist gesagt habe, daß er bauen sollte, sonst könnte man es nicht vertheidigen. Aber das hilft wohl dazu, daß er GOttes Wort gehabt habe, weil der Text spricht: „Er bauete dem HErrn einen Altar, der ihm erschienen war“; aus dem Geiste hat er geschlossen, weil er ihm dasselbe Land geben wollte, daß es ihm gefalle, daselbst Altäre zu bauen. Sollte er in dem Lande sein, so mußte er je predigen; darum mußte er auch einen Altar bauen. Derhalben thut er es nicht aus eigenem Vorwitz, sondern im Gehorsam. Also müssen wir darnach bleiben und uns stöhnen\footnote{stöhnen = stützen.}, daß wir immer sagen, daß sie nichts gethan haben, das GOtt gefallen hat, ohne GOttes Wort. Denn wenn man zuließe, daß man etwas mehr thun möge, so hätten schon Mönche und Pfaffen Recht gewonnen.

\begin{schriftzitat}
V. 10–16. Es kam aber eine Theurung in das Land. Da zog Abram hinab in Egypten,
\quellseite{1002}
daß er sich daselbst als ein Fremdling enthielte; denn die Theurung war schwer im Lande. Und da er nahe bei Egypten kam, sprach er zu seinem Weibe Sarai: Siehe, ich weiß, daß du ein schön Weib von Angesicht bist; wenn dich nun die Egypter sehen werden, so werden sie sagen, das ist sein Weib, und werden mich erwürgen, und dich behalten. Lieber, so sage doch, du seiest meine Schwester, auf daß mir’s desto baß gehe um deinetwillen, und meine Seele bei dem Leben bleibe um deinetwillen. Als er nun in Egypten kam, sahen die Egypter das Weib, daß sie fast schön war. Und die Fürsten des Pharao sahen sie und preiseten sie vor ihm. Da ward sie in des Pharao Haus gebracht. Und es ging Abram wohl um ihretwillen, und er hatte Schafe, Rinder, Esel, Knechte und Mägde, Eselinnen und Kameele.
\end{schriftzitat}

47. Das ist die letzte und schwerste Reise, die Abram thut; hat eine Zeitlang Trost gehabt, nun geht wieder her eine größere Versuchung, denn zuvor. Ist abermal ein groß Exempel des Glaubens. Das Land wird beschwert mit theurer Zeit, die Einwohner können sich kaum ernähren, wo will der Pilgrim bleiben? Er muß wohl hinaus. Also wird er aber[mal] verjagt, weiß nirgend hin, GOtt schweigt stille dazu, und läßt ihn aus dem Lande stoßen, das ihm verheißen ist, da er meint, er sitze am festesten. Also läßt GOtt die Seinen nicht ruhen, daß ja niemand am zeitlichen Gute hange; spielt also mit uns, gibt und rückt es wieder hinweg; und wiederum. Ueber das muß er sein Weib auch in die Fahr setzen, ja, erstlich sich ergeben seines Lebens, weiß keinen Rath, sich zu retten, denn daß er sein Weib dahin gibt.

48. Thut er aber auch recht, daß er sein Weib in die Fahr setzt und wagt ihre Ehre? Item, daß er so böse Gedanken von Leuten hat, so man doch das Beste soll vom Nächsten denken? Antwort: Er thut so viel, als er vermag, daß er GOtt nicht versuche. Weil er das thun kann, daß er sich mit seinem Weibe kann retten, so braucht er desselbigen. Befiehlt es doch alles GOtt. Denkt also: Wenn ich gleich das Weib nicht bei mir hätte, kann mich GOtt wohl schützen, und sie sowohl bei Ehren behalten, als bei mir; weil ich sie aber habe, und solch Mittel kann gebrauchen, will ich GOtt nicht versuchen.

49. Denn der Glaube ist also gethan, daß er sich frei auf das bloße GOttes Wort erwage, und ihm alle Dinge heimstelle, und doch nicht versuche. Was heißt aber versuchen? Es ist zweierlei: Das erste ist auf der linken Seite, wenn man GOttes Verheißung nicht annimmt, er gebe es uns denn in die Hände. Als, wenn Abram, über die Verheißung, das Land so bald in den Händen haben, und nicht eher glauben wollte, denn er es inne hätte, so wäre das Wort hinweg geschlagen und GOtt versucht, als die Juden in der Wüste thaten, welchen GOtt zugesagt hatte, er wollte sie führen und keinen Mangel leiden lassen; aber als sie nun sahen, daß [es] wollte fehlen, fingen sie an zu murren, wollten nicht glauben, sie hätten es denn in der Hand [2 Mos. 16, 2. 3.]. Das ist, GOtt versucht im Mangel, wenn nicht vorhanden ist, was man haben soll.

50. Das andere ist zur rechten Seite im Ueberfluß, wenn GOtt Vorrath gegeben hat, den man brauchen kann, und man läßt das fahren und sperrt das Maul gen Himmel, will ein sonderlich Zeichen haben. Wenn uns hungert, könnte er uns wohl ohne Speise erhalten; wie er denn [5 Mos. 8, 3.] gesagt hat, daß er uns durch sein Wort wolle ernähren; weil er aber auch Brod gibt, will er nicht, daß wir es lassen liegen. So heißt hier GOtt versucht, wenn man Noth und Mangel hat, und der Sache wohl helfen kann, und nicht gebraucht, das GOtt gegeben hat. Der Art sind auch, so da Keuschheit geloben. GOtt hat den ehelichen Stand vorgeschlagen und eingesetzt, daß wir der Hurerei entgehen; wer nun das nicht will annehmen, und die Natur nicht gibt, Keuschheit zu halten, der läßt stehen GOttes Wort, und das er gegeben hat, als sollte GOtt ein sonderlich, übernatürlich Werk und Wunder mit ihm üben.

51. Darum soll der Glaube also gehen und fahren: Wenn Mangel da ist, daß er frei traue auf GOttes Wort; ist aber Fülle und genug da, soll er es brauchen, und nicht etwas Anderes suchen. So geht er zwischen der linken und rechten Seite, bleibt auf der richtigen Mittelstraße. Das lehrt uns dies Exempel. Mangel hat er wohl gehabt, das befiehlt er GOtt, ob er jetzt gleich leben muß in der Fahr; hat aber das Weib, und sieht wohl, wenn er sagte, daß [sie] sein Weib wäre, möchten sie ihn tödten, so nimmt er, das ihm GOtt zu Schutze gibt, und braucht sein.

52. Es ist aber ein groß Ding, sein Weib
\quellseite{126}
von sich thun, und [es] so bloß auf GOtt wagen. Er muß ja keinen Zweifel gehabt haben, GOtt würde sie wohl behalten, muß auch ein fromm Weib sein, daß so gehorsam ist, daß sie sich unter fremde Leute gibt. Summa: es geht recht brüderlich zu, daß sie auch Leib und Ehre dahin setzt, und thut, was ihrem Manne gefällt. Es sind da zwei Leute gleich zu Hause gekommen: Abram, der Vater, Sarai, die Mutter, haben wohl verdient, daß man sie vorzüge, [zum] Exempel und Vorbilde des Glaubens; haben aber viel und große Anstöße erlitten, daß der Glaube je fein rein und lauter bliebe, beide im Mangel und Genüge, und konnten beide darben und reich sein, hoch fahren und niedrig sein, und sich in allerlei Gutes und Böses schicken.

53. So pflegt nun GOtt mit den Seinen zu fahren, daß er sie nicht lange läßt an Einem Orte bleiben, jagt sie hieher und daher; nicht allein um ihretwillen, daß ihr Glaube bewährt werde, sondern auch andern Leuten zu Nutze. Denn Abram hat freilich nicht können schweigen, hat sich auch nicht gebühret, daß er den Leuten nicht predigte von GOttes Gnaden. Darum hat ihn GOtt getrieben durch den Hunger in Egyptenland, auf daß er da auch Nutzen schaffete, und etliche erleuchtete mit rechtem Erkenntniß GOttes. Welches er auch ohne Zweifel gethan hat. Denn es leidet sich nicht, daß jemand mit Leuten umgehe, und nicht offenbare, was ihnen dient zur Seelen Seligkeit. Weil er nun spricht, daß er in Egypten gewohnt hat, und ihm wohl gegangen ist um Sarai willen, hat er es nicht unterlassen, sie zu lehren. So handelt GOtt wunderlich auf Erden, schickt Apostel und Prediger zu den Leuten, ehe sie sich’s versehen, oder eines daran gedenkt; auch die geschickt werden, wissen selbst nicht, wie sie hinkommen.

\begin{schriftzitat}
V. 17—20.

Aber der HErr plagte den Pharao mit großen Plagen und sein Haus, um Sarai, Abrams Weib, willen. Da rief Pharao Abram zu sich, und sprach zu ihm: Warum hast du mir das gethan? Warum sagest\footnote{In den Ausgaben: sagest.} du mir es nicht an, daß [sie] dein Weib wäre? Warum sprachest du denn, sie wäre deine Schwester? derhalben ich sie mir zum Weibe nahm. Und nun siehe, da hast du dein Weib, nimm sie und zeuch hin. Und er befahl seinen Gewaltigen über ihm, daß sie ihn geleiteten, und sein Weib, und alles, was er hatte.
\end{schriftzitat}

54. Da siehst du, wie GOtt dem Glauben immer antwortet, und läßt ihn nicht sinken. Da geht er hin, setzt sein Weib und ihre Ehre in die Fahr. Das war hoch gewagt, sonderlich in des Königes Hofe. Er sorget nicht, setzt GOtt sein Ziel, wie er sie soll wieder überkommen, stimmt seine Zeit, Weise, noch Person, sondern geht frei dahin, und denkt also: Es sei GOtt befohlen, wo ich bleibe und hinkomme mit dem Weibe, ich will es nicht wissen, wie oder wann sie wieder kommen werde. Hangt allein an dem Segen, in welchem ihm GOtt verheißen hatte, er wollte ihn nicht lassen; und wie er versucht wird, so geht er, und zweifelt nicht, er werde ihr wohl wieder zu ihm\footnote{In den Ausgaben (nach lateinischer Redeweise): sich.} helfen, oder ja sonst zum Besten wenden. Wie nun sein Herz steht, so findet er auch. Denn GOtt sorgt so hoch für ihn, daß er auch den König und alle sein Haus angreift mit großen Plagen.

55. Was es für Plagen gewesen sind, zeigt Moses nicht an; aber so weit bringt er den König, daß er Abram sein Weib selbst wieder gibt, befiehlt dazu seinen Gewaltigen und Amtleuten, ihn mit allem, was er hat, zu geleiten. So genau hat GOtt Achtung auf diesen Mann, der vor der Welt scheint, als kennete er oder sähe ihn nicht, ließ ihm sein Weib nehmen. Ehe er sich umsieht, greift er den König Pharao um seinetwillen an; welches der Prophet David hoch angezucht im 105. Psalm, V. 14. Derhalben ist es ein merklich Exempel, daß GOtt sich seiner so hoch annimmt, und mehr auf ihn gibt, denn auf das ganze Königreich. Was in des Königes Reich ist, gilt hundertmal weniger denn der arme, fremde Abram; daß wohl die Propheten solches so haben angesehen, daß wir lernen, wie sich GOtt hält gegen die, so ihm trauen, und glauben, wie seine Gewalt so groß und schrecklich ist, die nicht einem Christen unterworfen sei.

56. Denn GOtt schlägt hier nicht einen schlechten Mann; er war zuvor nicht gewarnt. Jetzt warnt er ihn so hoch, daß er froh wird, daß er Abram zum Freunde behalte; richtet so viel an, daß sich der König muß vor ihm demüthigen. Wie wäre es jetzt so übel zu leiden, daß ein armer Bettler in eines Königes Land käme, daß ihn Land und Leute gnädiger Herr\footnote{In den alten Ausgaben: „gnad herr“.} heißen müßten, [sich] vor ihm fürchten, und froh werden, daß sie
\quellseite{127}
ihm alle Dienste erzeigeten? Noch kann GOtt solches so bald ausrichten. Darum ist es nicht möglich, daß er jemand verlasse, der auf ihn trauet und an ihn hanget; ja, er hilft mehr und besser, denn wir immer wünschen können.

57. Denn Abram hätte solches nie begehrt noch gewünscht, daß der König nach ihm schickt, gibt ihm sein Weib wieder, und läßt ihn aufs beste geleiten. Darum sollen die Gläubigen unverzagt sein in allen Nöthen. Denn so spricht Salomo in Sprüchen [Cap. 16, 7.]: „Wenn jemands Wege dem HErrn wohlgefallen, so macht er auch seine Feinde mit ihm zufrieden.“

58. Der Spruch fließt aus vielen Historien. Wenn GOtt ein Wohlgefallen an uns hat, so lasse ihn sorgen für die Feinde. Er kann sie in einem Hui herum werfen und zu Freunden machen. Es ist nur der heillose, verdammte Unglaube, denn nicht zu wehren ist; wie er mit Augen sieht, so richtet er und geht, will nicht in die Finsterniß treten.

59. Da wird nun aus folgen ein großer Nutzen und Besserung, daß Könige und Fürsten durch ihn gelehrt und bekehrt werden. Denn solche Strafe ist wohl gelungen, daß der König Abrams Freund ist worden. Und ist ein Zeichen, daß er GOtt erkannt habe; also, daß Abram sein Amt ausgerichtet hat mit großem Nutzen. Wie GOtt allewege mit den Seinen handelt, daß er ihnen zu schaffen gibt auf Erden. Darum wirft er sie so hin und her. Das sei von diesem Capitel genug.

60. Aber was nun zwischen Abram und seiner Sarai zu sagen ist, befehle ich andern. [Es] ist auch angezeigt, wie eine feine Ehe es gewesen ist, und wie sie soll gethan sein. Denn Abram seinem Weibe ein Gebot auflegt, zu sagen, daß sie seine Schwester sei, und sich so hin [zu] stellen in die Fahr für ihren Mann. Da ist kein Widerstreben, sondern fein freundlicher Gehorsam, [sie] hält sich nach ihres Herrn Wort, daß rechte eheliche Liebe da gewesen ist. Darum hat auch sie allda ein mächtig groß Werck des Glaubens gethan, denn sie hat sich ebensowohl sammt ihrer Ehre müssen in GOttes Gewalt ergeben. Darum haben beide, Weiber und Männer, reichliche Exempel und Reizung des Glaubens und christlicher ehelicher Liebe.

\section{Das dreizehnte Capitel}

\begin{schriftzitat}
V. 1—18. Also zog Abram herauf aus Egypten mit seinem Weibe, und mit allem, das er hatte, und Lot auch mit ihm, gegen dem Mittag. Abram aber war wohl beladen mit Vieh, Silber und Gold. Und er zog immer fort von Mittag bis gen Bethel, an die Stätte, da am ersten seine Hütte war, zwischen Bethel und Ai, eben an den Ort, da er vorhin den Altar gemacht hatte, und er rief allda den Namen des HErrn an. Lot aber, der mit Abram zog, der hatte auch Schafe und Rinder, und Hütten, und das Land mochte es nicht ertragen, daß sie bei einander wohneten; denn ihre Habe war groß, und konnten nicht bei einander wohnen. Und war immer Zank zwischen den Hirten über Abrams Vieh, und den Hirten über Lots Vieh. So wohnten auch zu der Zeit Cananiter und Pheresiter im Lande. Da sprach Abram zu Lot: Lieber, laß nicht Zank sein zwischen mir und dir, und zwischen meinen Hirten und deinen Hirten, denn wir sind Gebrüder. Steht dir nicht alles Land offen? Lieber, scheide dich von mir. Willst du zur Linken, so will ich zur Rechten, oder willst du zur Rechten, so will ich zur Linken. Da hub Lot seine Augen auf, und besahe die ganze Gegend am Jordan; denn ehe daß GOtt Sodoma und Gomorra verderbete, war sie wasserreich, bis man gen Zoar kommt, als ein GOttesgarten, gleichwie Egyptenland. Da erwählte ihm Lot die ganze Gegend am Jordan, und zog gegen Morgen. Also scheidete sich ein Bruder von dem andern, daß Abram wohnete im Lande Canaan, und Lot in den Städten derselben Gegend, und setzte seine Hütte gen Sodom. Aber die Leute zu Sodom waren böse, und sündigten sehr wider den HErrn. Da nun Lot sich von Abram geschieden hatte, sprach der HErr zu Abram: Hebe deine Augen auf, und siehe von der Stätte an, da du wohnest, gegen Mitternacht, gegen Mittag
\quellseite{128}
, gegen Morgen, und gegen Abend; denn alle das Land, das du siehest, will ich dir geben, und deinem Samen ewiglich. Und will deinen Samen machen wie den Staub auf Erden. Kann ein Mensch den Staub auf Erden zählen, der wird auch deinen Samen zählen. Darum so mache dich auf, und zeuch durch das Land, in die Länge und Breite; denn dir will ich es geben. Also erhub Abram seine Hütten, kam und wohnte im Hain Mamre, der zu Hebron ist, und bauete daselbst dem HErrn einen Altar.
\end{schriftzitat}

1. Dies ist ein kurzes, leichtes Capitel, und eigentlich geschrieben um der Zusagung willen GOttes. Denn hier siehst du die allererste Verheißung, so er Abram thut auf dies Land. Denn im vorigen Capitel haben wir gehört, daß ihn GOtt hieß ausziehen, deutete aber kein anderes Land, ließ ihn schweben zwischen Himmel und Erde, hin und wieder ziehen, daß er noch nicht wußte, welches eigentlich das Land wäre, bis daß er wieder aus Egypten kam, und sich vom Bruder geschieden hatte. Da sagte er erst: Dies ist das Land, das ich dir verheißen habe zu geben; bestätigt und erklärt also seine Verheißung, die er im vorigen Capitel gethan hat.

2. Dies ist erstlich die Ursache, warum es geschrieben ist; da steht es im Grunde. Denn die Verheißung, daß die Juden das Land sollten einnehmen und besitzen, ist daher genommen, und\footnote{„und“ fehlt in der Erlanger.} wird immer hernach weiter angezogen. (Es ist) aber auch darum geschrieben, daß man abermal sehe, wie GOtt nicht abläßt von denen, so da glauben, so lange sie nicht ablassen; sondern, wie er ein Leiden auf das andere schickt, so folgt auch ein Trost nach dem andern; also, daß es fein gemengt ist, daß jetzt die Sonne scheint, jetzt regnet und finster ist, damit sie zugleich in Gutem und Bösem versucht werden. Er hat Abram eine gute Weile her lassen gehen nach dem Troste in Egypten, daß es abermal ein Jahr oder zwei gewähret hat, bis sich die Hirten miteinander zankten, daß er sich von seinem Bruder scheidet. Da kommt nun ein anderer Trost; gibt ihm aber Raum, daß er wandert, ist hier und da ein Gast und Fremdling, bis gen Hebron; da bleibt er hernach. So zeigt dies Capitel auch die Wallfahrt des heiligen Vaters an.

3. Das gehört nun auch zu seinem Glauben, daß er bisher mit seines Bruders Sohn gewandelt hat, muß sich aber jetzt auch von ihm scheiden; ist dazu ein Exempel der Liebe. Da sich seine und seines Vettern Hirten nicht vertragen konnten, welches er ohne Zweifel oft versucht hat, thut er, was er kann, setzt daran sein Gut und Freundschaft mit Lot, ehe er die Liebe zertrenne, und Zorn und Hader lasse werden. Wir sind also geschickt, daß wir wohl um einen Trunk Bier dürfen hadern und rumoren. Er ist im fremden Lande mit Weib und Gesinde, noch hält er so fest an der Liebe und Frieden, daß er alles daran setzt, was er hat, bleibt so in Eintracht mit dem Bruder, läßt ihm zu, seinen Vortheil zu wählen. Ein anderer hätte sich so leichtlich nicht lassen scheiden. So sieht man, wie sie rechtschaffene GOttes Kinder gewesen sind.

4. Aufs letzte ist auch, hier beschrieben das Land Sodom und Gomorra. Das zu verstehen, muß man die Gelegenheit wissen. Das gelobte Land liegt also, daß es Egypten gegen Abend, und den Jordan gegen dem Morgen hat. Am Ende des Jordans liegt das Meer, das man das todte Meer heißt, darinne die fünf Städte ersäuft sind, wie wir hören werden. Dasselbe Land, ehe es umgekehrt und verkehrt ward (sagt der Text), war es wasserreich als ein GOttes Garten, gleichwie Egyptenland. Da preist er es hoch, und zeigt sonderlich das an, daß Egypten köstlicher sei denn das gelobte Land, das so lustig sei, als hätte es GOtt selbst gepflanzt und zugerichtet; ist allerdings eine Schmalzgrube, rechter Kern und Boden des Landes gewesen, da alles genug war; darnach Bethlehem nicht weit davon auch köstlich gelegen ist.

5. Das zeucht er nun darum an, neben der Historie von Abram, anzuzeigen, daß weltlich Gut und Glaube schwerlich bei einander sind. Denn weil das Land so köstlich, reich und voll war, wurden sie böse Buben, die nach GOtt und seinem Worte nichts fragten. Wären sie frömmer gewesen, hätten sie vielleicht Kummer und Noth müssen leiden. Darum bringt sie der Ueberfluß und volle Genüge in solche Sünde und Untugend, daß sie GOtt vertilgen mußte. Solches hat auch der Prophet Ezechiel [Cap. 16, 49. 50.] angezeigt: „Das war die Sünde deiner Schwester Sodom“, sagt er zu Jerusalem, „Hoffart, Fülle des Brodes, und Ueberfluß, und ihr Müßiggang, und ihrer Töchter“, das
\quellseite{129}
ist, der Dörfer und des Landes umher, „und daß sie den Armen keine Hand reichte, niemand konnte bei ihr auskommen; darum ist sie stolz worden, und hat solchen Greuel angerichtet, darum ich sie verderbet habe.“ 6. Darum, sage ich, wo vollauf ist, da ist nicht viel Gnade; sondern wo es schmal ist, da ist man am besten geschickt. Wie wir sehen: Wenn man voll ist und keinen Mangel hat, ist niemand, der GOtt ansehe; wiederum, wenn man einen Tag Hunger leidet, so schreiet jedermann. Also thut die Welt nichts, denn daß sie GOtt erzürnt mit Gut und Reichthum; gibt er es nicht, so zürnt sie auch. So zeigt die Schrift an, daß GOtt gerne mit Armen zu schaffen hat, und wie gefährlich es ist, im Ueberflusse leben; wenn er viel gibt, können wir uns nicht mäßigen, wir müssen es mißbrauchen. Darum ist es eine große Gnade, daß er läßt den gemeinen Pöbel wenig haben, sonst wird er\footnote{In den Ausgaben: wirds.} so ungezogen, daß [es] niemand kann erleiden, wie Salomo sagt in den Sprüchen.

\section{Das vierzehnte Capitel}

\begin{schriftzitat}
V. 1—16. Und es begab sich zu der Zeit des Königes Amraphel von Sinear, Ariok, des Königes von Ellasar, Kedor-Laomor, des Königes von Elam, und Thideal, des Königes der Heiden, daß sie kriegten mit Bera, dem Könige von Sodom, und mit Birsa, dem Könige von Gomorra, und mit Sineab, dem Könige von Adama, und mit Semeber, dem Könige von Zeboiim, und mit dem Könige von Bela, die heißt Zoar. Diese kamen alle zusammen in das breite Thal, da nun das Salzmeer ist. Denn sie waren zwölf Jahre unter dem Könige Kedor-Laomor gewesen, und im dreizehnten Jahre waren sie von ihm gefallen. Darum kam Kedor-Laomor, und die Könige, die mit ihm waren, im vierzehnten Jahre, und schlugen die Riesen zu Astharoth Karnaim und die Susim zu Ham, und die Emim in der Fläche Kiriathaim, und die Horiter auf dem Gebirge Seir, bis an die Breite Pharan, welche an die Wüste stößt. Darnach wandten sie um, und kamen an den Neththorn, das ist Kades, und schlugen das ganze Land der Amalekiter, dazu die Amoriter, die zu HazezenThamar wohnten. Da zogen aus der König von Sodom, der König von Gomorra, der König von Adama, der König von Zeboiim, und der König von Bela, die Zoar heißt, und rüsteten sich, zu streiten im breiten Thal mit Kedor-Laomor, dem Könige von Elam, und mit Thideal, dem Könige der Heiden, und mit Amraphel, dem Könige von Sinear, und mit Ariok, dem Könige von Ellasar: vier Könige mit fünfen. Und das breite Thal hatte viel Thongruben. Aber der König von Sodom und Gomorra wurden daselbst in die Flucht geschlagen und niedergelegt, und was überblieb, floh auf das Gebirge. Da nahmen sie alle Habe zu Sodoma und Gomorra, und alle Fütterung, und zogen davon. Sie nahmen auch mit sich Lot, Abrams Bruder-Sohn, und seine Habe, denn er wohnte zu Sodom, und zogen davon. Da kam einer, der entronnen war, und sagte es Abram an, dem Ausländer, der da wohnte im Hain Mamre, des Amoriters, welcher ein Bruder war Escol und Aner. Diese waren mit Abram im Bunde. Als nun Abram hörte, daß sein Bruder gefangen war, wappnete er seine eigenen Knechte, dreihundert und achtzehn, in seinem Hause geboren, und jagte ihnen nach bis gen Dan, und theilete sich, fiel des Nachts über sie mit seinen Knechten, und schlug sie, und jagte sie bis gen Hoba, die zur Linken der Stadt Damascus liegt, und brachte alle Habe wieder, dazu auch Lot, seinen Bruder, mit seiner Habe, auch die Weiber und das Volk.
\end{schriftzitat}

1. Dies Capitel ist auch leicht und schlecht. Wer da wissen will, wie die Schlacht zugegangen ist, mag die Figur des jüdischen Landes gedruckt ansehen. Warum es geschrieben ist, wollen wir hernach sehen. Erstlich muß man wissen, wie das Land liegt, die Historie recht zu verstehen. Das Land, da die Moabiter, Seir und Edom gewohnet haben, bis an das rothe Meer, ist alles, das man nennt Arabiam Petreäm. Abram
\quellseite{130}
aber hat gewohnt zu Hebron, die zu der Zeit eine hübsche, große Hauptstadt gewesen ist, als hernach Rom in Welschland, also, daß alles auf sie hat gesehen; nun liegt sie auf der Grenze, daß sie das Salzmeer, das ist, das todte Meer gegen dem Morgen hat; item, Amalek, desgleichen die Karnaim, Riesen, Susim und Emim, Horiter, Amoriter auch gegen Morgen. Deren allen sind sie in das Land gefallen, und haben geschlagen alles, was sie gefunden haben auf denselben Flecken und Gebirge.

2. Da werden nun die Könige sonderlich genannt. Der König Amraphel von Sinear und KedorLaomor von Elam sind die Könige von Persien und Medenland, daher gekommen sind und noch bleiben die Elamiter. Wiewohl es in aller Welt geschieht, daß die Namen verändert werden, darnach ein Volk das andere austreibt, und sich in ein Land setzt, als hier die Horiter haben vor Zeiten gewohnt, da hernach die Edomiter saßen. Darum sind die Namen eines Theils\footnote{In den alten Ausgaben: ein Theil.} ungewiß.

3. So ist es nun zugegangen auf die Weise, daß die vier Könige sind gefallen von Morgen herein, durch der Amoriter und Moabiter Land, die haben sie geschlagen; dazu auch die Susim, das sind gewesen die Junker und Edelleute im Lande; die sind auch geschlagen, darum, daß sie wollten die Könige im Lande wieder frei und ihrer selbst Herren machen. Es hat die Zeit bisher jegliche Stadt einen König gehabt, also, daß es ein mächtig Volk gewesen ist. So haben sich nun die fünf Könige gerüstet in die Schlacht, wider jene viere.

4. Diese Schlacht wird nun erzählt, nicht darum, daß GOtt Lust habe zu sagen, wie sich die Heiden mit einander hadern und schlagen, sondern um Abrams willen und des Wunderwerkes, das da geschehen ist. Zwölf Jahre hatten sie gedient der Persen Könige, spricht der Text, im dreizehnten aber, da sie reich worden, meinten,\footnote{In den Ausgaben: „meinten sie, sie“.} sie säßen fest, fielen sie ab, gaben keinen Zehnten noch Zins mehr; das trieben sie ein Jahr. Im vierzehnten aber machten sich die andern auf, und schlugen sie; das andere nahmen sie alles hinweg, wie das Land voll und reich war.

5. Das ist nun das Wunderzeichen, daß die fünf Könige, die doch stark und mächtig waren, item, die Riesen, ein mächtig Volk, mit den Amoritern, auch die Amalekiter, Susim, und alles, was um die Grenze war, geschlagen wurden von jenen vier Königen, daß ihnen niemand widerstehen konnte; und Abram, der einige Mann, sollte sie schlagen, ohne alles Zuthun und Hülfe irgend eines Königes oder Stadt. Denn der Text spricht, er hat seiner eigenen Knechte genommen dreihundert und achtzehn, und niemand mehr, denn drei Brüder, Mamre, Escol und Aner, welche vielleicht Edelleute waren bei Hebron, da er wohnte.

6. Damit Moses anzeigen will, was Abrams Regiment sei gewesen, daß er dennoch sehr groß Gesinde habe gehabt, weil allein an Knechten so viel gezählt werden, welche dazu der mehrere Theil alle Weib und Kind gehabt haben. Mit denen hat er im Lande müssen umziehen, darinnen er nichts Eigenes hatte, hat sie regiert leiblich und geistlich, mit Predigen, Opfern und Beten, wie oben [Cap. 12, § 37 ff.] gesagt. Daneben haben sich dennoch auch etliche zu Hebron zu ihm gethan, und mit ihm gehalten, als diese drei, die hier genannt werden.

7. Also ist nun der fremde Mann aufgebrochen, und darf mit seinem Gesinde allein so viel und mächtigen Feinden nachjagen, von Hebron bis gen Dan, am Ende des gelobten Landes; also, daß er es durchzogen hat von Mittag gegen Mitternacht. Da er nun bis gen Dan gekommen ist, welche zu der Zeit noch nicht Dan, sondern Lesem hat geheißen [Jos. 19, 47.], da hat er sich auch mit Vernunft beholfen, wiewohl sie\footnote{„sie“, nämlich die Vernunft.} der Geist regiert, hat sie heimlich überfallen, nicht bald nachgeeilt, sondern weit vorher lassen ziehen, daß sie sicher waren, und meinten, sie wären hindurch. Daß man sehe, wie GOtt also spielt, als sei es natürlicher Weise zugegangen und durch menschliche Klugheit, und ist doch GOttes Wunderzeichen. So hat er sie nun plötzlich überfallen, da sie sich nicht versahen. Denn wo fünf Könige bei einander lägen, hielten sie es noch für Schande, daß sie sich dreihundert und achtzehn Mann sollten lassen schrecken, aufbrechen, und hinter sich lassen, was sie gefangen und geraubt hätten. Weil aber das hier geschehen ist, weist die Geschichte selbst aus, daß es nicht natürlicher Weise zugegangen sei, sondern müsse ein Wunderzeichen sein. Als sie nun
\quellseite{131}
so übereilt sind, sich getrennt, geflohen und geschlagen, hat er allen Raub, den sie hinweggeführt hatten, wiedergebracht, und freilich viel mehr reiche Beute. Das ist ein Stück von der Historie, wie es zugegangen ist.

8. Das erste, so man hier mag fragen, ist: Weil Abram ein evangelischer Mann ist, der in lauterm Glauben daher geht, und alle sein Leben steht in GOttes Wort, wie geht es denn zu, daß er hier das Schwert führt, und solchen Mord begeht, denn er wird je viel Leute müssen erschlagen haben? Wie kann nun ein solcher Christenmann solch Werk thun, das einem Henker zugehört, so viel Blut vergießen? Sollte er es nicht billig gelitten haben, wie Christus Matth. 5, 40. lehrt, und also sagen: Haben sie mir den Bruder entführt, so muß ich leiden, daß sie mich auch hinweg führen?

9. Antwort: So haben wir gelehrt, ein Christ ist ein solcher Mensch, der da ganz steht in GOttes Willen, thut alles, was er weiß, daß [es] ihm gefällt; und wiederum, ist also geschickt, daß er die Augen gar zuthut von den Werken, sieht nicht, wie groß oder klein, kurz oder lang, wie hübsch oder närrisch es scheint, wie die Vernunft thut, die unter den Werken immerdar Unterschied macht, spricht: Das ist köstlich, das ist gering; der Glaube aber sieht nur auf GOttes Willen. Ist es GOttes Gebot, so thut er’s; wo nicht, so thut er’s nicht, wenn es gleich gülden wäre; weiter steht er nicht. Wenn nun GOtt von mir forderte, meinen Nächsten zu schlagen, so müßte ich es thun, und nach dem Werk nicht richten, sondern nach GOttes Willen. Wo wollte sonst weltliche Gewalt herkommen, wenn er’s nicht geheißen hätte, das Schwert zu brauchen? Weil er es denn heißt, so muß man es thun. Obwohl ein Fürst so gesinnt sein muß, daß er für sich selbst niemand Leid thue, dennoch muß er die Bösen würgen, ob er gleich ihnen lieber das Leben gönnt.

10. Solches muß man nun wohl fassen, wenn man die Historien des Alten Testaments liest, denn da findet man viel Exempel, wie die Heiligen Blut vergossen und Leute gewürget haben. Wer nun hier nach den Werken der Väter richten will, der wird darüber zum Narren, denn sie sind so wunderlich, daß sich keine menschliche Vernunft daraus richten kann. Also ist das auch ein recht närrisch Werk anzusehen, wie auch droben [Cap. 12, 13.], da er sein Weib hingab. Solcher ist die Schrift überall voll. Als, da wir lesen im ersten Buch der Könige [Cap. 20, 31—42.], als Ahab, der König Israel, Benhadad, den König von Syrien, so geschlagen hatte und eingetrieben, daß er ihn in seiner Hand hatte, wollte er barmherzig sein, und seinem Feinde Gnade erzeigen, ließ ihn los, und hieß ihn Bruder; that es darum, daß er meinte, er thäte ein gut Werk. Das wollte GOtt nicht leiden, ließ sobald einen Propheten aufstehen, der trat zu seinem Nächsten, und sprach durch das Wort des HErrn [V. 35—37.]: „Lieber, schlage mich.“ Er aber weigerte sich, ihn zu schlagen. Da sprach er zu ihm: Darum, daß du dem Wort des HErrn nicht gehorchet hast, siehe, so wird dich ein Löwe schlagen, wenn du von mir gehest. Und da er von ihm ging, fand ihn ein Löwe, und schlug ihn. Und er fand einen andern Mann, und sprach: Lieber, schlage mich. Und der Mann schlug ihn wund. Danach ging der Prophet zum Könige, und sprach [V. 42.]: „Darum, daß du den verbannten Mann hast lebend gelassen, sollst du und dein Volk umkommen.“

11. Siehe, ist das nicht ein närrisch Werk vor der Vernunft von dem Propheten gewesen? Wenn es auch ein anderer thäte, ohne Geist und Wort, wäre [es] unmöglich. Was ist hier anders, daß solch Werk löblich macht, denn das: Weil GOtt will, daß man ihn soll\footnote{Jenaer: solt.} schlagen, so muß man ihn schlagen, und nur die Augen zuthun. An solchen Historien müssen sich viel hohe Leute stoßen, wenn das Licht der Vernunft darein sieht; die ist blind und toll, kann nicht anders handeln mit GOtt, denn daß sie ihm will vorschlagen, was köstlich und gut sei; und wie sie richtet, so soll man thun. Das ist GOtt nicht zu leiden, darum richtet er oft solche Werke an, und will sie haben, die Vernunft zu blenden, und sagt also: Glaube du mir, und siehe auf kein Werk, es sei, wie es wolle; heiße ich es dich, so thue es. Also werden wir hernach [Cap. 22, 2.] abermal hören, daß Abram seinen Sohn sollte opfern. War es nicht wunderlich, daß er den einigen Sohn, den er lieb hatte, ohne alle Schuld sollte erwürgen, so doch GOtt zuvor [Cap. 21, 12.] gesagt hatte, er wollte ihm von Isaak einen Samen geben? Noch mußte er die Augen ganz zuthun, nicht sehen, ob es sein Sohn
\quellseite{132}
wäre oder nicht, und frisch mit hin, daß er ihn schlachtete. Das sind rechte Leute, die nicht ansehen das Gebot, sondern den Willen deß, der es geboten hat. So haben sie es wunderlich getrieben, daß die ganze Welt daran zu Narren wird.

12. So antworte nun zu der Frage: ob ein Christ möge das Schwert führen, so er doch jedermann lieben soll? Ein Christenmens<em>ch</em> darf des Schwertes gar nichts um seinetwillen, auch nicht um der andern Christen willen; aber wenn es GOtt will von ihm geführet haben und ihm in die Hand gibt, soll er es thun, frisch hinrichten und würgen.

13. Also ist es auch hier gegangen, daß der Geist Abram so geführt hat, daß er diese Schlacht that als ein Christenmann, nicht als ein Heide, seinen Muth zu fühlen oder Lust zu büßen, sondern GOttes Ordnung, Wort und Befehl nach. Daß aber das da gewesen sei, zeigt an das hernach folgt, wie ihn der Hohepriester Melchisedek lobt und preist, als er von der Schlacht kam; auch beweist es das Werk und Wunderzeichen an ihm selbst. Wenn es GOtt nicht geheißen hätte, hätte er es nimmermehr gethan, noch sich dürfen unterstehen, auf zu sein wider so mächtige Könige, die alles, was sich wider sie setzte, geschlagen hatten. Weil es aber GOtt befiehlt, sonderlich seines Bruders halben zu retten, so thut er es, und geht hin im Glauben. Wenn nun solche Leute kommen, die können durchdringen und gewinnen. So hat GOtt jene wunderlich wollen schlagen, daß auch das Wunderwerk genug bestätigt ist, daß solches aus seinem Willen geschehen, und ihn gefallen habe.

14. Das sei geredet, zu antworten auf die Frage, daß die Christen für sich weder Schwert noch Spieß dürfen führen, noch für die, die recht glauben; weil es aber noth ist um böser Leute willen, daß man denen wehre und die Frommen schütze, mag ein Christ, so er dazu gefordert wird von Gott, und von denen, die an GOttes Statt sind, hingehen und würgen, wie die andern. Darum, wo es käme, daß ein Fürst im Lande oder Obrigkeit sich wehren und schützen muß, und aufgeboten in die Heerfahrt, so gehe hin im Namen GOttes, brenne, raube, würge, was dir vorkommt unter den Feinden; was du thun kannst, das thue frisch mit gutem Gewissen und Glauben. Da mußt du nicht schonen, sondern thun wie Krieges Art ist; nicht denken, daß\footnote{Erlanger: was.} du werdest Wittwen und Waisen machen, sondern ansehen, daß [es] GOtt so ordnet, das Land oder Volk zu strafen. Wirst du aber darüber geschlagen, so laß es GOtt walten.

15. Also wirst du immer lesen in den Historien von den Kindern Israel, wie ihnen geboten ward, alles, was da lebte unter den Feinden, zu würgen, daß sie nicht einen Hund sollten lassen leben. Es ist kein Zorn über die Bösen; dazu braucht er ebensowohl der Frommen als Bösen. Darum hat Abram recht gethan, daß er des Schwertes nicht geschont hat; nur frisch gefochten und gehauen, keine Barmherzigkeit bewiesen. So wirst du wiederum lesen im ersten Buch Samuelis [Cap. 15, 3. 9.], wie die Juden wollten fromm sein, wiewohl sie GOttes Gebot hatten, sie sollten Amalek gar vertilgen, was im ganzen Lande war, dürre und rein ab, nicht ein Vieh, Schaf, Haus, oder Stall, leben lassen. Denn also hatte er zuvor gesagt im andern Buch Mosis [Cap. 17, 14.]: „Ich will den Amalek unter dem Himmel austilgen“; denn sollten sie nachgekommen sein, und alles frisch hinweg geräumt haben, daß kein Stecken geblieben wäre. Da sie es aber nicht thaten, ward er zornig, daß er den König Saul vom Reich stieß [1 Sam. 15, 23.], also, daß auch sein ganzes Blut mußte ausgerottet werden und des Ungehorsams willen.

16. Also ist nun der fromme Abram geschickt gewesen, hat ein Herz gehabt voller Liebe gegen jedermann, wie gegen GOtt. Daher er hernach [Cap. 18, 23. ff.] so emsig bittet für die Sodomiter, daß er auch wohl für sie gestorben wäre. Wo bleibt denn hier die Liebe? Weil GOtt solches haben will, so geht GOttes Liebe über die Liebe des Nächsten. Darum wären diese Könige seinetwillen wohl geblieben, denn er hat sie gehalten wie seine Nächsten; aber weil ihn GOtt heißt schlagen, muß da keine Liebe noch Gnade mehr gehen.

17. Darum lasset uns lernen, was da heißen rechte Werke, dahin alle Historien sehen, als man gar viel findet, von Simson, David, Elia und Elisäo etc., welche vor der Vernunft so eigentlich scheinen, als sei es aus Zorn, Rache und bösem Willen gethan, daß die Vernunft nicht unterscheiden kann zwischen einem zornigen Tyrannen und dem heiligen Propheten; also müssen sie. Noch ist ein großer Unterschied. Hier geht
\quellseite{133}
Gottes Wort und Gehorsam, anderswo geht eigener Muthwille. Darum sehen wir das Widerspiel in der Schrift, daß, die den Feinden wollten Gnade thun, wollten je geistlich und fromm sein, haben GOtt erzürnt, und sind untergegangen.

18. So lerne nun aus solchen Historien, daß gar nicht nach den Werken zu sehen ist, sondern allein auf GOttes Gebot; daß du keines thuest, du hast denn Zeugniß seines Willens; so gehe denn mit Freuden hinan, ob du dich, und alles, was du hast, müssest daran setzen. Da werden denn rechte Leute aus, sind an kein Werk gebunden, thun nichts, denn das GOtt geheißen hat; darum sind sie sicher, daß GOtt alles gefällt, was sie thun. So ist der Patriarch hingegangen mit gutem Gewissen und rechtem Glauben, und wußte, daß sein Morden GOtt aufs allerbeste gefiele. Mit solchen Leuten ist gut streiten, sind beherzt und freudig dazu. Da kann zuzeiten Ein Mann zehn, zwanzig oder dreißig schlagen.\footnote{Erlanger: geschlagen werden.} Sterben sie, so sterben sie wohl. Die andern sind faul und verzagt dazu. Derhalben hat er gut kriegen gehabt mit GOttes Wort, welches sein rechter Harnisch war; darum mußte kein Schwert nachrücken. So sei beschlossen, daß ein Christenmann, er sei, wie ihm wer er wolle, soll\footnote{In den Ausgaben: soll er.} das Schwert führen, wenn er dazu gefordert wird; wo nicht, so lasse er es liegen. Darum, wenn die Obrigkeit von Pfaffen und geistlichem Stande haben will, daß sie mit ins Feld treten, sind sie es schuldig wie die andern, wiewohl sie sich wider Gottes Ordnung sein aus der Schlinge ziehen, und von aller Obrigkeit Gebot wollen frei sein.

19. Das ist diese Historie, darinne wir sehen, daß Abraham die Liebe des Nächsten behalten hat, also, daß er sie dennoch unter Gott gesetzt hat, daß der Glaube und Liebe gegen GOtt soll die Liebe gegen dem Nächsten regieren, daß wir den Menschen nicht höher lieben denn GOtt. Wo die recht geht, so geht sie auch recht gegen dem Nächsten.

20. Darnach soll die Liebe des Nächsten regieren alle äußerlichen Werke, daß man alles thue, was die Liebe fordert; also, daß alle Gebote durch dieselbe sich meistern lassen, damit wir alles thun, was wir sollen, andere Leute zu bekehren und zum Glauben zu bringen, als, mit Juden und Heiden; also, daß man ihnen allerlei Liebe und Dienst thue, was man kann, dem Glauben ohne Schaden. Denn nach der Liebe muß man alles richten, was geboten ist gegen dem Nächsten; wo es dawider ist, soll es aufhören. Ueber die Liebe kann kein Gesetz gestellt werden; sie soll regieren mit Gewalt, daß immer frei bleibe, es sei geboten oder nicht, alles dem Nächsten zugut; wo nicht, daß nur alles zu Trümmern gehe. So soll man nur frei wider des\footnote{„des“ fehlt in der Erlanger.} Nächsten Gebot thun, wider die Liebe gesetzt. Darum haben wir so mancherlei Exempel vorgeschrieben, darinne wir sehen den Glauben gegen GOtt und Liebe gegen dem Nächsten, so wunderlich durch einander, daß es scheint wider alle Vernunft.

21. Das ist der erste Theil dieses Capitels, geschrieben zu unterweisen die mit Kriegsleuten umgehen, nach GOttes Wort oder Ordnung, daß sie da nicht ansehen der Feinde Schuld oder Unschuld, sondern allein darauf, daß [es] GOtt also haben will. Denn dieser heilige Vater hat freilich zugeschlagen, und viel Leute erwürgt. Und ob er es gleich nicht selbst gethan hat, so hat er doch Ursach gegeben aller Schlacht und Mord, der da begangen ist, daß er muß dafür antworten: und hat doch wohl gethan, wird gelobt und gepreiset vor Gott. Ohne daß man allemal je so streite, daß es geschehe den Unsern zu Schutz, welche\footnote{Erlanger: welchen.} die Noth belangt.

22. Daß aber jetzt die tollen Fürsten in der Welt sich hiermit behelfen wollten, gilt nicht. Denn sie nur Krieg anfangen aus lauter Muthwillen, nicht Land und Leute zu schützen, sondern ihre Lust zu büßen; das hebt sie gegen einander. Da muß der gemeine Mann herhalten, um ihres Muthwillens willen Land und Leute verderben.

23. Wo es aber noth ist, ist es nicht ein Muthwille, sondern die Bedrückten zu schützen; denn ist es der Liebe Werk, so ist es denn nicht mehr Sünde, daß man Witwen und Waisen macht, auch zuzeiten die Unschuldigen mit umbringt, so es nicht anders geschehen kann. Denn da geht es zu, wie man spricht: Ein Nachbar ist dem andern einen Brand schuldig. Wenn wir bei einander wohnen, müssen wir auch gemeines Unglücks erwarten. Und ob wir gleich nicht
\quellseite{134}
Ursach geben, doch, weil wir mit in dem Haufen sind, die Ursach geben, müssen wir mit leiden. Wen GOtt gestraft haben will, den straft er; darum müssen wir nichts denn die Noth ansehen, welche GOttes Willen und Zorn bringt.

\begin{schriftzitat}
V. 17—24. Als er nun wieder kam von der Schlacht des KedorLaomor, und der König mit ihm, ging ihm entgegen der König von Sodom in die Fläche, die Königsthal heißt. Aber Melchisedek, der König von Salem, trug Brod und Wein hervor.\footnote{Randglosse: Nicht daß er es opferte, sondern daß er die Gäste speisete und ehrete. Dadurch Christus bedeutet ist, der die Welt mit dem Evangelio speiset.} Und er war ein Priester GOttes des Allerhöchsten, und segnete ihn, und sprach: Gesegnet seist du Abram, dem Allerhöchsten GOtt, der Himmel und Erde besitzt; und gelobt sei Gott der Allerhöchste, der deine Feinde in deine Hand beschlossen hat. Und demselben gab Abram den Zehnten von aller Beute. Da sprach der König von Sodom zu Abram: Gib mir die Seelen, die Habe behalte dir. Aber Abram sprach zu dem Könige von Sodom: Ich hebe meine Hände auf zu dem HErrn, dem allerhöchsten GOtt, der Himmel und Erde besitzt, daß ich von allem, das dein ist, nicht einen Faden noch einen Schuhriemen nehmen will, daß du nicht sagest, du habest Abram reich gemacht. Ausgenommen, was die Jünglinge verzehrt haben, und die Männer, Aner, Escol und Mamre, die mit mir gezogen sind, die laß ihr Theil nehmen.
\end{schriftzitat}

24. Dies ist darum geschehen, daß Abram desto sicherer wäre, daß er mit GOtt gehandelt hätte, und GOtt sein Werk von Herzen wohlgefiele, daß er die Könige strafte, die GOtt wollte gestraft haben. Wie man mehr in der Schrift findet, als [2 Kön. 10, 11. 23.] von Jehu, der des Königes Ahab Geschlecht und alle Priester des Baal so rein austilgte, daß nicht Einer überblieb, wie der HErr geredet hatte. Er zeigte da seine Barmherzigkeit, und wird darum gelobt. Denn solches ist eben vor GOtt, als wenn ein Vater seine Kinder mit Ruthen hauet, den lobt man. Darum hat er Lust dazu, denen das Schwert gegeben wird, daß sie es nur frisch brauchen über die, so es verdienen. Also schlug Moses [2 Mos. 32, 28.] auf einmal dreitausend Mann todt, der doch der sanftmüthigste Mann war, wie die Schrift sagt, darum, daß er Gottes Zorn hinweg nähme. So geschiehet hier dem frommen Abram, daß ihn der Priester segnet, und stärkt ihn in solchem Werk.

25. Zum andern zeigt der Hohepriester an, daß GOtt habe ein Wunder erzeigt, als er sagt: „Gelobet sei GOtt, der deine Feinde in deine Hand beschlossen hat.“ Abram würde es wohl gelassen haben, daß er mit dreihundert Mann vier mächtige Könige sollte geschlagen haben, die da zuvor fünf Könige, und alles, wo sie hinkamen, niedergelegt hatten. Es ist ein öffentlich Wunder, darum sagt er, daß es GOtt gethan habe.

26. So sagt nun der Text: „Melchisedek, der König von Salem, trug Brod und Wein hervor.“ Sodom und Gomorra ist nicht weit gewesen von Jerusalem. Jerusalem aber hat zu der Zeit Salem geheißen. Zu der Zeit sind nicht so große Könige gewesen, als jetzt; [es] hat fast jegliche Stadt ihren König gehabt. So sind nun bei einander gewesen der König von Sodom und Salem, und Abram. Der König aber von Salem war auch ein Priester GOttes; darum brauchte er seines Amts und segnete Abram. Diesen Text hat niemand so angesehen und ausgestrichen, als die Epistel an die Hebräer [Cap. 7, 1. ff.]. Daraus man sehen kann, was es ist, wenn Gott das Wort in das Maul gibt, daß er da findet und sieht, daß kein Mensch gesehen hatte. Am ersten achtet sie das nicht hoch, daß er Wein und Brod aufträgt; aber unsere Lehrer und Prediger haben es feinbildlich getrieben und aufgeworfen, und also gesagt: er habe Brod und Wein geopfert, darum sei es eine Figur des HErrn Christi, welcher auch ein Opfer habe eingesetzt, seines Leibes und Blutes im Brod und Wein. So hat man uns vorgeplaudert und bei der Nase geführt; haben daraus gemacht, was sie nur gelüstet.

27. Wider solchen Irrthum sollen wir uns setzen, und sagen, zum ersten: Im hebräischen Text stehet kein Wort, das da opfern heiße. Die hebräische Sprache ist so reich, als keine Sprache, sonderlich von geistlichen Sachen und Gottesdienst zu reden; darum sie auch sehr viele Worte hat, die da aufs Opfern deuten. Deren steht aber keines hier, sondern ein gemein Wort, Hozij, das heißt eigentlich producere, exprimere, das ist, hervortragen, als aus einem Keller oder Speisekammer, daß [es] so viel gesagt ist: er hat Abram und den Leuten, so von der Schlacht
\quellseite{135}
fangen, zu essen und trinken gegeben, ließ sie wohlleben und fröhlich sein, GOtt zu danken, daß so viel Könige geschlagen waren. Das muß unsern Träumern deuten, daß Christus Wein und Brod solle geopfert haben, als er seinen Jüngern das Sacrament gab.

28. Er hat sich allein einmal selbst geopfert [Hebr. 7, 27.] also, daß er selbst der Priester und auch das Opfer ist, der Altar aber das Kreuz. Köstlicher Opfer konnte er nicht GOtt thun, denn daß er sich hingab, ließ sich würgen, und in der Liebe Feuer verbrennen. Das ist das rechte Opfer. Das Sacrament aber und die Messe ist nur ein Zeichen solches Opfers, wie die Taufe neben dem Worte GOttes und Glauben, da wir nicht wirken, sondern nur nehmen, und GOtt nur gibt. Darum haben wir ihnen diesen Spruch aus dem Maul gerissen, daß er nicht von Opfern, sondern schlecht von Essen und Trinken sagt, den Leuten vorgezeigt, die von der Schlacht gekommen waren.

29. Der führt es aber recht, daß die Epistel zum Hebräern ist; spricht [Cap. 7, 3.] also, daß dieser Melchisedeck sei Christi Figur, nicht daß er Wein und Brod opfert, sondern hervorträgt. Sein Wein und Brod, das er vorbringt, sind die Worte des heiligen Evangelii, damit er die starken Christen speist. Denn die jungen Christen soll man säuberlich speisen mit Milch, als die Kinder; denn sie sind noch nicht mannhaftig genug, in Anfechtung zu stehen, und sich mit dem Teufel zu schlagen; da gehören gerüstete Leute zu, die das Schwert können in der Hand führen, sich wehren und kämpfen. Das sind sie, denen man Brod und Wein vortragen soll, die man stärkt auf das allerhöchste. Also, daß man immer darauf bleibe, wie oft gesagt, daß man die Figuren auf das Predigtamt ziehe, und auf der Lehre lasse bleiben. [Man] hüte sich nur, daß man es nicht auf Werke, die man vor GOtt will tragen und opfern, deute, wie sie auch hier thun.

30. Darnach wird angezogen, daß der König von Salem hier beschrieben ist, und doch nicht angezeigt, wer sein Vater oder Mutter gewesen sei, [es] wird weder Geschlecht noch Freundschaft genannt, auch nicht sein Tod ausgedrückt. Alles fährt plumps herein, mitten in der Historie, bricht auch kurz wieder ab: läßt es ansehen zu sagen, wo oder wie er herkomme, und wo er bleibe; zeigt nur an, was er mit Abram, und Abram mit ihm gethan habe. Das deutet nun die Epistel zum Hebräern meisterlich, sagt [Cap. 7, 3.] also: Daß er ohne Geschlecht, Vater und Mutter, ohne Anfang und Ende seines Lebens ist, ist er vergleichet dem Sohne GOttes, der ewig ein Priester bleibt. Wer nun Christum nicht versteht, der versteht auch diese Figur nicht. Der Prophet David hat sie fein verstanden, da er [Ps. 110, 4.] spricht: „Der HErr hat geschworen, und wird ihn nicht gereuen, du bist ein Priester ewiglich, nach der Weise Melchisedek“; da drückt er eben das aus, daß die Epistel anzeucht, als er spricht: „ewiglich“; doch nicht also, wie Melchisedek, nach der Weise, als er Wein und Brod opferte, sondern wie von ihm geschrieben ist, so bist du (sagt er [Hebr. 7, 3.]) ein Priester ohne Vater und Mutter, und ohne\footnote{„ohne“ fehlt in der Wittenberger und in der Erlanger.} Anfang und Ende.

31. Darum ist nun angezeigt, daß Christus gar viel ein anderer Priester, denn die levitischen Priester gewesen sind. Da war einer zwanzig, dreißig, vierzig oder fünfzig Jahre ein\footnote{„ein“ fehlt in der Erlanger.} Priester; starben aber allzumal dahin. Dieses Priestertum aber währt von der Zeit an, da die Welt angefangen hat, bis zum Ende. Durch diesen Priester haben Adam und Eva müssen absolviret werden; so auch der letzte Christ, der noch geboren werden soll. Er ist eine ewige Person, die immerdar bleibt, denn GOtt (spricht David [Ps. 110, 4.]) hat es geschworen, will fest darüber halten, und sich’s nicht lassen reuen. Wie fein hat der Prophet davon geredet und tief gesehen! [Es] ist alles in den kurzen Worten verfasset, was Christus ist. Er sagt nicht: Ein Priester wie Aaron, sondern wie Melchisedek. Jenes sind zeitliche Priester, haben auch nicht mehr denn das Gesetz gepredigt, aber niemand geholfen; dieser Priester stirbt nimmermehr, kann ewig helfen. Das heißt nach der Weise Melchisedek, daß es betreffe ein ewig Priestertum.

32. Nun ist oft gesagt, was ein Priester sei. Es ist gar ein mächtig, lieblich Wort, daß kein lieblicherer, freundlicherer Name auf Erden sein mag; und viel lieber zu hören, daß man Christum einen Priester heißt, denn einen HErrn, oder anders; ohne daß zu erbarmen ist, daß der Teufel in die Welt gebracht hat, daß man den
\quellseite{136}
Namen so beschmeißt hat und mißbraucht, daß er kein nütze ist gewesen. Unsere papistische Rotte hat den Namen geführt als Priester vor GOtt, aber gehandelt wie die Buben, und nur Seelen gemordet, mit äußerlichem Schein und Gewalt.

33. Priesterthum ist eine geistliche Gewalt, welche nichts anders ist, denn daß der Priester daher tritt, nimmt alle Gebrechen des Volkes auf sich, nicht anders, als wäre es\footnote{Jenaer: wären sie es.} sein eigen, und bittet GOtt für sie, nimmt von ihm das Wort, damit er jedermann tröste und helfe; ist noch lieblicher und tröstlicher denn Vater- und Muttername; ja, der Name bringt uns alle anderen [Namen]. Denn damit, daß er Priester ist, macht er uns GOtt zum Vater, und sich zum HErrn. Wenn ich ihn für einen Priester halte, so weiß ich, daß er nichts thut, denn daß er droben im Himmel sitzt als unser Gnadenstuhl, und da ohne Unterlaß uns vertritt vor dem Vater [Röm. 8, 34.], bittet für uns [Hebr. 7, 25. 1 Joh. 2, 1.], und redet das Beste. Dies ist der höchste Trost, der einem Menschen widerfahren, und keine süßere Predigt ins Herz gepredigt werden kann.

34. Solches hat er nun beweiset im Evangelio mit allen Worten und Werken. Denn er nichts thut, denn daß er den Leuten dient und hilft, und sich jedermann ergibt, dazu auf sich ladet allen Zorn, so wir verdient haben, läßt sich’s sein Blut und Leben kosten, daß er uns versöhne; auch jetzt, wiewohl er droben bei dem Vater sitzt in seiner Herrlichkeit, noch predigt er hier ohne Unterlaß, fegt, reinigt, trägt, und bessert uns durch seine Gliedmaßen, wie er leiblich auf Erden gethan hat, also, daß er uns auch alle zu Priestern macht [Offenb. 1, 6. 5, 10.]. Was könnte immermehr Tröstlicheres erdacht werden den betrübten Gewissen?

35. Was du nun hiergegen hältst, ist eitel Dreck und Unlust. Da hat sich die päbstische Rotte selbst unterstanden, Pfaffen zu machen, die mit ihren Messen und Werken unsere Sünde aus sollten löschen und GOtt versöhnen, Christo und seinem süßesten Namen zur höchsten Schande und Schmach, daß seine greuliche Lästerung und Verleugnung Christi auf Erden gekommen ist.

36. Darum lasset uns je diesen Namen rein machen und in höchsten Ehren halten; [er] ist so herrlich, daß wir wohl davor sollten niederfallen; man hat es auch wohl gethan, da die Christenheit noch recht stund. Wie Paulo geschah; als er die Letze gab\footnote{Erlanger: setzt Gabe. Jenaer: setzt gabe. Wittenberger: letzt gabe. Letze = Valete wird auch häufig „Letze“ geschrieben.} und seinen Abschied machte, gingen ihnen allen die Augen über und fielen ihm um den Hals [Apost. 20, 37.]. Denn er hatte mit ihnen gehandelt und gelehrt, wie ein rechter Priester handeln soll. Das sollte man auch noch wohl thun. Aber die also fahren, wie der Papst durch den Teufel fährt, die Welt nur mit Gesetzen drücken, und voll Sünden machen, daß man die erst soll Priester heißen, dazu die obersten Priester; da lasset uns für segnen. Stockmeister, Henker und Mörder sollen sie heißen. Pfaffen und Priester sollen ihrem Amte genug thun, daß sie sich ganz hingeben und Knechte werden mit Leib und Leben, beten vor GOtt, und predigen vor den Leuten.

37. Das ist der Spruch Davids [Ps. 110, 4.], aus diesem Texte gezogen mächtig, gewaltig, daß Christus ein ewiger Priester ist, nach der Weise Melchisedek, das ist, ein ewiger, unsterblicher Priester; denn er will das Gewissen trösten, nicht auf Werke, oder irgend zeitlich. Denn es ist je der größte Trost, daß sein Priestertum ewig steht; steht es ewig, so ist es nicht möglich, daß wir verdammt werden. Hier ist unser Gnadenstuhl und Priester, der hört nicht auf; wenn wir schon straucheln und zurück fallen, daß uns der Teufel übertäubt, so steht er und fällt nicht: daß GOtt nicht allein geredet, sondern auch geschworen hat, daß er je unsere Gewissen aufs höchste versichere, daß wir an der Zusagung halten; thut einen Eid dazu, und sagt, es soll ihn nicht gereuen, so doch genug wäre, daß GOtt [es] allein sagte. Nun ist Melchisedek ja ein Mensch gewesen, daß er auch gestorben ist; aber weil es die Schrift nicht meldet, so ist er gleich, sagt der Prophet, dem rechten Melchisedek, der wahrhaftig keinen Anfang noch Ende hat.

38. Nun, vom Priestertum habe ich auch vor mehr gesagt, wie wir in Christo alle Könige und Priester sind [Offenb. 1, 6.]. Welches auch der große, mächtige Schatz ist, daß wir durch Christum dürfen vor GOtt treten [Hebr. 4, 16.], bitten, lehren und predigen, und die Gewalt haben, daß wir GOttes Kinder heißen [Joh.
\quellseite{137}
[1, 12.]. Welches sich so weit streckt, daß, was wir bitten, sind wir sicher, daß er uns erhört [1 Joh. 5, 14.], wie ein Vater sein Kind. Es ist je ein überschwänglich Gut, daß er mich armen Madensack so hoch ehrt, daß er mir zuhört und läßt mein Wort gelten. Dazu haben wir auch die Gewalt, was wir predigen, daß [es] so viel gilt, als [ob] es GOtt selbst sagte. Wenn ein Christ predigt, tauft oder absolviert, ist [es] eben so viel, als [ob] Gott selbst herab käme, redete, und alles selbst thäte; [er] hat nichts behalten, das er uns nicht gegeben habe. Was soll er mehr thun?

39. Ist es nicht über die Maße groß, daß ein Mensch soll Gott werden? Noch hat er es auch in der Schrift [Ps. 82, 6.] gesagt: „Ich habe gesagt, ihr seid Götter, und alle Kinder des Allerhöchsten“, denn weil wir seine Kinder werden, müssen wir auch Mitgenossen werden der göttlichen Natur und Namens, wie Petrus [2. Ep. 1, 4.] sagt. Unser Wort und Werk ist alles GOttes Wort und Werk, ohne Unterschied. So nahe ist er bei uns, und so hoch segnet er uns. Das ist das Priestertum Christi, welches alles steht und liegt im Glauben, das die Epistel an die Hebräer [Cap. 11, 6. ff.] so reichlich führt.

40. Dazu zeucht sie auch an dieses Königes Namen, Melchisedek, rex Salem, das heißt auf deutsch, ein König der Gerechtigkeit, König des Friedens. Die zwei Titel gehören eigentlich Christo zu, daß er ist beide, ein König der Gerechtigkeit, und des Friedens. Wie denn? So haben wir gesagt, Gerechtigkeit ist, die uns GOtt gibt im Herzen; also, daß wir dadurch wahrhaftig und gerecht werden, damit er wahrhaftig und gerecht ist; das ist nichts anders, denn seine Gnade, dadurch wir auch im Herzen gesinnt werden wie er, so weise und stark, als er ist; alle Tugenden, die er an sich hat, werden uns mitgeteilt durch das Evangelium, so wir daran glauben, daß wir mit ihm Erben sind im himmlischen Wesen. So ist Christus ein König der Gerechtigkeit darum, daß er das Haupt ist, von welchem alle Gerechtigkeit fließt in die Christen. Denn durch ihn müssen wir allein dazu kommen.

41. Merke nun, welch ein großer, schöner Titel das ist, und welch ein lieblich Reich. Christus ist der König der Gerechtigkeit, der Wahrheit, Lichtes und Lebens; wiederum, der Teufel König der Ungerechtigkeit, des Todes, Finsterniß und Lügen. Alles, das Christus in seinem Reiche wirkt, tauget dahin, daß er die Leute rechtschaffen mache; soll das geschehen, so muß er die Sünde abthun und hinweg nehmen. Darum ist sein Reich nichts anders, denn daß er uns der Sünde entledige. Darum, willst du Christum recht fassen, so fasse die zwei Namen. Trotz nun dem, der ihm den\footnote{Wittenberger und Erlanger: die.} Namen nehme, daß jemand mehr ein König der Gerechtigkeit heiße. Man findet wohl Könige der Hoheit, weltlicher Ehre und Gutes etc., aber keinen König der Gerechtigkeit findet man, denn allein diesen. Derhalben, welche die Gerechtigkeit durch Werke und den freien Willen suchen, sind schon außer Christo, in des Teufels Reiche. Nun ist Christi Reich nichts denn eine stete Uebung auf Erden. Daß wir von Sünden frei werden, das thut niemand denn dieser Priester, der richtet es allein aus, daß wir des Unglückes los werden; nicht mit Treiben und Schlägen durch das\footnote{„das“ fehlt in der Erlanger.} Gesetz, sondern durch das tröstliche Evangelium und die\footnote{„die“ fehlt in der Erlanger.} süße Gnade und Güte.

42. Weiter ist auch der vornehmste Preis und Tugend dieses Königreiches, daß es ein fried- sam Reich heißt. Wo sein Reich ist, da soll auch rechter Friede sein [Röm. 5, 1.]. Nicht äußerlich, sondern im Gewissen, daß daselbig sicher, fröhlich und unerschrocken ist. Das fühlt das Herz, wenn es mit GOtt wohl steht und eins ist. Wer den Frieden noch nicht hat, gehört in das Reich nicht. Also ist kürzlich abgemalt, wie Melchisedek Christi Figur ist mit dem ewigen Königreiche und Priestertum.

43. Zuletzt sagt der Text, wie Abram dem Melchisedek den Zehnten gegeben hat von aller Beute. Das zeucht nun die Epistel [Hebr. 7, 5. ff.] mit vielen Worten an; sie will also schließen: Weil Abram dem Melchisedek den Zehnten gibt, dazu Abram von ihm den Segen, als von einem Priester, nimmt, ist es ein Zeichen, daß Melchisedek größer sei denn Abram. Wiewohl die Juden des Sinnes sind, daß sie halten, Abram sei mehr gewesen, ist doch so viel aus dem Texte bewiesen, daß sich Abram für geringer bekennt. Daraus wird angezeigt, wie das levitische Priestertum geringer sei, denn das Priestertum Christi. Die Priester, so von Levi sind, kommen von Abram, sind auch durch ihn verzehnet dem Priester Melchisedek, dieweil sie noch in des
\quellseite{138}
Abram Lenden waren. Da ist kein ewiges Priesterthum. Soll denn Melchisedeks Priesterthum ewig sein, nicht nach den Zeiten, so muß es je köstlicher sein denn Abrams, und was von ihm herkommt.

44. So ist nun hier das Letzte nach der Historie, daß bereits zu der Zeit die Priester den Zehnten genommen haben. Denn GOtt hat von Anbeginn der Welt das Amt versorgt, daß sie ihre Nahrung von anderen Leuten nehmen; auch hat es Christus Matth. 10, 10. bestätigt, und Paulus, als, Gal. 6, 6., da er sagt: „Der da unterrichtet wird mit dem Worte, der theile mit allerlei Gutes dem, der ihn unterrichtet.“ Es ist eine feine, göttliche Ordnung, wäre wohl werth, daß man dabei geblieben wäre in der Christenheit, daß man noch so einen Pfarrherrn versorgte mit dem Zehnten. Nun aber die Freiheit eingerissen ist, ist den Predigern befohlen auf ihr Gewissen, daß sie nehmen, so viel ihnen noth ist; wiewohl St. Paulus solches nicht hat gebraucht noch gefordert, daß er doch Recht hatte. Weil nun der Zwang aus ist im neuen Testamente, so ist es nicht noth, den Zehnten zu geben. Wäre doch solche Ordnung, sage ich, noch fein, wo wir’s selbst unter einander eins würden; man müßte aber auch die Priester, die da predigen, allein solches lassen nehmen, weil GOtt geordnet hat, daß sie sich nähren vom Evangelio [1 Cor. 9, 14.], aber die da geistliches Standes sein wollen, und nicht das Evangelium predigen, sollten es nicht haben. Wenn man solche findet, könnte man sie wohl damit nähren; man wird ihrer doch wenig genug finden.

45. Der andern ist die Welt voll, die alle genug haben, und kein Evangelium predigen; haben mehr denn die Hälfte der Welt Güter, wenn man es gleich austheilt; wenden nicht mehr vor, denn daß sie geistlich sind. Wozu sind sie geistlich? Daß sie schlemmen und demmen, predigen kein Wort GOttes. Die sind es, die den edlen, reinen Namen so schänden, den man sollte für Heiligthum halten.

46. Das sage ich darum, daß man wisse, wie es GOtt selbst geordnet hat da zu der Zeit, und hernach lange gehalten ist, daß man den Priestern den Zehnten mußte geben. Wie es aber jetzt steht, wird es noch dahin kommen, daß man das Evangelium muß fallen lassen, allein darum, daß man denen, so es predigen, ihre Nahrung nicht geben will. Daß es andere umsonst thun sollen, ist nicht recht; denn es weist wohl aus, daß solche Leute keine Lust zum Evangelio haben, denn jedermann umsonst predigen soll. Aber die des Teufels Prediger sind, denen soll man genug geben, wie man auch gethan hat; dazu gibt die ganze Welt mit Haufen. Wir stellen uns jetzt schon auch also: dem Evangelio will und kann niemand geben, haben alle Sorge, wir werden arm und verhungern.

47. Das letzte Stück in diesem Capitel ist, daß der König von Sodom zu Abram spricht, er soll die Habe behalten, und ihm allein die Seelen geben; Abram aber schwört, er wolle von dem allen nicht einen Faden behalten, daß er nicht sage, er habe ihn reich gemacht. Das beschreibt Moses auch, als liege große Macht daran; scheint doch als gar eine schlechte Historie. Es ist ohne Zweifel großes Gut da gewesen; ist zu denken, da vier Könige geschlagen sind, die so viel Landes hatten gewonnen, und groß, mächtig Gut zusammen gebracht, und Abram ihnen alles wieder abnimmt. Es muß eine große Beute gewesen sein, daß [er], wenn ers behalten hätte, freilich so reich wäre worden als der König einer. Nun erkannte der König von Sodom solch Wunder, daß Gott durch ihn gethan hatte, wollte ihm dankbar sein, daß er ihm alle Habe schenkte, ohne die lebendigen Leiber. Deß weigerte er sich, will auch nicht einen Schubriemen davon nehmen.

48. Ich habe vor [Cap. 2, § 14] gesagt, daß man müsse in der Schrift gewohnen des Wortes „Seele“, daß es nicht allein den Geist, der da von dem Leibe scheidet, wie wir es in unserer Sprache führen, sondern den lebendigen Leib, lebendige Leute, heißt. Das ist die Weise auf hebräisch zu reden, wie auch Christus im Johanne [Cap. 10, 12.] redet: „Ein guter Hirte lässet seine Seele für seine Schafe“; item [Cap. 12, 25.]: „Wer seine Seele hasset“ etc.

49. Was lehrt nun der Heilige Geist in dem Exempel? Eine große, edle, tapfere Frucht des Glaubens, wie fein frei der Mann hangt an GOtt, so gar ohne allen Geiz; daß solch Exempel der Heilige Geist für gut angesehen, daß es geschrieben würde, uns zu stärken wider das Laster. Denn siehe, wie große Ursache er hat, das Gut zu sich zu nehmen. Erstlich ist er fremd im Lande, so wäre es je ein großer Vortheil, daß er auch so viel hätte, als der König, daß er sich möchte setzen, und im Lande bleiben mit


\quellseite{139}
allen Ehren; [er] hat oft viel müssen leiden von seinen Nachbarn, daß er nimmer konnte sicher und ruhig sitzen. Und noch viel redlicher Ursache, denn er hat es selbst erworben mit seiner Arbeit, Leib und Leben daran gesetzt, daß es auch vor GOtt ihm gebührte; dazu, daß es ihm der König gerne gibt, und selbst anbeut, ja, froh wird, daß er nur errettet ist; also, daß er auf allen Seiten guten Fug und Recht dazu hat.

50. Warum scheuet er sich’s denn zu nehmen? Nicht anders, denn daß er bösen Schein hat wollen meiden; wie er selbst sagt: „Daß du nicht sprechest, du hast Abram reich gemacht.“ Die Ehre und Ruhm konnte er dem Könige nicht lassen. Denn wenn er es hätte angenommen, würde er gesagt haben: er hätte Abram so viel geschenkt, und hätte müssen sein Schuldner sein; wie allemal geschieht, wenn einer solches annimmt, so ist es lieblich; aber es sähet bald. Wie man auch spricht: Es wird nicht theuer gekauft, denn das man geschenkt nimmt; gibt man es nicht zwiefältig wieder, so spricht man, er sei undankbar, und muß verbunden sein, daß man Recht zu einem hat. So ist es auch im weltlichen Regimente: daß [es] etwas Besseres\footnote{In den Ausgaben: „etwas besser“ in der Jenaer mit der Conjectur „etwa“ am Rande.} ist, einer kaufe und bezahle es frei, denn daß er es ihm schenken lasse. So will er hier auch thun, dem Könige unverbunden sein, ob er wohl Recht dazu hat.

51. Aber die größte Ursache ist diese: daß er GOtt seine Ehre rein behält, und beweist seinen Glauben, wie er gar nicht hange an zeitlichem Gut, und um böses Ansehens willen lieber läßt sein Recht fahren. Nun lässet uns die Welt ansehen, wie viele würde man finden, die es ihm nachthun? Wir sind so gar im Gut ersoffen, daß man eher alles in die Schanze schlägt, ehe jemand etwas läßt umsonst fahren, zuvorans, so man Recht dazu hat. Da siehe, was [es] für ein Geist gewesen ist, ob er in das Schwert gezückt habe, daß er Land und Leute und Gut kriegte; wie jetzt die tollen Fürsten thun. Darum zwingt die Historia, daß Abram gar einen hohen, aufrichtigen Geist in dem gehabt habe; gar nicht seinen Nutzen gesucht, sondern allein seines Nächsten, daß er sich hier solch groß Gut und Nutzen verzeiht. Darum spricht er nun: „Ich habe einen HErrn, zu dem hebe ich meine Hände auf“, der da gegeben hat, daß die Könige geschlagen sind, wird auch wohl geben, was mir noth ist. Ich will GOtt allein den Ruhm geben, und nicht dir, daß ich etwas von dir habe.

52. Das ist dies Capitel; das ist uns geschrieben zur Unterweisung, daß wir nicht thun wie die Narren, so nichts vom Alten Testamente halten. [Es] ist eitel großer Vorrath und Schatz, der uns gegeben ist, damit wir unsern Glauben stärken, tausendmal besser, denn alle Legenden, so je geschrieben sind.

\section{Das fünfzehnte Capitel}

\begin{schriftzitat}
V. 1—5. Nach diesen Geschichten begab es sich, daß zu Abram kam das Wort des HErrn im Gesichte, und sprach: Fürchte dich nicht, Abram, ich bin dein Schild, und dein sehr großer Lohn. Abram sprach aber: HErr, HErr, was willst du mir geben? Ich gehe dahin ohne Kinder, und mein Hausknecht hat einen Sohn, dieser Elieser von Damasco. Und Abram sprach weiter: Mir hast du keinen Samen gegeben, und siehe, der Sohn meines Gesindes soll mein Erbe sein. Und siehe, das Wort des HErrn redete mit ihm: Er soll nicht der Erbe sein, sondern der von deinem Leibe kommen wird, der soll dein Erbe sein. Und er hieß ihn hinaus gehen und sprach: Siehe gen Himmel, und zähle die Sterne, kannst du sie zählen? Und sprach zu ihm: Also soll dein Same werden.
\end{schriftzitat}

1. Dies ist auch ein reich Capitel, um des einigen\footnote{Erlanger: einen.} Spruches willen wohl zu merken, den Paulus anzeucht, nicht an Einem Orte. Siehe, Abram hat noch kein Kind, und ihm wird zugesagt, nicht allein, daß er soll Kinder kriegen, sondern auch, daß sein Same so groß werden soll, daß GOtt um ihretwillen das Land will
\quellseite{140}
strafen; nicht allein eines, als Egypten, da er sie heraus führte, sondern auch andere, nämlich das, da er ein Gast in ist; also, daß er sie vertilgen und seinen Samen hinein setzen will. Hier muß man nun die Historie fassen, was GOtt damit meint. Wie er die Kinder Israel aus Egypten führt, und in das Land setzt, werden wir hernach hören im andern Buche. Die Verheißungen stehen hier, sind aber noch finster, so lange bis es dahin kommt, daß\footnote{Erlanger: da.} es alles erfüllt und geschehen ist.

2. Zum ersten, daß er spricht: Nach diesen Geschichten kam zu Abram das Wort des HErrn, und sprach: „Fürchte dich nicht“ etc., ist aber[mals] nicht des Gesetzes Wort, sondern eine Verheißung und evangelisch Wort. Denn wir haben gehört den Unterschied unter dem Evangelio und Gesetze. Gesetz ist, wenn GOtt heißt thun oder lassen, und etwas von uns fordert. Evangelium ist, wenn er etwas Gutes verkündigt, das er selbst thun und uns geben will. So gehen die zwei gar wider einander: eines treibt, das andere gibt; das fordert, jenes schenkt. So soll man alle Sprüche, die GOttes Gnade und Huld verheißen, nicht anders heißen, denn Evangelium, welches so viel heißt als eine gute Botschaft.

3. Zum andern, siehe abermal, wie GOtt so genau Aufsehen hat auf die Seinen. Abram ist fremd im Lande, hält so an GOtt, daß er auch des Gutes nicht achtet, dazu er groß Recht hatte, bleibt williglich in der Fahr unter den Feinden; darum läßt ihn GOtt nicht ungetröstet. Wenn wir auch also glaubten, würde er uns auch solchen Trost hören lassen; wir können aber nicht dazu kommen, denn wir versuchen es nicht.

4. Das Evangelium hat nun Abram gefaßt, dadurch sein Glaube stärker worden ist denn zuvor; wie GOtt pflegt immer, je mehr und mehr man predigt, stärker zu machen. Darum ist er der Verheißung froh worden, daß GOtt sagt, er soll sich nicht fürchten; hat ihn sicher gemacht, daß ihm nichts schaden würde. Item, daß er weiter redet: „Ich bin dein Schild“ und Schußherr, der dich beschirmen will allenthalben; dazu nicht allein dein Schild, sondern ich will auch „dein Lohn“ selbst sein, und dir es reichlich vergelten. Das\footnote{Erlanger: was.} ist dasselbige, daß er ihm so einen großen Samen will geben, und dazu das Land.

5. Auf diese Worte gehört nun ein Glaube; denn hätte er es nicht geglaubt, wäre es verloren. So hat er nun müssen also thun, daß er sich hat erwogen,\footnote{In den alten Ausgaben: „erwegen“; sich erwegen = sich darauf verlassen.} daß er keines Menschen Schutz bedürfte, und einen gewissen Schutzherrn habe vom Himmel. Es sind treffliche, große Worte, auch ein mächtiger Glaube, der sie faßt; [er] sieht der keines, sondern hört nur die Worte, und verläßt sich darauf. GOtt hatte ihn gerufen aus seinem Vaterlande, und zugesagt, er wollte\footnote{Erlanger: wolle.} ihn segnen, und zum großen Volke machen, wie oben [Cap. 12, 2.] gehört. Das hat er nicht verstanden, wie es sollte zugehen, noch nicht gewußt, ob das Volk sein natürlich Fleisch und Blut sein sollte, oder sonst ein groß Volk zufügen. Denn GOtt kann wohl einen Mann aufwerfen, und ihm viel Volkes zubringen, als er hier schon so groß worden ist, daß er dreihundert Knechte im Hause hat, derer keiner sein Blut noch Fleisch ist.

6. Da forchte er nun nach, und begehrte zu wissen, ob es von seinem Blute und Fleische, oder allein seines Gesindes sein sollte; spricht also: „HErr, was willst du mir doch geben?“ Als wollte er sagen: Ich begehre zwar nicht viel Gutes, habe dem Könige von Sodom alle Leute wieder gegeben, hätte wohl Begierde, wenn es möchte sein, einen Erben zu haben; du gibst mir aber keinen, sondern mein Hausknecht mit seinen Kindern, hoffe ich, ist mein Gut und Erbe. Als sollte er sagen: Du gehst je wunderlich mit mir um, gibst so viel Leuten Kinder, denen du gar nichts verheißest; mir aber, dem du so große Gnade erzeigest und solche Verheißung thust, gibst du keines; wie soll ichs nun gewarten? Da gibt ihm GOtt einen Trost und sagt: „Deines Knechtes Sohn soll nicht dein Erbe sein“; du sollst nicht durch dein Gesinde zum Volke werden, sondern ich will dir aus deinem Leibe einen Sohn geben. Und noch weiter sagt er: „Siehe gen Himmel, und zähle die Sterne, ob du sie zählen kannst; so groß soll dein Same auch sein.“ Das sind die Worte der Zusagung, darauf folgt Mosis Wort und spricht:

\begin{schriftzitat}
V. 6. 7. Abram glaubte an den HErrn, und das ward ihm zur Gerechtigkeit gerechnet. Und
\quellseite{141}
er sprach zu ihm: Ich bin der HErr, der dich von Ur aus Chaldäa geführet hat, daß ich dir dies Land zu besitzen gäbe.
\end{schriftzitat}

7. Dies ist der große Hauptspruch in diesem Buche, den Paulus [Röm. 4, 3.] führt und recht ausstreicht, daß [es] Wunder ist; sagt, es sei nicht um seinetwillen geschrieben, denn es lange nach seinem Tode geschrieben ist, sondern um unserwillen, daß wir auch so glauben. Darum ist der Text nicht in Wind zu schlagen. So siehst du, daß hier gar kein Werk erstlich angezeigt wird, nicht mehr vorgetragen, denn daß GOtt Abram zusagt, er sollte von seinem Leibe ein Kind kriegen. Da thut er nichts, betet noch wirkt nichts, sondern glaubt allein; und der Glaube wird so hoch gerechnet, daß er darum vor GOtt fromm und heilig ist.

8. Derhalben schließen wir gewaltiglich mit Paulo [Röm. 3, 28.], daß der Glaube ohne alle Werke rechtfertigt. Warum das? Darum, sagt St. Paulus, daß es Gott die höchste Ehre gab, so man ihm geben kann, daß er ihn für wahrhaftig hielt und gehorsam war. Denn er hätte wohl Ursache gehabt zu zweifeln, weil er nun war bei neunzig Jahr alt, und auch sein Weib nicht viel jünger, daß sie beide natürlich nicht mehr tüchtig waren, so lange dahin in dem Alter,\footnote{„so lange dahin in dem Alter“ = so weit vorgerückt im Alter. „so“ von uns gesetzt statt „zug“ in den Ausgaben.} daß [es] unmöglich war, daß sie sollten Frucht zeugen; wie er auch hier sich bekümmert hat, daß er müsse ohne Kinder hingehen.

9. Da streitet er nun wider Vernunft, Natur und wider Hoffnung, daß er ein Vater solches reichen Samens, der so groß würde werden, sein sollte. Was thut er aber? Er hat die Worte gefaßt, die GOtt hier spricht, „siehe nicht an (spricht Paulus Röm. 4, 19.) den erstorbenen\footnote{In den Ausgaben: „verstorbenen“.} Leib“, und daß ihr Mutterleib ganz veraltet war; da war keine Hoffnung, eine Frucht zu haben; doch thut er die Augen zu, sieht nicht, wie unmöglich es ist, fährt im Finsterniß, hält sich an das Wort; da setzt er seine ganze Zuversicht auf, und war gewiß, es könnte nicht lügen. Damit gibt er GOtt die Ehre, daß er sich für einen Narren, GOtt für weise, gütig, wahrhaftig, gewaltig, allmächtig hält. Solches thun ist der rechte Gottesdienst. Denn das ist die rechte Ehre, wenn du ihn von Herzen ehrst, das ist, wenn du so von ihm hältst, als von GOtt zu halten ist, daß er nicht lügen kann, und alles thun, was er sagt, es scheine, wie es wolle. Das kann der Unglaube nicht. Also hat er das erste und höchste Gebot erfüllt.



10. Siehe, so bauet St. Paulus auf diesen Text seine Predigt, führt ihn dürr, rein und stark, und schließt, daß der Glaube nichts anders ist, denn eine Tödtung des alten Adams, der so nach aller Vernunft und Sinne müßte sagen: Ist es doch unmöglich, ungehört, und wider die Natur, darum wird es nicht wahr sein, wird ein Traum oder Betrug sein. Also urteilt der alte Mensch von GOttes Wort. Aber solcher Dünkel stirbt ganz in dem Abram, und er wird stark und lebendig im Wort, daß er darauf trotzt, es werde und müsse also geschehen, wenn auch alle Welt anders sagt, und alle Sinne trügen. So gewinnt er gar einen andern, neuen Wahn und Verstand über die Natur. Das heißt nun, allein durch den Glauben rechtfertig werden.

11. So befehle ich nun, dies Exempel wohl zu bewegen aus St. Pauli Epistel [Röm. 4, 1. 2. 3.], der es mit höhern Worten führt. „Was wollen wir sagen (spricht er) von unserm Vater Abram, daß er funden habe nach dem Fleisch? Das sagen wir: Ist Abram durch die Werke rechtfertig worden, so hat er wohl Ruhm, aber nicht vor GOtt“; als wollte er sagen: GOtt fragt nicht darnach, daß ein Mensch durch Werke sich will rechtfertig machen. „Wie spricht aber die Schrift? Abram hat GOtt geglaubet, und das ist ihm zur Gerechtigkeit gerechnet.“ Da steht je die Schrift dürr und klar, daß er durch den Glauben, ohne Werke, ist rechtfertig worden.

12. Das beweist er darnach weiter durch das, so folgen wird im 17. Capitel, V. 11., von der Beschneidung, welche lange hernach gekommen ist. Und sollte es irgend ein Werk gethan haben, so sollte es auch das gethan haben. [Es] ist aber lange zuvor beschrieben, daß er rechtfertig sei, ehe er dieses Werkes gedenkt. So geht der heilige Vater her in GOttes Worte, der da sagt, er wolle sein Lohn sein, und ihm ein Kind und großes Volk von seinem eigenen Leibe geben, der doch nun todt war, Frucht zu zeugen, nach der Natur. Spricht nun GOtt abermal zu ihm: „Ich bin der HErr, der dich aus Chaldäa geführet hat“ etc. Da erklärt er sich, was er meint, und beschließt, daß dies das Land sei, das er ihm
\quellseite{142}
verheißen hat, wie auch droben [Cap. 13, 15.]. Das sind alles eitel starke evangelische Worte und Predigten, damit angezeigt ist, daß Abram betrübt gewesen ist, und noth gewesen, ihn mit so viel Worten zu trösten. Denn solche Worte GOttes und hoher Trost kommen nicht vergeblich oder ohne Noth.

\begin{schriftzitat}
V. 8—21. Abram aber sprach: HErr, HErr, wobei soll ich es merken, daß ich es besitzen werde? Und er sprach zu ihm: Nimm eine dreijährige Kuh, und eine dreijährige Ziege, und einen dreijährigen Widder, und eine Turteltaube, und eine junge Taube. Und er nahm solches alles, und zertheilte es mitten von einander, und legte einen Theil gegen das andere über; aber die Vögel zertheilte er nicht. Und das Gevögel fiel\footnote{Randglosse: Gevögel fiel. Das Gevögel und der rauchende Ofen, und der feurige Brand bedeuten die Egypter, die Abrams Kinder verfolgen sollten. Abram scheucht sie davon, das ist, GOtt erlöst sie, um der Verheißung willen, Abram versprochen. Daß er aber nach der Sonne Untergang erschrickt, bedeutet, daß GOtt seinen Samen eine Zeit verlassen wollte, daß sie verfolgt würden; wie der HErr selbst hier deutet. Also geht es auch allen Gläubigen, daß sie verlassen, und doch erlöst werden.} auf das Aas, Abram aber scheuchte sie davon. Da nun die Sonne untergegangen war, fiel ein tiefer Schlaf auf Abram, und siehe, Schrecken und große Finsterniß überfielen ihn. Da sprach er zu ihm:\footnote{Jenaer: Abram.} Das sollst du wissen, daß dein Same wird fremde sein in einem Lande, das nicht sein ist, und da wird man sie zu Knechten machen, und übel handeln vierhundert Jahre. Aber ich will Richter sein über das Volk, deß Knechte sie werden sein. Danach sollen sie ausziehen mit großem Gut. Und du sollst fahren zu deinen Vätern mit Frieden, und in gutem Alter begraben werden. Sie aber sollen nach vier Mannszeiten wieder hierher kommen; denn die Misethat der Amoriter ist noch nicht gar hie. Als nun die Sonne untergegangen und finster worden war, siehe, da rauchte ein Ofen, und ein feuriger Brand fuhr zwischen den Stücken hin. An dem Tage machte der HErr einen Bund mit Abram, und sprach: Deinem Samen will ich dies Land geben, von dem Wasser Egypti an, bis an das große Wasser Phrath, die Keniter, die Kinisiter, die Kadmoniter, die Hethiter, die Pherisiter, die Riesen, die Amoriter, die Cananiter, die Gergesiter, die Jebusiter.
\end{schriftzitat}



13. Abram ist nun gewiß, daß der Same, der ihm verheißen war, von seinem Leibe kommen sollte. Der Vater ist gedeutet, aber die Mutter ist noch im Zweifel, als wir hören werden. Denn GOtt übt ihn noch immer mehr und länger. Was ist’s nun, daß er spricht: „Wobei soll ich es merken, daß ich dies Land besitzen werde?“ Ich habe oft gesagt, wenn GOtt eine Verheißung thut, so thut er auch ein Zeichen dazu; also auch hier, da er ihm hat verheißen, ein Kind zu geben, das dies Land sollte besitzen, gibt er ihm ein solch Zeichen, daß er dreierlei Thiere soll nehmen, und zertheilen und gegen einander legen. Es ist ein seltsam Zeichen, als uns dünkt, weil er nun lange todt ist; aber ihm ist es tröstlich gewesen, hat seinen Glauben damit bestätigt und gestärkt. Was aber dies Zeichen bedeute, legt der HErr selbst aus; wie wir hören werden.

14. Auf das erste: Diejenigen, so die Bibel oberflächlich lesen, halten es dafür, als rede Clauß Schmid mit Hans Wurst, ein Mensch mit dem andern; aber, wenn man recht hinein sieht, ist es groß und trefflich, daß die göttliche Majestät mit dem Menschen so redet, als mit einem Kinde. [Es] ist kein theurer Zeichen GOttes Gnade und Huld, denn wenn er sich hören läßt, als wiederum keine größere Plage ist, denn wenn er stille schweigt und sich nicht läßt hören. Niemand soll man solche Historien, daß GOtt mit den Vätern geredet hat, hoch halten, nicht um der Väter willen, die da Leute gewesen sind als wir; sondern darum, daß sie GOtt so hoch gehalten hat, daß er mit ihnen redete. Wenn Abram kein ander Lob hätte, wäre es genug, daß GOtt einmal oder zwei mit ihm geredet hätte. Nun siehst du, wie er ihn so hoch hebt, daß er so oft und mit viel Worten mit ihm redet und ihn tröstet, und läßt ihn doch daneben stecken in aller Fährlichkeit Leibes, Gutes und Ehre, dazu auch des Weibes. Derhalben die Historien nicht so zu überlaufen und für schlecht zu halten sind.

15. Also ist dies Zeichen auch an ihm selbst fast unnütz für den etlichen Geistern, denen die Schrift nicht schmeckt, aber sonst an ihm selbst fast reich. Denn da stärkt GOtt Abrams Glauben mit dem Zeichen, daß er nicht zweifele, er werde nicht allein einen Samen haben, sondern auch das ganze Land einnehmen; legt es dazu selbst aus. Wie weit sich aber die Figur und Deutung streckt, können wir nicht besser treffen, denn der Text
\quellseite{143}
selbst gibt. Denn der Text geht nicht damit um, was die dreijährige Kuh, Ziege und Widder bedeutete; aber so viel will er, daß ihn das Zeichen sollte so gewiß machen, nicht allein daß er den Samen haben sollte, sondern auch wissen, wie es ihnen gehen würde. Als sollte er sagen: Dies Land habe ich dir geredet, daran sollst du nicht zweifeln; aber so lange wird sich’s verziehen und der Glaube so wohl versucht werden, daß sich’s allerdings stellen wird, als werde nichts daraus. Darum verkündigt er ihm, daß sein Same fremd sein soll in fremden Landen, und übel gehandelt werden so lange, bis die Zeit komme, daß er das Land einnehme, nämlich, nach vierhundert Jahren, wenn niemand glauben würde, daß GOtt an sein Wort gedächte.

16. Denn so habe ich gesagt, daß allezeit GOttes Werk so gethan und geschickt ist, daß, wenn er sein Wort erfüllen will, [es] so närrisch angreift, als sollte es gar zurück gehen; als, hier verheißt er ihm so groß Ding, daß sein Same das Land soll inne haben. Da das sollte angehen, konnte seine Vernunft noch Sinn anders sehen und fühlen, denn daß es gar mit ihm sollte untergehen. Denn so ist 2 Mos. 1, 7. 10. geschrieben, als sich die Kinder Israel so mehrten, daß es wimmelte von Kindern, und einen Schein hatte, als sollte es nun wahr werden, kehrte sich’s eben um, daß ein König aufkam, ließ alle Kinder umbringen, schwächte und täubte das Volk, so daß schier alle Hoffnung aus war. Ja, da sie nun schon heraus aus Egypten gegangen waren [2 Mos. 14, 9. 10.] und an das rothe Meer kamen, jagten ihnen die Feinde nach, daß sie auf keiner Seite hinaus konnten, und gleich den Feinden und dem Tode mitten im Rachen staken. So handelt GOtt mit allem seinem Thun, Worten und Zeichen. Denn seine Werke sind nicht wie Menschenwerke, Jes. 55, 8, sondern gar widersinnig. Hier geht es also, daß, wenn etwas aufgehen soll, so geht es zuvor unter; und wiederum.

17. So lehren solche Historien alles, was der Glaube ist. Wer sich nun nicht hierein richten kann, der versteht nicht GOttes Werk, muß sich daran ärgern und zurück fallen. Wo nun da Juden gewesen sind, die nicht geglaubt haben, haben gedacht: O es ist längst aus, was sollte nun daraus werden? hätte er solches wollen thun, so ließe er uns nicht so lange und schwerlich plagen. Welcher ohne Zweifel viel gewesen sind, die GOttes Werk mit der Vernunft gemessen, nach den Augen und Sinnen gerichtet, und nicht haben können verstehen.

18. Um deswillen gibt nun der Herr des ein Zeichen, kommt zuvor und warnt als ein freundlicher Vater; als sollte er sagen: Siehe zu, so wird es gehen, daß du dich nicht ärgerst an meinem Werk oder Wort. Wie Christus auch that, als er Joh. 16, 1. spricht: „Solches habe ich zu euch geredet, auf daß ihr euch nicht ärgert.“ Will also sagen: Da habt ihr das Wort, haltet aber fest daran; denn wenn sich’s anheben wird, als solle es erfüllt werden, so wird es vor aller Leute Augen zurück gehen; seid aber getrost, und sehet nur zu, daß ihr nicht wanket, sondern euch an das Wort haltet, ich will es wohl ausführen, höher, denn ihr denken noch bitten könnet. So ist dies auch eine Warnung: Ob es gleich scheint, als habe ich meiner Verheißung vergessen, daß ich euch unterdrücken lasse, so will ich euch doch mit großer Gewalt und Wunderhand einsetzen, und euch an euren Feinden rächen.

19. Siehe, so ist es mit Christo auch gegangen, und geht noch täglich in seinem Reiche also. Da jedermann meinte, er sollte das Volk erlösen und König werden, so griff er es also an, daß das Volk mußte sagen: Nun ist es aus; [er] legte sich nieder, ließ sich an das Kreuz schlagen, und zurichten als den ärgsten Buben, daß sie ihn anspieen, und sprachen [Matth. 27, 42.]: Ja, ist das der König und Erlöser, der anderen Leuten hilft, und kann ihm selbst nicht helfen? Wir haben ihn geachtet als für den Allerverachtetsten, und nichts denn Schmerzen und Schwachheit an ihm gesehen, sagt Jesaias [Cap. 53, 3.], und Lucä am letzten, V. 21., sagten die zwei Jünger: „Wir hofften, er sollte Israel erlösen.“ Darum, da er in sein Reich wollte gehen, ging er also, daß sich alle Vernunft und Sinne mußten daran ärgern, auch die Apostel selbst konnten es nicht ermessen, wiewohl sie Christus zuvor so fleißig warnte.

20. So geht es allemal zu, in allen Historien. Wo GOtt wirkt, stellt er sich so närrisch dazu, daß [es] nach der Vernunft gar in Dreck fällt, aber nach dem Geiste geht es auf das allerstärkste auf. Dies ist die Warnung, so hier GOtt thut, nicht allein um Abrams, sondern seines Samens willen.

21. Denn Abram wird die Zusagung gar
\quellseite{144}
fleißig\footnote{Erlanger: feindlich. Wittenberger und Jenaer: fleißig.} gepredigt haben, darnach Isaac, und so fortan, als das ihr Evangelium war: „Wir sind Abrahams Same“ [Joh. 8, 33.], haben die Zusage. Abraham hat geglaubt, und ist dadurch vor Gott rechtfertig geachtet“ [Röm. 4, 3.], darum lasset uns auch glauben. So sagt Joseph zu ihnen, da sie schon in Egypten waren [1 Mos. 50, 24. 25.]: „GOtt wird euch heimsuchen, und aus diesem Lande führen, in ein Land, das er Abram, Isaac und Jakob geschworen hat. Darum, wenn euch GOtt heimsuchen wird, so führet meine Gebeine von dannen.“ Das Evangelium haben sie wohl gewußt und getrieben; aber da es sollte angehen, stellte sich’s so widersinnisch, daß da noth war, die Frommen zu trösten und stärken mit diesen Worten, so GOtt hier zuvor sagt. Da dienen nun die Exempel zu durch und durch, daß wir gerüstet seien: Wenn wir das Evangelium haben, und es angehen soll und helfen, soll sich’s finden, als habe er unser gar vergessen; wenn das Leben soll angehen, wird der Tod mit aller Macht herfallen.

22. Was ist nun hier das Zeichen? Die Thiere, so er theilt, und die Tauben, ist sein Same und Kinder in Egypten. Die Raben und Vögel, die es fressen wollen, sind die Egypter, der König mit seinen Tyrannen, die Abrams Fleisch und Blut haben wollen verschlingen und verzehren. Er aber scheucht sie davon, nicht um seiner Heiligkeit, sondern um der Verheißung willen. Denn er muß da stehen an Statt der göttlichen Verheißung. Denn das Wort, das GOtt sagt, das ist, die Verheißung, hat sie weggetrieben, daß [es] so viel gesagt sei: Darum, daß GOtt an seine Verheißung dachte, dem Abram gethan, darum half er den Juden und erlöste sie. Also ziehen es die Propheten selbst oft an, und sonderlich Moses [2 Mos. 2, 24.], daß GOtt spricht: „Ich habe gedacht an meinen Bund mit euren Vätern, Abraham, Isaac“ etc., ich bin wiedergekommen, euch zu helfen. Das heißt die Vögel hinweg gescheuchet, daß er um seinetwillen, das ist um der Verheißung willen, ihm versprochen, will seinem Samen helfen. Also, wo in der Schrift angezogen wird, daß GOtt Abrams und der Väter gedenkt, wird alles um seiner Verheißung willen angezogen.

23. Hier ist aber eine Frage: wie die Rechnung überein trifft von den vierhundert Jahren, davon er hier sagt? Denn wenn man es rechnet, findet sich’s, daß sie fast bei zweihundert Jahren und irgend sechzehn in Egypten gewesen sind; aber die Verfolgung hat gewährt fünfzig oder sechzig Jahre; also, daß von der Zeit an, da Jakob in Egypten zog, bis an Moses, der sie wieder ausführte, sind zweihundert und sechzehn Jahre. Warum sagt er denn hier von vierhundert Jahren? Darum, daß er Abram mit einzeucht mit seinem Samen. Denn von Abram bis auf den Ausgang aus Egypten rechnet man vierhundert und dreißig Jahre. Als sollte er sagen: Es soll nicht eher geschehen, denn bis vierhundert Jahre um seien. Daß er aber hernach sagt: „Sie sollen nach vier Mannes Zeiten wieder hieher kommen“, das ist, vier Menschen Gedenken, das sind fast zweihundert Jahre. Denn eines Mannes Gedenken sind nicht viel über fünfzig Jahre. Das befehle ich den Gelehrten weiter zu handeln.

24. Daß aber hernach folgt, daß ein groß Schrecken auf Abram fiel, da die Sonne untergegangen war etc., ist eben das vorige Zeichen. [Er] gibt ihm zweierlei Zeichen, wie der König Pharao zweierlei Träume hatte, daß doch einerlei deutete, wie hernach [Cap. 41, 1. 5.] folgen wird; dazu, daß er je gewiß und sicher wäre, daß es so geschehen sollte. GOtt hat große Sorge für das zukünftige Volk, daß er es in dem Glauben behielte, auf daß sie wüßten, wie es viel anders sollte angesehen sein. So bedeutet nun der feurige Ofen aber die Egypter; daher GOtt hin und wieder in der Schrift [5 Mos. 4, 20. Jer. 11, 4.] sagt: „Ich habe euch aus dem eisernen Feuerofen Egypti geführt“, darinnen ihr gebacken wurdet. Daß aber Abram ein Schrecken ankommt, und eine große Finsterniß, bedeutet, daß sein Blut und Fleisch so gedrungen, gezwungen, und verfolgt sollte\footnote{Wittenberger und Erlanger: sollten.} werden in Egyptenlande; daß müßte er eine Figur sein. Der feurige Brand oder Fackel ist Moses, oder der Engel, der sie ausführte. Das ist das Zeichen oder beide Zeichen, die GOtt gibt zu sicherem Glauben.

25. Also auch wir, wenn wir das Wort haben, haben wir auch Zeichen daneben. Und wie ich von Abram gesagt habe, daß wir an ihm spüren, daß GOtt groß mit ihm gethan habe, und für einen
\quellseite{145}
großen Heiligen hatte, daß er mit ihm geredet hat: so ist es mit uns auch, wir sind auch in die Gnade gekommen, daß wir GOttes Wort hören. Wie jämmerlicher Schade [und] Würgen\footnote{„Würgen“ fehlt in der Wittenberger.} in aller Welt ist, weil das Wort GOttes nicht geht, kann niemand ermessen. Wiederum, den Reichthum, daß man sein Wort hat, kann niemand genug preisen, seine Zunge aussprechen. Es ist ein gnädiges Zeichen, wenn sich die Majestät hören läßt. Nun läßt er sich jetzt viel reichlicher hören, denn zu Abram. Das ist aber der Mangel, Abram hat [es] mit Ohren gefaßt, aber auch verstanden; wir hören es, können es aber nicht. Die es aber fassen mit dem Herzen, die sind eben so würdig, als Abram; denn wie das Wort einerlei ist, so ist auch die Würde einerlei.

\section{Das sechzehnte Capitel}

\begin{schriftzitat}
V. 1—16. Sarai, Abrams Weib, gebar ihm nichts; sie hatte aber eine egyptische Magd, die hieß Hagar, und sprach zu Abram: Siehe, der HErr hat mich verschlossen, daß ich nichts gebären kann. Lieber, lege dich zu meiner Magd, ob ich doch vielleicht aus ihr mich bauen möge, mehr denn aus mir. Abram gehorchte der Stimme Sarai. Da nahm Sarai, Abrams Weib, ihre egyptische Magd, Hagar, und gab sie Abram ihrem Manne zum Weibe, nachdem sie zehn Jahr im Lande Canaan gewohnt hatten. Und er legte sich zu Hagar, die ward schwanger. Als sie nun sahe, daß sie schwanger war, achtete sie ihre Frau gering vor ihr. Da sprach Sarai zu Abram: Ich muß Unrecht leiden um deinetwillen; ich habe meine Magd dir zugelegt, nun sie aber siehet, daß sie schwanger worden ist, muß ich geringe geachtet sein vor ihr; der HErr sei Richter zwischen mir und dir. Abram aber sprach zu Sarai: Siehe, deine Magd ist unter deiner Gewalt, thue mit ihr, wie dir’s gefällt.\footnote{Der folgende Text bis zum Ende des Capitels fehlt in der Wittenberger.} Da sie nun Sarai wollte demüthigen, flohe sie vor ihr. Aber der Engel des HErrn fand sie bei einem Wasserbrunnen in der Wüste, nämlich, bei dem Brunnen am Wege zu Sur, der sprach zu ihr: Hagar, Sarai Magd, wo kommst du her? und\footnote{Die Worte: „wo ... und“ fehlen in der Erlanger.} wo willst du hin? Sie sprach: Ich bin von meiner Frau Sarai geflohen. Und der Engel des HErrn sprach zu ihr: Kehre um wieder zu deiner Frau, und demüthige dich unter ihre Hand. Und der Engel des HErrn sprach zu ihr: Siehe, ich will deinen Samen also mehren, daß er vor großer Menge nicht soll gezählt werden. Weiter sprach der Engel des HErrn zu ihr: Siehe, du bist schwanger worden, und wirst einen Sohn gebären, des Namen sollt du Ismael heißen, darum, daß der HErr deine Armseligkeit erhöret hat. Der wird ein wilder Mensch sein, seine Hand wider jedermann, und jedermanns Hand wider ihn, und wird gegen allen seinen Brüdern wohnen. Und sie hieß den Namen des HErrn, der mit ihr redete: Du GOtt siehest mich; denn sie sprach: Gewißlich, hier habe ich den Rücken gesehen deß, der mich sah.\footnote{So die Erlanger. Jenaer: siehet.} Darum hieß sie den Brunnen einen Brunnen des Lebendigen, der mich sah, der da ist zwischen Kades und Bared. Und Hagar gebar Abram einen Sohn, und Abram hieß den Sohn, den ihm Hagar gebar, Ismael. Und Abram war sechs und achtzig Jahr alt, da ihm Hagar den Ismael gebar.
\end{schriftzitat}

1. Welche schlechte Historien beschreibt er doch? Hat er denn nichts anders zu schreiben gehabt? Ich habe vor gesagt, man soll die Schrift, GOttes Werk und Wort, und GOttes Gericht mit Furcht lesen, und nicht achten als andere Historien, sondern immer Acht haben auf den, der das Wort redet und das Werk thut, da findest du allewege GOtt, es scheine so gering es wolle.

2. Bisher haben wir nun gehört, wie GOtt Abram verheißen hat, daß von seinem Blut und Fleisch ein groß Volk sollte kommen, und das ganze Land einnehmen. Das hat er gepredigt, Sarai hat es gehört, welche nun fünf und siebenzig Jahr alt war, als Ismael geboren ward:


\quellseite{146}
indeß hat sie geharrt zwei, drei, vier, fünf, sechs Jahre, ob sie GOtt noch fruchtbar machen wollte; ist nichts daraus worden. Da sie nun sieht, daß sie keine Frucht trägt, wie der Text sagt, gedachte sie so: Mein Mann hat die Zusagung, daß er gewiß ist, daß er so reich soll werden vom Samen aus seinem eigenen Leibe; es will aber nichts daraus werden, vielleicht will es GOtt von mir nicht haben; wirft sich herunter, denkt also: Ich bin das Weib nicht, GOtt sieht mich nicht dafür an, daß ich des Samens Mutter sein soll; verzeihet sich der Ehre mit großer Demuth, dadurch sie hernach wieder höchlich geehrt wird.

3. Was thut sie denn? Sie wollte auch gerne, daß die Verheißung erfüllt würde, ist ebenso wohl in dem Glauben gestanden als Abram, auch dadurch rechtfertig worden, gibt GOtt auch die Ehre, daß er wahrhaftig sei, läßt ihr an dem genügen, des andern verzeihet sie sich. Alles sollte sie sagen: Soll ich es nicht sein, so sei es eine andere; doch will sie keine Fremde dazu nehmen, sondern nimmt ihre Magd. Es sind aber da zur Zeit die Knechte und Mägde leibeigen gewesen, wie das Vieh, also, daß, was sie zeugeten, das war der Herren und Frauen eigen, und wenn sie wollten, mochten sie sie ausstoßen, und die Kinder behalten. Welches jetzt aufgehoben ist.

4. So dachte nun Sarai: Ich will keine Fremde, sondern die nehmen, die ich am nächsten habe, die doch sonst mir Kinder trüge, daß sie dennoch mein sind; wie sie spricht: „Ob ich mich aus ihr erbauen möchte.“ Bauen heißt im Alten Testamente Kinder zeugen, wie wir in diesem Buche mehr hören werden; und im andern Buche Mosis Cap. 1, 21. steht so geschrieben: „GOtt bauete den hebräischen Weibern Häuser“, das ist, er gab ihnen Früchte, daß sie viel Kinder trugen, machte ihnen die Häuser voll.

5. Hier ist abermal ein großer Glaube, beide Abrams und Sarai. Er ist in das Land gekommen, da er fünf und siebenzig Jahr alt war, und Ismael ist geboren im zehnten Jahr hernach. Da sind sie neun ganzer Jahre im Lande, dazu kommt die letzte Verheißung ein Jahr oder vier oder sechs zuvor, daß sie allenthalben reichlich versucht sind. [Sie] hätten wohl mögen denken: Wann will es denn werden? noch haben sie immer fest daran gehalten, wie lange GOtt auch verzogen hat. Doch hätte sie auch gern geeilt, wie die Weiber alle thun, hätte gern bald ein Kind gehabt, darum treibt sie es, gibt ihm ihre Magd; solches bewilligt Abram. Das scheint allemal für Narrwerk denen, so die Schrift überhin laufen, als sei Abram ein Mann gewesen, der um Lust und Vorwitz willen sich zur Magd legte. Aber GOtt hat ihm sein zuvor kommen. Hätte er Lust zur Büberei gehabt, so hätte er wohl eine andere gefunden, der Weiben eine zu Hebron; geht aber nicht weiter, denn Sarai haben will, und ihn bittet. Daher klar ist, daß er nicht gegeizt hat nach einer andern, und nicht mehr sucht, denn die Frucht.

6. Hier wäre auch zu reden von dem Stücke: ob ein Mann auch mehr denn ein Weib haben möge. Es ist genug gesagt. Was hier geschehen und beschrieben ist, muß man lassen bleiben als GOttes Geschichte; [es] gebühret niemand, solche Historie zu strafen; darum muß Abram hierinne nicht gesündigt haben. Weiter haben wir gehört, daß Abram ein rechter, ja, vollkommener Christ gewesen ist, auf das altevangelische gelebt im Geiste GOttes und Glauben; darum müssen wir sein Leben so lassen gehen, daß es ein Exempel sei, darnach zu thun, so sichs begäbe im selben Glauben. So fragt man nun: Weil es verboten, daß er mehr denn ein Weib hätte, wie wollen wir ihn denn verantworten, daß er ein Christ geblieben sei?

7. Zum ersten, wenn man wil einfältig dazu reden, spricht man also: es sei durch das Evangelium aufgehoben, da Christus spricht im Matthäo [Cap. 19, 8.]: „Von Anfang war es nicht also“; item, Paulus [1 Cor. 7, 2.]: „Ein jeglicher habe sein Weib, und eine jegliche ihren Mann“, daß man möchte sagen, Abram hat etliche äußerliche Werke gethan, die nun aufgehoben sind.

8. Aber das wird nicht genug thun, denn wir müssen uns so bewahren, daß er redlich verantwortet sei. Denn das ist je wahr, daß alles, so wir finden im alten Testamente, von den Vätern äußerlich gethan, frei sein solle, nicht verboten. Als, die Beschneidung ist aufgehoben; aber nicht also, daß es Sünde wäre, wenn man es thäte, sondern frei, weder Sünde noch wohlgethan. Also, das Osterlamm essen, und dergleichen alles, ohne daß man es nicht in dem Namen thue, daß man dadurch wolle rechtfertig
\quellseite{147}
werden, wie es die Väter auch nicht darum gethan haben. So muß auch unter andern Exempeln der Väter mitgehen, daß sie viel Weiber genommen haben, daß es auch frei sei gewesen. Und Summa Summarum: Abram hat es nicht so weit geführt, als andere Patriarchen, nimmt die Magd nicht eher, denn es ihn sein Weib heißt und bittet. Andere Exempel werden weiter bringen; dies könnte noch wohl so hingehen, als ein Sonderliches vor andern Vätern.

9. So sage ich nun und schließe: Wo nicht Christen, sondern heidnische Leute sind, wollte ich noch, daß man dem Gesetze nach hätte, vom Scheiden, daß einer ein Weib möchte von sich thun, und eine andere nehmen. Christus hat es ja aufgehoben; sagt also im Matthäo [Cap. 19, 8. 9.]: „Moses hat euch erlaubt zu scheiden von euren Weibern, von eures Herzens Härtigkeit wegen; von Anbeginn aber ist es nicht also gewesen. Ich sage aber euch: wer sich von seinem\footnote{Erlanger: einem.} Weibe scheidet, es sei denn um Hurerei willen, und freiet eine andere, der bricht die Ehe.“ Desgleichen auch Paulus zum Corinthiern [1. Ep. 7, 10. 11.]: „Den Ehelichen gebiete nicht ich, sondern der HErr, daß das Weib sich nicht scheiden lasse von dem Manne. So sie sich aber scheiden lässet, daß sie ohne Ehe bleibe, oder sich wieder mit ihm versöhne, und daß der Mann das Weib nicht von sich lasse.“ Aber die Christum nicht hören, wäre noch wohl so gut, daß Mosis Gesetz ginge, ehe man das leiden müßte, daß zwei Eheleute keine gute Stunde bei einander hätten. Aber dabei müßte man ihnen sagen, daß sie nimmer Christen wären, sondern im heidnischen Regimente. Bist du aber ein Christ, mußt du dich nicht scheiden.

10. Aber es ist nicht verboten, daß ein Mann nicht mehr denn ein Weib dürfte haben. Ich könnte es noch heute nicht wehren, aber rathen wollte ich es nicht; denn damit bleibt gleichwohl noch, daß man sich nicht scheiden solle, sondern seinem Weibe anhangen. Darum bringen die vorigen Sprüche solches nicht. Doch wollte ich es nicht aufbringen, sondern darum sage ich es, wenn es zur Schärfe käme, daß man recht müßte zu antworten, daß man die Väter nicht verwerfe, als hätte sich es nicht geziemet, zu thun; wie die Manichäi sagten.

11. So haben wir nun gesehen die große Demuth der heiligen Mutter Sarai, daß sie sich der Ehre und des Kindes verzeihet, und einem andern Weibe läßt den Samen und Segen, Abram verheißen. Er redet nicht viel Worte von dem heiligen Weibe; aber das Exempel gibt genug Anzeige, daß sie von reichem Geiste und hoch begnadet ist gewesen; darum sie die Schrift an mehr Orten anzeucht. [Sie] hat sich nicht erhoben, sondern geniedriget, darum ist sie auch erhöht worden [Matth. 23, 12.].

12. Nun die Historia, wiewohl sie [den] Grund legt eines heiligen Lebens, ist sie doch nicht allein darum geschrieben; Paulus [Gal. 4, 22. ff.] hat es gar meisterlich ausgestrichen. Es sind auch [die] Worte so gesetzt, muß etwas dahinter sein, nicht allein die Historie. Denn er hätte sonst viel Dinge können anziehen und des schweigen; läßt es aber stehen, sagt, wie die Magd stolz ist worden, als sie sahe, daß sie schwanger ging, daß sie die Frau mußte ausstoßen; und wie sie wieder zu ihr kam und den Sohn Ismael gebar, davon der Engel verkündigte, wie er ein wilder Mensch würde sein und streitbar, und sich setzen wider alle seine Brüder. Daher man es noch dafür hält, daß der Türke von diesem Ismael herkomme, der gewohnt hat in Arabia, daher sie auch Ursprung haben, daß man Hagarenos heißt die Saracener.

13. Darum wollen wir die Figur handeln, wie es die Historie an ihr selbst bringt. So hast du gehört, daß die Schrift allenthalben vorbildet die zwei Predigten, des Gesetzes und Evangelii; nach denselbigen zwei Predigten hält sie auch vor zweierlei Volk, des Gesetzes und der Gnade. In dem Ismael wird angezeigt das Volk, so das Gesetz hält; darum führt er auch den Namen, Art und Weise, so solche Leute führen, daß er ein Vorbild ist aller Gesetz- und Werkheiligen. Derhalben soll es uns auch bekannt sein, weil es auch der Apostel Paulus hat sonderlich ausgestrichen, was [das] für Leute sind, die durch Werke, das ist, ohne Gnade wollen fromm werden.

14. Am ersten, wird er nicht geboren von der rechten Hausmutter, die das Regiment hat, sondern von der Magd; wiewohl sie auch ehelich ist mit Abram, doch ist es ein großer Unterschied. Denn der andere Sohn, Isaac, ist nicht von der Magd, noch leibeigen, sondern von der Freien, wie Paulus sagt zum Galatern [Cap. 4, 22.], die die rechte Frau ist und regieret im Hause.


\quellseite{148}
Und, daß das Größte ist, wird er geboren vom Weibe, da die Verheißung GOttes auf deutet. Das ist trefflich; und St. Paulus [Gal. 4, 23.] führt es mit großer Gewalt, daß Ismael wohl geboren ist aus Abrahams Blut und Fleisch; aber es ist geschehen ohne Gottes Wort; Isaac aber durch GOttes Wort.

15. Denn da Sarai sagte: „Siehe, ich kann keine Frucht tragen. Lieber, nimm doch meine Magd, ob ich von ihr möchte Kinder zeugen“, da ist kein GOttes Wort noch Befehl, sondern nicht mehr, denn ein Werk der Liebe, daß sie weicht, läßt ihre Ehre fahren, die sie gerne gehabt hätte; so ist er ihr zu Willen. Dazu ist sie, die Hagar, ein jung Weib, daß die Geburt gar natürlich ist. Aber hernach [Cap. 17, 16.] kommt GOttes Wort zu Sarai und spricht: Du sollst einen Sohn haben, den will ich segnen. Das war ein GOttes Wunder. Das muß man nun merken, daß Paulus gewaltig führt, wie Ismael geboren ist durch die Magd, ohne Gottes Verheißung; Isaac von Sarai, durch die Verheißung. Nun haben wir gesagt, daß gar groß zu pochen ist, da GOttes Wort ist; wo das ist, da ist es köstlich Ding; denn es ist ein ewig Ding. Wenn das Kleidung und Schah nicht wäre, so wäre Ismael nicht mehr, denn ein anderer.

16. So geht es auch im neuen Testamente zu, und von Anfang der Welt; im geistlichen Regimente wird alles geboren entweder von GOttes Gesetz oder durch die Gnade. Da das Siegesvolk regiert, da sind eitel Ismael, empfangen auch von GOtt rechten Samen, wie Ismael von Abram empfangen wird; hören das Wort und Predigt, wachsen mehr davon denn die rechten Christen, aber werden nicht rechtschaffen, fallen auf Werke, greifen es mit dem freien Willen an, verlassen sich darauf, haben GOttes Gebot zuvor. Da gehen sie in einem scheinbarlichen Leben vor der Welt, daß [es] niemand tadeln kann. Das sind alles der Magd Kinder. Denn es ist noch kein GOttes Wort der Verheißung da, sondern allein Fleisch und Blut, natürlich Ding; GOttes Gnade wirkt nicht da, [es] ist kein Glaube noch Geist da, dadurch sie vom Gesetz frei werden, darum werden nichts denn Ismael daraus.

17. Ismael aber heißt auf deutsch ein Hörer GOttes. Den Titel hat das Volk alles, sonderlich die Juden, die sind damit fein abgemalt, haben GOttes Wort gehört, wie Christus wird kommen; des hatten sie beide Ohren voll, waren hoch begnadet mit vielen Propheten, haben das Gesetz Mosis getrieben durch und durch mit Fleiß, sind das Volk, das vor aller Welt GOttes Gesetz alle Tage hatte; darum haben sie wohl solchen Namen, daß sie viel hören von GOtt, sehen ihn aber nicht. Gleich als wenn einer auf dem Wege geht, der hört und sieht wohl, was vor ihm ist, aber im trachen\footnote{Wittenberger und Jenaer: starcken.} Gange kann er nicht sehen, was hinter ihm ist. So war es hier auch. Christus sollte kommen nach Abram; das hörten sie wohl, sahen es aber nicht; ja, da er gegenwärtig war, und noch immer, bleiben sie nur Ismael, hören immerdar.

18. Der Art sind alle, so mit Werken umgehen, kennen und sehen GOtt nicht, ist ihnen verborgen, hören allein von ihm, und stehen im Wahne und Vorsatz, so zu leben in ihren Werken, bis GOtt kommt und gibt Gnade zu. Denen ist er noch nicht gekommen, ja, kommt nimmermehr. Dieweil hören sie immer, er wird kommen; also, daß er immer hinter ihnen bleibt, wie den Juden. So ist das Volk gar meisterlich vorgemalt mit Namen und allen Zügen. Wir Christen aber haben ihn vor Augen, und kennen ihn, sagen also: Mit Werken traue ich nichts zu erlangen; ich weiß aber, daß Christus da ist, der mir alles erworben hat, sehe ihn vor mir im Spiegel, bin nicht mehr ein Ismael. Das ist ein Stück dieser Figur.

19. Die Mutter aber heißt Hagar. Hagar heißt ein Gast; daher Hagareni die Gäste. So sagt nun Paulus [Gal. 4, 30.], alle, die mit Werken umgehen, sind freilich Gäste und Fremdlinge, nicht Kinder im Hause. Wer ein Gast ist, deß sind weder Haus noch Güter eigen, isset und trinkt, und geht wieder davon. So sind alle zeitlichen Werkheiligen; ihr schön Leben bezahlt GOtt mit zeitlichen Gütern; als er dem jüdischen Volke das Land Canaan gab, hob sie hoch, daß sie ein berühmt und mächtig Volk waren; wie es auch war. Da wir äußerlich fromm waren, gab er uns den Bauch voll zu fressen, groß Gut und Ehre vor der Welt. Solchen Heiligen, die mit zeitlichem, vergänglichem Wesen umgehen, gebührt auch zeitlicher Lohn. Der Christen Lohn aber der Gerechtigkeit ist ewig, wie der Prophet [Jes. 32, 17.] sagt.
\quellseite{149}
Wenn wir nun sterben, wird es erst recht angehen: da kommt der Tod hinweg, und dafür eitel Leben und Seligkeit. Auf Erden haben sie nicht den Lohn, sondern alles Unglück, Schmach und Schande. Darum sind wir nicht mehr Gäste und\footnote{Erlanger: noch.} Fremdlinge, wie Paulus [Eph. 2, 19.] sagt, sondern Bürger mit den Heiligen und GOttes Hausgenossen, ja, die Kinder und das Erbe.

20. Die Werkheiligen aber haben auch wohl mit uns den Namen, daß sie Abrams Kinder und Christen heißen, und bleiben doch in zeitlicher Heiligkeit, die nicht ewig bleibt. Der Glaube aber, so GOttes Wort faßt, ist ein ewig Ding, nämlich an Christum, der nicht stirbt. So heißt nun der Name Hagar den ganzen Haufen der Werkheiligen; die sind Gäste, bleiben nicht endlich bei uns in der Christenheit, sondern fallen auf Werke, so da vergehen. Wir aber hängen uns an Christum, der da ewig ist. Weiter sagt der Text also:

\begin{schriftzitat}
V. 12. Siehe, er wird ein wilder Mensch sein, seine Hand wider jedermann, und jedermanns Hand wider ihn, und wird gegen allen seinen Brüdern wohnen.
\end{schriftzitat}

21. Das ist die rechte Art und Farbe, so diese Heiligen führen, die nicht im Glauben wandeln. Je schöneren Wandel sie führen, je ärger sie sind; boshaftiger, stolzer und unleidlicher Volk findet man nicht auf Erden, auch kein frecher, roher und freveler. Denn sie lassen sich dünken, sie sind rein, darum können sie niemand leiden, noch Barmherzigkeit mit andern haben, können es nicht lassen, müssen jedermann urteilen und richten; allerwenigst aber können sie nicht leiden, daß man ihr Ding nicht gelten lasse und hoch hebe.

22. Also haben wir auch gethan, haben gesagt: Wir sind geistlich, der andere Haufe ist die Welt; sind so aufgeblasene Leute worden, daß man ehe einen Kaiser hätte eingetrieben, denn einen schäbichten Mönch. Wenn man der einen antastete, so hatte\footnote{Jenaer: hat.} man Gott ein Auge ausgestochen; wenn es sich auch jemand unterstanden hat, den haben sie verbannt bis unter die Hölle. Das ist es, daß der Text nennt, daß sie frech und ungehorsam sind, wollen recht haben und müssen recht haben; man lache oder weine, bitte oder bräue, [so] bleiben sie mit dem Kopf härter denn kein Ambos, wie Hiob [Cap. 41, 15.] sagt, darauf wenn man gleich lange schmiedet, wird er nicht weich, macht eher ander Eisen weich.

23. So sagt Christus im Evangelio [Matth. 11, 18. 19.]: „Johannes ist gekommen, aß nicht und trank nicht, so sagen sie: Er hat den Teufel. Des Menschen Sohn ist gekommen, isset und trinket, so sagen sie: Siehe, wie ist er ein Fresser und ein Weinsäufer.“ Wie man es mit ihnen macht, kann man ihnen nicht recht thun: „Pfeift man, so tanzen sie nicht; klagt man, so weinen sie nicht“ [Matth. 11, 17.]. [Es] gilt nichts, denn was sie gut dünkt, und was ihnen wohlgefällt; wo nicht, so soll es GOtt nicht gefallen. Das sind rechte Ismael, Unbändige, Halsstarrige und Widerspenstige; ein Volk, das nicht zu lenken noch zu beugen ist; wie man sie mit Gottes Wort zwingen will, müssen sie sich heraus zu drehen. Solche muß man unter sich leiden, [wir] können\footnote{Erlanger: könnten.} nicht dawider, daß sie den Namen haben, den wir; wie Abram den Ismael im Hause leiden muß.

24. Weiter: „Seine Hand wider jedermann, und jedermanns Hand wider ihn.“ Man sagt noch, daß daselbige Volk, von Ismael hergekommen, ein wild, untreulich Volk sei, das sich nur mit Raube ernährt, reißt zu sich, was es ergreift. Solches wird nun Ismael erfüllt haben. Dem folgen diese Unseren auch nach. Die Hände sind die Werke, die müssen andere sein denn der Christen. Was soll denn daraus werden? Es sind viel weltkluge Leute, die so sagen: Man sollte das Evangelium predigen, daß kein Aufruhr und Unlust würde in der Welt, sonst wird das Volk irre. Aber das haben wir eben unsern Predigern zu danken. Warum haben sie nicht also das Evangelium vorgetragen, daß man wisse, was recht sei oder nicht? Ist doch geboten, daß jedermann zuhöre, und wisse zu urteilen; [sie] haben uns eben das Recht genommen. Weil sie das erhalten, so ist kein Aufruhr. Wenn man aber dem Urtheil nachgehen soll, ob der Papst und seine Lehre recht sei, da hat GOtt Hader eingesetzt, daß [es] nicht anders gehen kann noch soll.

25. Darum darf niemand dessen nach Ruhen: Soll das Evangelium angehen, so ist es recht, daß man sich getrost dawider legt. Wir, die
\quellseite{150}
mit dem Himmel leben, müssen uns darein ergeben, und fröhlich deß erwägen, daß seine Hand wider uns ist; so wird er müssen leiden, daß unsere Hand auch wider ihn ist; [wir] müssen den Hals daran setzen und nicht schweigen, sondern immer widerstreben. So ist er frevel und frech, läßt ihn nicht sagen, wird zornig, und würgt uns zuletzt; so sterben wir in GOttes Namen um GOttes Wortes willen. Also soll sich niemand entsetzen, sondern noch fröhlich sein, ob wir sehen, wie schändlich das Evangelium verfolgt wird. Denn es ist so zuvor gesagt, daß wir uns deß trösten und [es] frisch darauf wagen. Man muß hier im Heer liegen, da wird kein Mensch Richter sein, weder Papst noch concilia, sondern Gott will selbst Richter sein, sagt er [Ps. 82, 1.], darum wird nichts daraus, daß die Sache gerichtet werde; [man] thue was man wolle, so wird doch Himmel auf seiner Seite stärker sein und schließen.

26. Es ist nicht wohl möglich, daß ein Concilium sei, das einträchtiglich schließe, ohne Fleiß und Sauerteig. Ich habe noch nie kein Concilium unter allen gesehen, da der Heilige Geist regierte. Es ist wohl in etlichen ein Stück oder zwei christlich gewesen, aber daß es ganz rein wäre, habe ich nicht gelesen. Darum ist es nichts, daß man sich noch berufen wolle auf ein Concilium, es wird doch nichts Gutes ausgerichtet. Es ging genau zu mit dem Concilio zu Jerusalem, daß es rein abging. [Es] waren nur drei Männer da, die es rein erhielten, alle anderen Christen schlossen straks wider Christum; da hielt niemand denn Petrus, Paulus und Barnabas, noch ließen sie etwas zu; wären die nicht gewesen, so wäre es alles verfälscht worden. Wer hat uns verblendet, [daß] das, was wir setzen und thun, müsse recht sein?

27. Es wird nimmer kein Concilium, da die Leute zusammen stimmen sollten aus dem Heiligen Geist. Das läßt auch GOtt eben darum geschehen, daß er selbst will Richter sein, und nicht leiden, daß Menschen richten. Darum befiehlt er jedermann, daß er wisse, was er glaube. Es wird nicht helfen, daß du sagest, du hast so gehört oder gepredigt; da fragt der Teufel nichts nach. Aber wenn du GOttes Wort hast, kannst du sagen: Da habe ich das Wort, was darf ich weiter fragen, was die Concilia sagen? So muß er zurück.

28. Weil denn nun GOtt haben will, als wir sehen, daß sein Wort keinen Richter leiden soll denn ihn; so soll niemand hier scheiden, kann auch keine Ruhe sein. Und, da GOtt für sei, wenn es stille und Ruhe würde, so wäre das Evangelium aus. Es muß rumoren, wo es kommt; thut es [das] nicht, so ist es nicht recht. Darum Christus [Luc. 12, 49. 51.] sagt: „Ich bin gekommen, ein Feuer anzuzünden auf Erden; was wollte ich lieber, denn es wäre schon angezündet? Meinet ihr, daß ich hergekommen bin, Frieden zu geben auf Erden? Ja sage ich Nein zu, sondern Zwietracht.“ Soll das gehen, so muß sich Himmel regen mit der Hand, und jene wieder. Da scheidet uns der Heilige Geist. Wenn Menschen scheiden wollen, wird es nur ärger.

29. Das rede ich alles darum, daß man gewohne, was Gottes Wortes Art ist, daß man nicht achte, als sei es ein still und ruhig Ding, wie es die höchste Vernunft vorschlägt, welche es so hart ansicht, daß [es] so rumorisch zugeht; wollte der Sache so gern helfen, daß doch Einigkeit würde. Gehen kann ich es lassen, daß man darinne handele; aber wenn es dazu kommt, wird Himmel immer seine Hand wider uns legen, wenn wir uns schon erbieten auf das allerfreundlichste. So muß man sie lassen unter einander geneigt bleiben, etliche Himmel, etliche Staat, und immer Hand gegen Hand gehen lassen. Weil wir denn solches sehen in der Schrift vorgebildet, wie es geht, sollen wir getrost sein, GOtt danken, der uns damit versichert, daß es recht ist. GOtt ist wohl so klug, daß er solche weisen Leute bald zu Narren macht; wie er auch thun wird. [Er] kann wohl durch verachtete geringe Leute, zu denen sichs niemand versieht, sein Wort lassen predigen. Seine Gerichte sind hoch, verborgen und unbegreiflich. Die Welt will allein nach der Person sehen; ist sie weise und scheint heilig, so will jedermann zufallen, und wiederum, [sie] hanget nicht bloß am Wort GOttes, wie die Gläubigen thun. So bleibt nun dies Himmel und seines Haufens Amt und Art, daß sie frech und stolz sind, jedermann richten, unbändig, unüberwindlich sind, nicht aufhören zu rumoren wider das Evangelium und die rechten Christen.

30. Das andere Volk ist Gnadenvolk; werden auch geboren von Abram, aber haben auch die Verheißung dazu, und heißen filii promissionis [Röm. 9, 8.], über welche GOttes Wort geht, und sagt: Du bist es. Das sind die, so nach dem heiligen Evangelio mit dem Geist verzeichnet
\quellseite{151}
werden, wie Paulus Röm. 8, 16. schreibt, vertrauen auf Christum, nicht auf Werke, halten an ihrem Heilande; die werden nicht durch das Fleisch geboren, sondern durch das Wort und Geist\footnote{Statt der vorhergehenden Worte hat die Erlanger: „die werden nicht geboren durch die Predigt“, man kann \&c.} Joh. 3, 5. 6.], man kann auch nicht sehen, wer sie sind, sondern GOtt hat sie beschlossen, daß es die\footnote{„es die“ fehlt in der Erlanger.} sind, welche er will. Man predigt wohl öffentlich jedermann; aber er weiß, welche er rühren will, und neu geboren machen. Das sind, die nicht an Werke gebunden sind, sondern: sein gelenkt, nicht stolz noch frech, können alles leiden, verdammen niemand, haben ihr Gut und Schatz im Himmel; wie ihre Gerechtigkeit ewig ist, so ist auch ihr Lohn ewig. Die sind nicht von Abrams Blut und Fleisch allein geboren, sondern von der Frau und Freien, davon Paulus [Gal. 4, 23.] sagt. 31. Die Freiheit ist, die Christus erworben hat und uns befrei[e]t, daß kein Werk noch Gesetz das Gewissen bindet. Denn Christus hat für alles genug gethan, durch den Heiligen Geist haben wir es gar, sind hinfort nichts schuldig, denn daß wir einander lieben und dienen. Das ist die rechte Mutter, und das sind\footnote{„das sind“ fehlt in der Erlanger.} die rechten Kinder. Denn Sarai, wie wir hernach hören werden, heißt eine Frau, die da regiert im Hause, der alles unterthänig ist, wie der Mann ein Herr heißt. So will nun Paulus sagen: Jene ist eine Magd und Gast, die ist Frau im Hause, da hat sie den Namen nach.

32. Das sind die zweierlei Leute, doch unter Einem Regiment, das alles fließt von GOttes Wort. Denn es ist nicht von groben, frechen Leuten gesagt, die unter das weltliche Schwert gehören, sondern von den Allerbesten, Frommen und Weisesten auf Erden; vor GOtt aber gilt es nicht, es sei denn anderweit aus dem Heiligen Geist geboren, und werde ein neuer Mensch. Es ist vor Augen, daß manch frommer Mann, Mönch und Nonne so leben, daß sie niemand tadeln kann. Von solchen Leuten sagt er hier, das sind die Rechtsschuldigen; wenn sie hören, daß ihr Ding nichts ist, wollen sie unsinnig werden; die sind das rechte Teufelsreich und seine Schuppen, dadurch er den größten Schaden thut. Die andern gehören unter den Henker und Fürsten, diese aber wollen auch mit dem Heiligen Geiste regiert sein; es ist aber nur Schein und Gleisen. Solches soll man sich gewöhnen zu verstehen. Der Vernunft und Menschen Dünkel steckt so tief in uns, daß es wenig Leute sehen und hingan gehen. Es ist schwer, daß man solche tapferen, großen Leute soll verachten; [es] kann’s niemand richten, denn der den Geist hat, es treugt Vernunft und aller Menschen Sinn.

\section{Das siebenzehnte Capitel}

\begin{schriftzitat}
V. 1—27. Als nun Abram neun und neunzig Jahr alt war, erschien ihm der HErr, und sprach zu ihm: Ich bin der allmächtige GOtt, wandele vor mir, und sei ohne Wandel; und ich will meinen Bund zwischen mir und dir machen, und will dich fast sehr mehren. Da fiel Abram auf sein Angesicht. Und GOtt redete weiter mit ihm, und sprach: Siehe, ich bin es, und habe meinen Bund mit dir, und du sollst ein Vater vieler Völker werden. Darum sollst du nicht mehr Abram\footnote{Randglosse: Abram heißt hoher Vater, Abraham aber, der Haufen Vater; wiewohl dieselben Haufen nur mit einem Buchstaben angezeigt werden in seinem Namen, nicht ohne Ursache.} heißen;

sondern Abraham soll dein Name sein; denn ich habe dich gemacht vieler Völker Vater. Und will dich fast sehr fruchtbar machen, und will von dir Völker machen, und sollen auch Könige von dir kommen.\footnote{Das Folgende, bis zu Ende des Capitels, fehlt in der Wittenberger.} Und ich will aufrichten meinen Bund zwischen mir und dir, und deinem Samen nach dir, bei ihren Nachkommen, daß es ein ewiger Bund sei; also, daß ich dein GOtt sei, und deines Samens nach dir. Und will dir und deinem Samen nach dir geben das Land, da du ein Fremdling innen bist, nämlich das ganze Land Canaan,
\quellseite{152}
zu ewiger Besitzung, und will ihr GOtt sein. Und GOtt sprach weiter zu Abraham: So halte nun meinen Bund, du und dein Same nach dir, bei ihren Nachkommen. Das ist aber mein Bund, den ihr halten sollt zwischen mir und dir, und deinem Samen nach dir, bei ihren Nachkommen: Alles, was männlich ist unter euch, soll beschnitten werden. Ihr sollt aber die Vorhaut an eurem Fleische beschneiden, dasselbe soll ein Zeichen sein des Bundes zwischen mir und euch. Ein jegliches Knäblein, wenn es acht Tage alt ist, sollt ihr beschneiden bei euren Nachkommen. Desselben gleichen auch alles, was Gesindes daheim geboren oder erkauft ist von allerlei Fremden, die nicht eures Samens sind. Also soll mein Bund an eurem Fleische sein, zum ewigen Bunde. Und wo ein Knäblein nicht wird beschnitten an der Vorhaut seines Fleisches, des Seele soll ausgerottet werden aus seinem Volke, darum, daß es meinen Bund unterlassen hat. Und GOtt sprach abermal zu Abraham: Du sollst dein Weib Sarai nicht mehr Sarai heißen, sondern Sara soll ihr Name sein. Denn ich will sie segnen, und von ihr will ich dir einen Sohn geben. Denn ich will sie segnen, und Völker sollen aus ihr werden, und Könige über viel Völker. Da fiel Abraham auf sein Angesicht und lachte, und sprach in seinem Herzen: Soll mir hundert Jahr alt ein Kind geboren werden, und Sara neunzig Jahr alt gebären? Und Abraham sprach zu GOtt: Ach, daß Ismael leben sollte vor dir! Da sprach GOtt: Ja, Sara, dein Weib, soll dir einen Sohn gebären, den sollst du Isaak heißen; denn mit ihm will ich meinen ewigen Bund aufrichten, und mit seinem Samen nach ihm. Dazu um Ismael habe ich dich auch erhöret. Siehe, ich habe ihn gesegnet, und will ihn fruchtbar machen, und mehren fast sehr. Zwölf Fürsten wird er zeugen, und will ihn zum großen Volke machen. Aber meinen Bund will ich aufrichten mit Isaak, den dir Sara gebären soll, um diese Zeit im andern Jahr. Und er hörte auf mit ihm zu reden, und GOtt fuhr auf von Abraham. Da nahm Abraham seinen Sohn Ismael, und alle Knechte, die daheim geboren, und alle, die er gekauft, und alles, was Mannes Name war in seinem Hause, und beschnitt die Vorhaut an ihrem Fleische eben desgleichen Tages, wie ihm GOtt gesagt hatte. Und Abraham war neun und neunzig Jahr alt, da er die Vorhaut an seinem Fleische beschnitt. Ismael aber, sein Sohn, war dreizehn Jahr alt, da seines Fleisches Vorhaut beschnitten ward. Eben auf Einen Tag wurden sie alle beschnitten, Abraham, sein Sohn Ismael, und was Mannes Name in seinem Hause war, daheim geboren und erkauft von Fremden; es ward alles mit ihm beschnitten.
\end{schriftzitat}

1. Wir haben bisher gehört, wie dem Patriarchen Abraham verheißen ist zum ersten GOttes Gnade und Schutz, darnach eine Frucht und Lohn von seinem rechten Weibe. Nun wird in dem Capitel beschrieben, wie GOtt abermal mit ihm redet, und seinen und seines Weibes Namen wandelt, und mit ihm aufrichtet den Bund der Beschneidung; und vornehmlich, daß er ihm den Sohn Isaak verheißt von seiner Frau Sara. Es ist ein groß Capitel, gehöret wohl den Gelehrten zu; doch wollen wir davon machen, so viel wir können. Die erste Verheißung ist, darinne ihm GOtt zusagt, wie [Cap. 15, 1.] gehört: er wolle sein Schutz und sein Lohn sein, an welchen Worten er mit dem Glauben gehangen hat, wie vorhin am andern. Aber, wie ich gesagt habe, wenn GOtt verneuet seine Verheißung und Wort, gilt es allemal eines neuen und bessern Glaubens. Denn weil wir auf Erden sind, wird keiner so vollkommen, muß immer zunehmen, und stärker werden. Also hat GOtt diesen heiligen Mann immer weiter geführt, und immer neue Stücke vorgehalten; so ist er immer geblieben in seinem Glauben, obwohl sich die Sache und Händel mancherlei begeben haben, um welcher willen er das Wort wieder gibt, denselbigen Glauben, wie gesagt, noch baß zu stärken.

2. Nun ist das ein großer mächtiger Glaube gewesen, daß ihn auch Paulus Röm. 4, 18—20. preiset: „Abraham hat GOtt gepreiset und die Ehre gegeben, als er glaubte wider alle Vernunft und Hoffnung, nicht angesehen seinen verstorbenen Leib, und seines Weibes“, da es nach der Natur unmöglich und nicht zu helfen war, daß sie sollten Kinder haben. Das hat er alles aus den Augen gesetzt, nicht darnach gerichtet, so, daß er fühlete und ihm unter Augen stieß, sondern nach dem bloßen Worte, und also gesagt: GOtt ist wahrhaftig, treu und allmächtig, darum frage ich nicht darnach, wie unmöglich es ist; da habe ich sein Wort und Verheißung, das kann mir nicht fehlen, wenn alles andere Ding trügt.


\quellseite{153}

3. Denn der Glaube ist ein Licht, das ihn führt, und leuchtet im Finsterniß, da die Vernunft blind ist und zum Narren wird. Der Natur ist es unmöglich, GOtt ist es aber möglich. Darum, wo das Wort ist: „GOtt sagt“, thue nur die Vernunft zu. In dem nun, daß du GOtt so preisest und ehrst, so wird dir der Glaube gerechnet zur Gerechtigkeit. Denn das ist, wie vor [Cap. 15, § 9] gesagt, die rechte Ehre GOttes, wenn man ihn dafür hält, daß er alles sei und vermöge, was GOtt zugehört; damit habe ich ihm gedienet und angebetet. Was gehört aber GOtt zu? Daß er allmächtig, gütig, gnädig, gerecht, wahrhaftig und weise sei, und alles, was gut ist. Solches alles ihm geben, ist der rechte Gottesdienst.

4. Das thut Abraham hier, wie Paulus [Röm. 4, 17. 21.] sagt, in dem, daß er glaubte dem, der da ruft die Dinge, die da nicht sind, daß sie seien. Denn er weiß wohl, Abraham, wenn alle Creaturen zusammen thäten mit allen Kräften, könnten sie ihm nicht einen Sohn geben; daran müssen alle Creaturen verzweifeln. Er denkt aber: Ja, GOtt kann es wohl thun, nicht allein von einem alten Weibe, sondern auch von einem Stein [Matth. 3, 9.]. [Es] ist aber etwas Großes,\footnote{In den Ausgaben: Ist aber groß.} daß er solches glauben kann, ob es wohl leicht scheint. Weil er nun das glaubte, so trauete er GOtt auch wohl mehr, was er sollte, daß er ihm das ewige Leben und alles Gut würde geben. Wenn er uns so vorschlüge, und spräche: Da hast du eine Hand voll oder ein Körnlein, davon will ich dich das ganze Jahr speisen; wer würde da glauben? Da würden wir sehen, wie unmöglich es der Natur wäre; denn weil du darnach sehen wolltest, so müßte der Glaube fallen.

5. So wird es gehen, wenn wir sterben sollen. Da ist GOttes Wort und das Sacrament oder Zeichen, daß dir GOtt will das ewige Leben geben. Da steht vor Augen das elende Wesen, Sünde, Tod und Hölle, und was man verläßt in diesem Leben. Dies ist gewiß, jenes ist ungewiß. Wer da könnte GOttes Wort fassen, das Petrus [2. Ep. 1, 19.] ein Licht in einem finstern Orte heißt, und die Augen zuthun, hindurch sehen durch den Tod in das Leben, wie hier Abraham thut: [der] tritt in das Finstere, springt nicht allein zurück, und erschrickt nicht, sondern ist getrost und fröhlich dazu, daß ihm das Herz lacht, daß es muß ein überschwänglicher Glaube gewesen sein, der auf das völlige gefaßt und gesagt hat: „Nun ist es gewiß. Als sollte er sagen: Wie groß seltsam Ding und Wunder ist das, wie unmöglich wird es vor den Leuten sein! Also, daß das Lachen nicht verstanden werde als aus einem Zweifel, sondern aus großem Wundern des Glaubens.

6. Nun, das Exempel ist auch uns geschrieben, daß wir allenthalben lernen die Art eines rechtschaffenen Glaubens, daß er kein Licht hat denn GOttes Wort, und daß der, so ihm folgt, rechten Gottesdienst thut. Darum Paulus [Röm. 4, 21.] sagt: „Er ward stark im Glauben, preiſete Gott, und wußte auf das allergewisseste, daß, was GOtt verheißt, das kann er auch thun“; als sollte er sagen: Hat es GOtt geredet, so wird er gewißlich nicht lügen. So siehst du nun, wie Abraham gar auf einen Haufen gibt GOtt, was er ihm geben soll, daß er allmächtig sei, und könne es thun; dazu weise, daß er wisse, wie er es soll thun; item, wahrhaftig, daß er gewiß halte, was er sagt etc. So hat er GOtt auf das höchste gedienet; denn wo das Herz solches nicht hält und glaubt von GOtt, ist kein rechter Gottesdienst, sondern alles Heuchelei, was man für Gottesdienst erdichtet. So ist der Glaube genugsam ausgestrichen, daß nicht noth ist, mehr davon zu sagen. [Wir] wollen nun sehen die Stücke, die das Capitel sonderlich handelt.

7. Am ersten nimmt und wandelt ihm GOtt den Namen, und spricht: „Du sollst nicht mehr Abram heißen, sondern Abraham“, daß der Name ein wenig verlängert wird; und gibt Ursach, warum er so soll heißen: „Ich habe dich gemacht vieler Völker Vater“; also, daß Abraham heiße ein Vater, davon mancherlei Volk komme. Der vorige Name [\hebraeisch{אַבְרָם}] Abram, von zwei Silben ist gemacht von dem Wort [\hebraeisch{אָב}] Ab, das heißt auf Deutsch ein Vater, und [\hebraeisch{רָם}] Ram heißt hoch. So heißt Abram ein hoher Vater. Den Namen wechselt er ihm, er soll nicht mehr hoher Vater heißen, sondern vieler Völker Vater; doch nicht also, daß dies Wort Abraham [\hebraeisch{אַבְרָהָם}] die Deutung völlig gäbe, die er hier setzt. Denn es ist nicht mehr denn ein Buchstabe, nämlich ein [\hebraeisch{ה}] h dazwischen gesetzt, daß der Name einer Silbe länger wird. Das muß nun ein mysterium sein. Der Buchstabe [\hebraeisch{ה}] h bedeutet in der
\quellseite{154}
hebräischen Sprache sonderlich vor den andern. Hier aber ist darinne beschlossen das Wort [\hebraeisch{הָמוֹן}] Hamon, das heißt einen Haufen; davon das Wort [\hebraeisch{מָהָמוֹן}] Mahamon, das wir sagen Mammon, kommet, welches Christus gern braucht im Evangelio; also, daß der einige Buchstabe deute einen Haufen Völker. Und hat so wunderlich den Namen gebrochen, daß schier nicht zu spüren wäre, daß es so viel sollte heißen, wo es GOtt nicht deutete, weil es so unter einander gebrochen und geworfen; will auch so bleiben. Ist ohne Zweifel darum gemacht, daß er soll sein ein Vater derer, die da rechtschaffen glauben sollten, wie St. Paulus Röm. 4, 11. anzeigt.

8. Also wandelt er dem Weibe auch den Namen. Zuvor hieß sie Sarai, jetzt, sagt er, soll sie Sara heißen. Da wird eines Buchstabens weniger. [\hebraeisch{שָׂרַי}] Sarai heißt meine Frau; denn der Buchstabe Jod am Ende gibt so viel, als das Wort mein: [\hebraeisch{אֵלִי אֵלִי}] Eli, Eli, mein Gott, [\hebraeisch{אָדוֹנַי}] Adonai, mein Herr, [\hebraeisch{אָבִי}] Abi, mein Vater. [\hebraeisch{שָׂרָה}] Sara aber heißt domina, eine Frau, nicht insgemein ein Weibsbild, sondern die da regieret im Hause; nicht wie der Mann, sondern wie eine Magd von einer sagt: Dies ist meine Frau. Da will GOtt mit eingezogen haben, daß Abrahams Frau, sein ehelich Weib, nicht seine Herrin, sondern unterthänig sein solle. Das hat er nicht darum gethan, daß viel daran gelegen wäre, sie hieße, wie sie hieße, ohne daß sie mußte sein eine Figur, Anfang und Deutung des neuen Testamentes Volkes. Denn so sagt Paulus Gal. 4, 24.: Hagar ist gewesen eine Mutter der Heiligen, die Ismael heißen, mit schönen Werken gefangen, die unter dem Gesetz, nicht lustig und willig, sondern Knechte sind, nicht des gesegneten Samens, sondern allein vom Fleisch und Blut, ohne das Wort Gottes.

9. Aber zu diesem Samen von der Frau kommt das Wort GOttes, das ist der Schatz, dagegen Fleisch und Blut nichts wäre. Es lautet wohl bei der Vernunft so schlecht, als hätte es ein Schuster geredet; vor GOttes Augen aber ist der Schatz so groß, daß [es] nicht auszudenken ist. Darum hat er sie nun Sara geheißen, daß sie die rechte Frau sei, das ist, daß wir Sara Kinder sind. Warum? Sara hat einen Sohn geboren, nicht von Fleisch und Blut, und ist doch Fleisch und Blut; [er] ist ihre Frucht, aber nicht Fleisches und Blutes Werks.\footnote{Jenaer: Werk.} Die göttliche Majestät gibt da mit dem Wort die Kraft aus, darum ist er ein Kind des göttlichen Wortes, doch im Fleisch und Blut erzogen. Das Wunder will GOtt ausrichten, daß er einen Sohn läßt geboren werden aus Fleisch und Blut, doch nicht durch Fleisch und Blut.

10. Darum sind\footnote{Erlanger: sind sie.} Gottes Kinder nicht, ohne die durch das Wort geboren werden. Jsaak ist auch leiblich durch das Wort geboren, [es] mußte gar nicht durch Fleisch und Blut geschehen; sonst wäre Hagar [und] Ismael auch Mutter und Sohn. So ist hier gar niedergelegt alles, was Fleisch und Blut vermag, und ein Urtheil geschlossen: Was nicht von GOttes Wort ist, das ist verloren, es sei ein geistlich Leben, wie es wolle; geht es im freien Willen, wenn es gleich nach GOttes Gesetz fährt, macht es nicht Kinder Gottes.\footnote{„Gottes“ fehlt in der Erlanger.} Sollst du aber ein Jsaak werden, mußt du anders geboren werden und auch die Verheißung haben, daß dich zum Erben mache. So schlecht liegt es da, wie ein todtes Wort, noch ist es so stark, daß es schließet und urtheilt über die ganze Welt. Das\footnote{Erlanger: Da.} heißt nun Paulus Sara Kinder, der Frau, die da Frau im Hause ist, des Gutes mächtig, nicht Magd, welche den Lohn nimmt und davon zeucht; sondern die man nicht kann aus dem Hause stoßen, und des Mannes Gut unter sich hat, regiert das Gesinde, und,\footnote{„und“ fehlt in der Erlanger.} was im Hause ist, darüber ist sie Herr. Darum heißt sie Paulus die Freie, das ist, niemandes Magd noch Unterthan, sondern das Haus ist ihr unterthan.

11. Das ist die christliche Kirche, die ist die rechte Sara, recht frei, hat niemand über sich, denn ihren Hausherrn, Christum, das ist ihr Mann, hat Recht zu ihr. So hat sie, was er hat, Wort, Geist, Taufe, Sacrament; was sie schafft, das ist gethan in Himmel und Erde; was sie sagt, das bestätiget er; was sie thut, das thut GOtt; [es] gilt eben so viel, als taufte GOtt selbst, oder selbst absolvierte. Die Gewalt ist überaus groß, daß sie niemand genugsam kann ermessen und preisen. GOtt hat ihr seine Gewalt in den Schooß gelegt, und als der Frau die Schlüssel gegeben, daß, was sie thut und
\quellseite{155}
schaft, soll gethan sein; nicht mehr, denn daß sie ihm gehorsam sei und nach ihm sich richte; sie darf alles thun, was sie will, allein, daß sie nicht über die Schnur trete und einem andern anhange, das ist, nicht weiter predige und mit Gewissen fahre, daß sie nichts Neues aufwerfe, sondern sich halte, das er thut, so ist sie die liebe Braut. Deß Kinder sind wir, werden von ihr geboren, so sind wir auch ihre Erben. Darum geht es hier fein zu, daß wir haben, was Christus und seine Braut hat. Das Evangelium aber ist ihr Bauch, damit trägt und gebiert sie; so werden wir GOttes Kinder, Isaaks Brüder, nach Isaak (das ist, nach der Verheißung) geboren von der Freien, die\footnote{Erlanger: wie.} mir Gottes Wort haben.

12. Darum sind wir nicht Knechte, spricht Paulus abermal [Gal. 4, 31.], sondern nur Kinder und Junker von der freien Frau; gehen aus freiem, willigem Geiste her, thun alles umsonst, nicht angesehen den Lohn, wie ein Knecht. Denn sie sind geboren in das Erbe, daß sie kein Werk, dasselbe zu erwerben, dürfen thun, sind der Mutter im Schooße; darum sehen sie ihre Werke nicht an, sondern thun mehr, denn man fordert, haben keinen andern Verdienst, denn daß sie geboren sind. Die Geburt bringt es mit sich, daß das Erbe ihr eigen ist; er ist der Vater, der muß denken, daß er das Erbe mache und sie in sein Gut setze; sie können nichts dazu thun.

13. Siehe, so sehen wir an allen Orten in der Schrift, was da sein sollte ein recht christlich Wesen, und wie GOtt solch Ding in der Schrift vorhin so versteckt und verborgen hat; doch so stark, daß jedermann muß sagen, es sei recht, und sehen kann, daß GOtt solches im Sinne habe gehabt, daß er einmal eine solche Predigt wollte lassen ausgehen vom Glauben. Denn hier siehst du gar kein Werk, kein Gebot, sondern lauter Verheißung, gnädig geschenkt.\footnote{Jenaer: „Geschenck“.} Er sieht kein Werk an, heißt ihn auch keines thun, sondern frei sagt er, sie soll ihm ein Kind gebären. Das sind Worte, nicht die gebieten, sondern die da schenken; darum wird er nicht geboren aus [des] Vaters oder [der] Mutter Verdienst, sondern kommt aus dem Worte, also, daß Abraham nichts dazu konnte thun. Wie er nun geboren wird, so müssen wir alle geboren werden; so kommen wir alle in das Erbe und werden Christen, haben es nie verdient noch gesucht, daß wir Gottes Wort hörten, sondern er kommt zuvor allen Gedanken, und spricht: Glaube an Christum, so wirst du selig. Davon werden wir andere Menschen, unverdient und unversehens.

14. Was will sich nun Isaak rühmen? Woher wird er ein Erbe? Abraham hat es wohl gesammelt, aber ihm nicht gemacht noch gedacht; es kommt allein daher, daß ihn GOttes Wort macht zu einem Kinde, und legt ihn in Abrahams Schooß. Darum kann er sich nichts annehmen, denn daß ihn GOtt hat in das Erbe gesetzt, nicht aus seinem Zuthun, sondern durch das lautere Wort. So darf er sich nicht rühmen, daß er GOttes Kind sei, denn dadurch, daß ihm GOtt hat das Evangelium predigen lassen, und den Glauben gegeben. So sind wir auch selig nirgend her, denn daß uns GOtt durch das Wort geboren hat, und zu Kindern gemacht.

15. Siehe, das sind die schlechten Historien, wenn man es nach der Vernunft achtet, daß man möchte GOttes spotten; wie auch etliche gethan haben. Aber laß sie narren und scherzen; GOtt hat die Art, daß er mit menschlichen Dingen umgeht, als, mit Kinder zeugen und dergleichen; ist aber mit solchen gnadenreichen Sprüchen und hohen Deutungen geschmückt und geziert darum allein, daß er die hohen Klüglinge zu Narren mache.

16. So schließen wir nun aus der Historie, was vor GOtt angenehm sei, damit er Kinder macht; keine Menschenkräfte, kein Gebot noch Werk. Er muß zuvor kommen mit seinem Worte, daraus wir Kinder werden mit\footnote{Erlanger: wie.} Isaak, daß wir wissen, daß lauter Gnade ist, und allein um den Glauben zu thun, der so groß ist, daß er rechte Erben macht, und allein der rechte Gottesdienst ist.

17. Weil nun dem also ist, was thut man denn nach dem Erbe oder neuen\footnote{„neuen“ fehlt in der Erlanger, steht jedoch in der Wittenberger im Text, und in der Jenaer in Klammern.} Geburt? Antwort: Darum muß man ja die Werke auf andere Dinge richten, und sagen: Du mußt gleichwohl nicht müßig gehen, sondern Werke thun, aber damit schaffst du nicht das Erbe, weil es vorhin da ist. Wie, soll man denn nichts thun? Nein, sondern hingehen und schaffen, daß das Erbe und Gut gemehrt und groß werde. So haben wir es gar durch den Glauben. Aber gute Werke thun wir nicht, denn darum, daß
\quellseite{156}
das Reich größer werde, predigen und bringen andere Leute dazu mit Worten und Werken, daß unser Leben darnach gerichtet sei auf andere Leute, ihnen auch zu helfen. Diese Lehre ist uns hier in dem Stücke gegeben, welches wir nun so oft gesehen und gehört haben, daß wir’s gewiß seien und gerüstet, und lassen uns nicht davon reißen.

18. So haben wir die mächtige Zusagung, und den großen Glauben Abrahams, und wie sein und seines Weibes Name verändert sind: seiner erlanget, und ihrer verkürzt, und kehret es\footnote{„es“ fehlt in der Erlanger.} Gott also um. Wenn das Wort kommt, so geht des Glaubens Reich an, und hören unsere Werke auf, und gilt nichts mehr, was an uns ist. Der Glaube nimmt zu, und die Werke ab; sonst sind sie groß, lang und viel, aber der Glaube klein und kurz. Darum muß sich’s also umkehren, daß hier etwas zugehe, daß Abrahams Name länger werde, und dort etwas abfalle.

19. Nun ist noch ein Stück hier, von der Beschneidung, welches auch ein recht närrisch und schimpflich Gebot anzusehen ist. Davon habe ich mehrmal gepredigt, wer da will, mag es weiter lesen in den Postillen. Es ist das Stück, daß St. Paulus allenthalben anzeucht, und hat wunder damit zu schaffen. Denn die Juden verließen sich hart darauf.

20. Siehe aber auf\footnote{Erlanger: darauf, auf.} die Ordnung: Abraham ist verheißen, daß GOtt wollte sein Schutz sein; item, ihm einen Sohn geben von seiner rechten Frau, und seinen Samen sehr mehren etc.; nach denselbigen Zusagungen allen beschneidet er ihn, zu führen dies Zeichen an seinem Fleische, damit er äußerlich bekenne, daß er sein GOtt, und er, Abraham, sein Knecht sei. Da muß er ein äußerlich Werk thun; noch ist es nicht seine Gerechtigkeit, sondern muß es frei umsonst thun. Denn der Glaube hat ihn schon zuvor gerecht und selig gemacht, gleichwie Isaac wird geboren durch Gottes Gnade; so bald er aber geboren ist, ist er schon der Erbe. Darum hilft ja die Beschneidung nichts dazu; wenn er schon nimmer beschnitten wäre, so wäre er dennoch der Sohn geblieben.

21. Das legt nun Paulus meisterlich aus, als er Röm. 4, 11. sagt: die Beschneidung sei nicht mehr, denn ein Zeichen, damit er vor den Leuten sich bekennete. Denn zuvor, weil er durch den Glauben rechtfertig war, mußte er das Zeichen tragen, zu bekennen den Gott, an den er glaubte. Darum heißt es Paulus sigillum justitiae fidei, ein Siegel der Gerechtigkeit des Glaubens. Denn gleichwie mit einem Siegel bewiesen wird, was Kraft haben soll, so wird auch durch die Beschneidung beweiset vor der Welt, daß jedermann sagte: Damit bekennt Abraham seinen GOtt und Glauben. Wie auch wir, wenn wir glauben, sind wir von allen Sünden frei und GOttes Kinder; das beweisen wir vor den Leuten mit der Taufe, und mit dem närrischen Zeichen, wie die Heiden sagen, daß wir unsern GOtt essen. Denn es sind alles solche Zeichen, derer\footnote{Wittenberger und Erlanger: bes.} die Vernunft muß lachen und spotten.

22. Aber GOtt nimmt darum das Zeichen, daß er die Deutung will darinne zeigen. Er hätte wohl Abraham an einem andern Orte mögen heißen beschneiden, denn an dem schändlichen Orte, daß [es] vor der Welt Sünde und Schande ist; darum auch die Heiden der Juden auf das schmählichste spotteten und sie anspeieten. Es dient ja gar nichts, Gnade dadurch zu verdienen, sondern allein zum Zeichen des Glaubens, gefordert zum Bekenntniß und Zeugniß vor der Welt, daß sie sich daran stoße und zum Narren werde, und dennoch niemand selig werde, er sei denn des Glaubens des beschnittenen Abrahams [Röm. 4, 16.].

23. Darum mußt du deine Klugheit herunter werfen und ein Narr werden. Lache, wie lange du kannst, daß man uns in das Wasser steckt; dennoch ist es beschlossen vor GOtt: Sollst du der Sünden los und selig werden, so mußt du werden wie der Narr, der aus der Taufe kommt, oder wirst des Himmelreichs fehlen.

24. Also nimmt GOtt gerne solche närrischen Zeichen, vor der Welt schändlich und schmälich, daß sie sich daran ärgere.\footnote{Wittenberger und Erlanger: ärgerte.} Denn wo er es nicht hätte, so bliebe der stolze Muth in uns stecken, daß wir und unsere Weisheit auch etwas wären; da würde denn jedermann zulaufen. Aber er will, daß wir den Sinn brechen, und in eine andere Haut kriechen. Und weil GOtt solches will, sollen wir es uns wohlgefallen lassen, und nicht für Narrheit halten. Das ist die Ursache
\quellseite{157}
, warum es leiblich eingesetzt ist; was es aber geistlich bedeutet, habe ich sonst gesagt. [Es] ist kürzlich nichts anders, denn ausziehen den alten Adam, durch das Kreuz, Schmach und Schande in Tod gehen und alles Unglück, wie [es] uns Gott zuschickt.

25. Aber daneben ist zu merken, daß wir können mit den Juden streiten und den Text recht führen, nämlich, daß wir die Beschneidung nicht höher halten, denn ein äußerlich Zeichen, damit Abraham beweist seinen inwendigen Glauben, wie wir in der Taufe. Denn so führt Paulus die Beschneidung wider die Juden, und beschließt endlich, daß sie sich nicht können rühmen des Werkes vor allen Heiden, davon sich Abraham selbst nicht rühmen kann, daß er dadurch vor GOtt etwas geworden sei, sondern war es bereits zuvor, durch den Glauben an GOttes Wort. Also auch wir, ehe wir das Sacrament der Taufe oder des Altars nehmen, müssen wir den Glauben haben, das Zeichen aber ist eine Beweisung vor den Leuten.

26. So sollst du auch sagen von allen guten Werken. Warum sind sie geboten, so doch der Glaube genugsam ist? Wahr ist es, der Glaube ist allein genugsam\footnote{Erlanger: genug.} vor Gott, daß man keines Werkes dazu bedarf; doch muß man die Werke darum thun, daß man den Glauben beweise vor der Welt, damit GOtt gepreiset werde und mein Nächster auch zum Glauben gereizt werde [Matth. 5, 16.]. Ich dürfte es für mich gar nichts, sondern muß um anderer willen lieben, GOtt zu Lobe und dem Nächsten zu Dienst, auf daß andere Leute auch zu uns kommen. Das nun von Abraham geboren, allein vom Fleische und Blute, und auch beschnitten ist, das thut nichts; denn Ismael ist auch von Abraham und auch beschnitten, noch wird er nicht gerechnet für Abrahams Samen. Woran fehlt es denn? Daran: Er hatte wohl Fleisch und Blut, daß er brachte von Abraham, hatte aber Glauben und Wort nicht. Denn wer sein Same sein soll, sagt Paulus, der muß den Glauben haben, den er hat. So wird er ein Vater aller derer, so an Christum glauben, da er an geglaubt hat. Davon wir nun genug haben gehört.

\section{Das achtzehnte Capitel}

\begin{schriftzitat}
V. 1—20. Und der HErr erschien ihm im Hain Mamre, da er saß an der Thüre seiner Hütte, da der Tag am heißesten war. Und als er seine Augen aufhub und sahe, da standen drei Männer gegen ihm. Und da er sie sahe, lief er ihnen entgegen von der Thüre seiner Hütte, und bückte sich vor ihm nieder\footnote{Randglosse: Vor ihm nieder. Vor Einem fällt er nieder, und redet auch als mit Einem, und doch mit Dreien; da ist die Dreifaltigkeit in Gott angezeigt.} auf die Erde, und sprach: HErr, habe ich Gnade funden vor deinen Augen, so gehe nicht vor deinem Knechte über. Man soll euch ein wenig Wassers bringen, und eure Füße waschen, und lehnet euch unter den Baum. Und ich will euch einen Bissen Brods bringen, daß ihr euer Herz labet, darnach sollt ihr fortgehen. Denn darum seid ihr zu eurem Knechte gekommen. Sie sprachen: Thue, wie du gesagt hast. Abraham eilte in die Hütte zu Sara, und sprach: Eile, und menge drei Maß Semmelmehl, knete und backe Kuchen. Er aber lief zu den Rindern, und holte ein zart, gut Kalb, und gab es dem Knaben, der eilte und bereitete es zu. Und er trug auf Butter und Milch, und von dem Kalbe, das er zubereitet hatte, und setzte es ihnen vor, und trat vor sie unter den Baum, und sie aßen.\footnote{Der folgende Text fehlt in der Wittenberger.} Da sprachen sie zu ihm: Wo ist dein Weib Sara? Er antwortete: Drinnen in der Hütte. Da sprach er: Ich will wieder zu dir kommen, so ich lebe, siehe,\footnote{Erlanger: „Uber ein Jahr, nach der Zeit die Frucht leben kann, will ich wieder zu dir kommen, so soll“ etc. — Außerdem sind in diesem Text noch einige Varianten; wir haben, ebenso wie die Jenaer, den Text der Bibel gegeben.} so soll Sara, dein Weib, einen Sohn haben. Das hörte Sara hinter ihm, hinter der Thüre der Hütte. Und sie waren beide, Abraham und Sara, alt und wohl betaget, also, daß es Sara nicht mehr ging nach
\quellseite{158}
der Weiber Weise. Darum lachte sie bei sich selbst und sprach: Nun ich alt bin, soll ich noch Wohllust pflegen, und mein Herr auch alt ist. Da sprach der HErr zu Abraham: Warum lacht denn Sara, und spricht: Meinst du, daß wahr sei, daß ich noch gebären werde, so ich doch alt bin? Sollte dem HErrn etwas unmöglich sein? Um diese Zeit will ich wieder zu dir kommen, so ich lebe, so soll Sara einen Sohn haben. Da leugnete Sara, und sprach: Ich habe nicht gelacht; denn sie fürchtete sich. Aber er sprach: Es ist nicht also, du hast gelacht. Da stunden die Männer auf von dannen, und wandten sich gegen Sodom, und Abraham ging mit ihnen, daß er sie geleitete. Da sprach der HErr: Wie kann ich Abraham verbergen, was ich thue? Sintemal er ein groß und mächtiges Volk soll werden, und alle Völker auf Erden in ihm gesegnet werden sollen. Denn ich weiß, er wird befehlen seinen Kindern und seinem Hause nach ihm, daß sie des HErrn Wege halten, und thun, was recht und gut ist, auf daß der HErr auf Abraham kommen lasse, was er ihm verheißen hat.
\end{schriftzitat}

1. Dies Capitel zeigt an, wie der heilige Vater Abraham vor GOtt ein großer Mann und hoch angesehen ist gewesen; daß der allmächtige GOtt anzeigen will aller Welt mit dem Exempel, wie gnädig und hold er ist denen, die da glauben. Es ist je wahr, daß nicht größer Ding auf Erden ist, ja, in allen GOttes Werken, denn der Glaube; der ist so mächtig, daß er alles hinwegnimmt, was sündlich ist, macht uns Kinder GOttes, daß er also mit uns umgeht, wie ein Bruder mit dem andern. Also hier, mit Abraham; wie freundlich thut er sich doch zu ihm? [Er] sieht nichts an, denn den hohen Glauben; damit hat er GOtt gefangen, gefällt ihm so wohl, daß er mit ihm scherzt wie mit einem besten Freunde, ja, wie mit einem Kinde: kommt selb dritte, ist mit ihm, scherzt mit Sara; doch, daß Ernst dabei ist. Zuletzt hebt er selbst an und spricht: „Wie kann ich Abraham verbergen, was ich thue?“ Als sollte er sagen: Er ist mein Freund, so gut, daß ich ihm alles muß sagen. Am Ende redet Abraham mit ihm, und er hört ihm so lange zu, und gewährt ihm, so oft er bittet.

2. Siehe, da ist erklärt und ausgelegt, was Paulus [Tit. 3, 4.] meint mit den Worten: Apparuit χρηστότης et φιλανθρωπία Dei, daß ist, die Freundlichkeit und Leutseligkeit GOttes, daß er so leutselig ist, daß jedermann gerne mit ihm umgeht und um ihn ist, und er wieder mit jedermann freundlich fährt. Die Tugend gibt er GOtt eigentlich im neuen Testamente, durch das Evangelium offenbart, hat es auch reichlich bewiesen; wie du hier siehst, daß ein Mensch mit dem andern nicht könnte freundlicher umgehen. Das ist nun der Preis des Glaubens Abrahams, uns zu einem Exempel, daß auch wir, wenn wir glauben, uns tröstlich zu ihm versehen, daß er auf das allerfreundlichste mit uns fährt; daß wir es aber nicht erfahren noch schmecken, das macht, daß wir nicht glauben. Wenn wir es fühlten, würden wir sehen, wie lieblich Historien das sind, die den Glauben also vorbilden, daß Abraham GOttes Freund heißt. Summa, es ist ein hohes Wunder und nicht zu ermessen, über alle Worte und Gedanken, daß Gemüssen muß es selbst erfahren. Wir achten es gering, wenn er redet; wenn wir aber dächten, was [es] für eine Person wäre, würden wir davor niederfallen. Weil wir denken, es sei gerade, als rede ein Schuster, erfahren wir nimmer nichts davon.

3. So laß nun das groß und hoch sein, daß GOtt selbst sich so zu ihm thut, geht in sein Haus und ist sein Gast, redet so viel Worte mit ihm; darnach auch, wiederholt die Verheißung, die er ihm im vorigen Capitel [V. 16.] gethan hat, daß seine Sara sollte einen Sohn gebären, welches sie nicht gehört hatte; jetzt aber wird es so stark getrieben, daß sie es auch hören muß, wiewohl sie nicht so bald glaubt, sondern noch zweifelt. Was nun mehr hierin angezeigt ist, weiß ich nicht; es sieht mich also an, als habe Sara dem Abraham nicht geglaubt, denn er wird es ihr ohne Zweifel oft gesagt haben. Jetzt, als sie es von den Dreien hört, däucht sie es so wunderlich, daß sie gleichsam spöttlich dazu redet, und spricht: Soll ich noch Wohllust treiben, wie eine junge Frau? Schämt sich noch, daß sie mit Wohllust soll umgehen. Da zeigt er auch an, was für ein fein, züchtig Weib sie ist gewesen, auch im äußerlichen Wandel; auch wie sie ihrem Manne unterthan ist, und heißt ihn Herr. Welches St. Petrus [1. Ep. 3, 6.] für gut angesehen hat, daß er es anzeige, und sie den Weibern zum Exempel setzte.

4. Nun, wir müssen auf dem bleiben, daß Sara selbst muß das Wort hören. Sie will
\quellseite{159}
nicht mehr Lust in der Welt wissen, steht in dem Sinne, daß sie nun des Todes wartet, und ihr keine Lust zu Muthe war, wie der Text sagt: Non erant ei muliebria, es ging ihr nicht mehr nach der Weiber Art. [Es] ist auch nach der Natur nicht da gewesen, daß sie ein Kind sollte haben; darum schlägt sie es in Wind, glaubt Abraham noch nicht, denkt also: Wer weiß, ob es mein Mann recht verstanden habe? Derhalben kommt nun GOtt selbst, redet mit ihr so freundlich. Welches er von einem andern nicht würde leiden: Sie hört es von Abraham, und zweifelt; jetzt, als sie es vom Engel selbst hört, hebt sie an zu lachen. Das würde ihm von einem andern nicht gefallen haben.

5. Weiter fährt er zu, und scherzt nicht, spricht: „Warum lachet Sara, dein Weib? Sollte es GOtt zu schwer sein?“ Dennoch soll es geschehen: „Wenn ich wieder komme, so soll Sara einen Sohn haben.“ Das sind eitel überflüssige Worte, daß er redet, wie mit seiner Tochter; welches er gegen andere nicht thäte. Das thut es, wo der Glaube rechtschaffen ist, da ist GOtt mit allen seinen Gütern und Segen, ist alles schlecht eitel Friede und Ruhe, ob gleich der Glaube noch schwach ist.

6. Denn hier haben wir ein groß, tröstlich Exempel gesetzt deß, so Paulus Röm. 14, 1. sagt: „Den Schwachen im Glauben nehmet auf.“ Es ist je hier ein rechtschaffener Glaube gewesen, weil sie zuvor mit Abraham ausgezogen, in starkem Glauben, in Egypten, in des Königs Hause, wie gehört, daß sie nimmer hätte gethan, wäre sie nicht in festem Glauben gestanden. Item, so stark war sie, daß sie Abraham ihre Magd gab zu einem Weibe, das sie auch nicht gethan hätte, wo sie nicht gewiß wäre gewesen, daß sie von Abraham Kinder würde haben; noch strauchelt und zweifelt sie hier. Das heißt ein schwacher Glaube; dennoch verwirft sie GOtt nicht darum, schont ihrer um des vorigen Glaubens willen, hält es ihr zugute, ob sie nicht so stark glaubt als Abraham; denn er geht mit ihr um, wie Paulus [Röm. 8, 26.] sagt, daß er ihrer Schwachheit aufhilft, und redet so viel mit ihr, daß sie muß glauben. Als sie leugnet: Nein, ich habe nicht gelacht; spricht er: Du hast gelacht. Da sahe sie, daß es GOttes Wort war, wiewohl sie die Person nicht kannte.

7. So ist nun die Summa über den Preis des Glaubens, der hier angezeigt wird; ist auch das Exempel, daß GOtt die Seinen immer führt von einem Glauben in einen andern, schlägt immer ein neu Stück vor zu Uebung, Besserung und Zunehmen des Glaubens von Tag zu Tag. Darum leben wir auch auf Erden, sonst wäre es besser, daß wir bald stürben. Dies sind die Stücke des Glaubens in dem Exempel.

8. Darnach sind auch hierinne moralia, das ist, Exempel der Liebe. Denn dies Exempel führt die Epistel zum Hebräern [Cap. 13, 2.]: „Vergesset nicht gastfrei zu sein, denn durch dasel[b]ige haben etliche, ohne ihr Wissen, Engel beherberget.“ Da meint er den Abraham und seinen Vetter Lot, im folgenden Capitel. Mit dem Werke hat GOtt uns Exempel gegeben und gepreiset die Mildigkeit, daß wir gerne herbergen sollen und den Fremdlingen Gutes thun. Zu welcher Tugend oft die Schrift, sonderlich das Neue Testament, vermahnet, und ist solch Exempel nicht umsonst geschrieben. Abraham war zuvor im Glauben, den übt er nun, und kommt fort in die Liebe auch; hat seine Thüre lassen offen stehen, wer gekommen ist, [den] hat er aufgenommen; ja, als er hier die drei Männer von ferne siehet, geht er ihnen entgegen, und nöthigt sie, daß sie bei ihm bleiben. Er ist da gesessen, sagt der Text, vor der Thüre, und hat geruht, da es um Eilens Zeit gewesen ist (denn sonst hat er freilich nicht viel müßig gegangen), und gewartet, wo Gäste kämen, die er aufnähme. Der\footnote{So die Wittenberger. In den anderen Ausgaben: „Das“.} Exempel führt man jetzt unter uns keines mehr, daß man arme Pilgrim oder Fremdlinge aufnähme, oder, wenn man es schon thut, gleichet es um Genieß willen. Wer es aber aus Liebe thäte, würde ohne Zweifel auch Heilige zu Gast kriegen, ja, GOtt selbst haben.

\begin{schriftzitat}
V. 20—33. Und der HErr sprach: Es ist ein Geschrei zu Sodom und Gomorra, das ist groß, und Sünden, die sind fast schwer. Darum will ich hinab fahren, und sehen, ob sie alles gethan haben nach dem Geschrei, das vor mich kommen ist, oder ob es nicht also sei, daß ich es wisse. Und die Männer wandten ihr Angesicht, und gingen gen Sodom. Aber Abraham blieb stehen vor dem HErrn, und trat zu ihm, und sprach: Willst du denn den Gerechten mit dem Gottlosen umbringen? Es möchten vielleicht fünfzig Gerechte in
\quellseite{160}
der Stadt sein, wolltest du die umbringen, und dem Orte nicht vergeben um fünfzig Gerechter willen, die darinnen wären? Das sei ferne von dir, daß du das thust, und tödtest den Gerechten mit dem Gottlosen, daß der Gerechte sei gleich wie der Gottlose; das sei ferne von dir, der du aller Welt Richter bist, du wirst solch Gericht nicht thun. Der HErr sprach: Finde ich fünfzig Gerechte zu Sodom in der Stadt, so will ich um ihrer willen alle den Orten vergeben. Abraham antwortete, und sprach: Ach siehe, ich habe mich unterwunden, zu reden mit dem HErrn, wiewohl ich Staub und Asche bin. Es möchten vielleicht fünf weniger denn fünfzig Gerechte darinnen sein, wolltest du denn die ganze Stadt verderben um der fünf willen? Er sprach: Fünde ich darinnen fünf und vierzig, so will ich sie nicht verderben. Und er fuhr weiter mit ihm zu reden, und sprach: Man möchte vielleicht vierzig darinnen finden. Er aber sprach: Ich will ihnen nichts thun, um vierziger willen. Abraham sprach: Zürne nicht, mein HErr, daß ich noch mehr rede, man möchte vielleicht dreißig darinnen finden. Er aber sprach: Fünde ich dreißig darinnen, so will ich ihnen nichts thun. Und er sprach: Ach siehe, ich habe mich unterwunden, mit meinem HErrn zu reden, man möchte vielleicht zwanzig darinnen finden. Er antwortete: Ich will sie nicht verderben um der zwanzig willen. Und er sprach: Ach zürne nicht, mein HErr, daß ich nur noch einmal rede, man möchte vielleicht zehn darinnen finden. Er aber sprach: Ich will sie nicht verderben um der zehn willen. Und der HErr ging hin, da er mit Abraham ausgeredet hatte, und Abraham kehrte wieder hin an seinen Ort.
\end{schriftzitat}

9. Siehe, da ist noch ein trefflich Exempel, darinnen man sieht den großen Ernst von Abraham, wie er voll Liebe gegen dem Nächsten gewesen ist, als er hört, daß GOtt mit Ernst daran will, und die Städte Sodom und Gomorra verderben. Da sorget er, daß GOtt seinen Ernst lasse gehen, und das Land strafe, wie zuvor [1 Mos. 7, 17.] die Welt mit der Sündfluth. Darum tritt er vor ihn, übt ein edel Werk der Liebe, daß er auf sich nimmt alle Sünden, die sie gethan haben, und dafür mit allem Ernste bittet, nicht einmal, sondern sechsmal. [Es] sind nichts anders denn eitel rechte feurige Worte, und ein so herzlich, heftig Gebet, als ich keines in der Schrift weiß. So steht sein Muth, daß er je gerne wollte den Zorn zurück treiben, hält so fest an, und bringt es GOtt so nahe, daß er muß von ihm weg gehen. Das ist je ein recht ausbündig Liebe-Werk. So sollten wir auch\footnote{„auch“ fehlt in der Erlanger.} thun und bitten, wo wir sehen jemand in Sünden liegen, da gewiß ist, daß GOttes Zorn und die Strafe vorhanden ist, auf daß wir erwehreten, daß er die Strafe nicht gehen ließe.

10. So nimmt er nun zum ersten vor sich GOttes Güte, aufs härteste; sagt nicht, daß sie nicht genügend haben, sondern ruft seinen Namen an, rückt ihm seine Ehre auf, will also sagen: Thust du das, so wirst du deinem Namen eine Unehre und Schmach aufthun; schone doch dein selbst. Was würde die Welt sagen, denn daß du ein ungnädiger, greulicher GOtt wärest? So greift er ihn auf das allerhöchste an, und gewinnt ihn damit ab, hält ihn so lange auf, bis er nimmer kann. Das macht, daß er also dachte: Zu einer solchen Gegend werden je viel frommer Leute sein; [er] meinte, er wollte es ihm enge spannen, und nur fünfzig zählen; ja, da bittet er das erste, andere, dritte, bis auf das sechste Mal; geht alles hin. Darüber thut er nicht weiter bitten, dachte also: Ist ihrer nur fünf, so wird er sie wohl behalten, wenn schon die andern untergehen.

11. Da merke nun: Man sagt viel, es seien noch viel fromme Leute in der Welt; ja, tröste dich darauf, Abraham hat sich auch darauf vertröstet, daß er in der Zuversicht das Gebet thäte. Wär es ihm gesagt hätte, daß in demselben Lande nicht zehn fromme Leute wären, hätte er es nicht geglaubt; noch findet er nicht mehr denn Lot, selb dritte, die andern verderben allzumal, ohne den Fremdling und Gast; wie oben zuvor [Cap. 7, 13.] in der Sündfluth nicht mehr denn acht Menschen in der großen, weiten Welt erhalten sind. Rühme sich nun, wer da will, daß man viel will schließen: Es sind noch viel frommer Leute auf Erden. Man verläßt sich darauf, und schließt darnach weiter: Sie können nicht irren, [es] sei nicht zu glauben, daß GOtt so zornig sei. So hat uns GOtt den Trotz weggerissen, daß wir darauf nicht bauen. Daß [es] nicht gilt, daß unsere Papisten vorgeben: So viel Leute können nicht irren. Wer dürfte jetzt solche Städte Sodom gleichen? wie Christus die Stadt
\quellseite{161}
Capernaum dagegen hält, und sagt [Matth. 11, 24.]: „Es soll Sodom erträglicher ergehen am jüngsten Gerichte, denn ihr.“

12. Also halte ich auch für gewiß, daß der Zorn GOttes, der da gewesen ist, nirgend gleich ist dem Zorn, der bisher gegangen ist unter dem Pabstthum, und noch geht, daß ich wohl mag sagen, daß so viel Bücher geschrieben, und Doctores der Schrift gemacht sind, und kein Buch noch Doctor das Evangelium gepredigt hat. Wer es auch sich unterstanden hat, mußte bald verbrennen. Dieweil haben wir immer gesagt: Ei, GOtt ist nicht so zornig, daß er so viel Leute lasse irren. Sind etliche fromm, so haben sie ja den Schein nicht noch das Ansehen, daß sichs\footnote{sichs = sich des. Wittenberger und Jenaer: sich.} niemand zu ihnen versieht.

13. Darum soll uns das schrecken, und nicht wundern, ob GOtt ein ganz Land verderbt, weil wir wissen, daß durch GOtt verkündigt ist, daß die Verkündigung so groß sein würde, daß Christus muß sagen [Luc. 18, 8.]: „Meinest du auch, wenn des Menschen Sohn kommen wird, daß er Glauben finden wird auf Erden?“ Ich habe Sorge, es sei allzumal und allzustark erfüllt, daß alles, was nur unter und in dem Pabstthum ist, des Teufels sei. Es ist wohl schrecklich, noch ist es GOttes Gericht; ist doch alles verkehrt, was je unter das Pabstthum kommen ist. Es ist nicht ein gemein Ding um einen frommen Mann; Abraham ist freilich das edelste Kleinod dieselbe Zeit auf Erden gewesen.

14. So ist nun dies geschrieben, eben zu schrecken alle, die in Verzweiflung gehen; nicht um derer willen, die blöde und furchtsam sind. GOtt will niemand freundlich noch hold sein, denn die sich demüthigen, an ihnen verzweifeln, sich bücken und Gnade suchen bei ihm; denen ist alles freundlich, den andern aber ist alles schrecklich.

15. Dies ist das Exempel von Abraham, dabei wir sehen, wie voll des Geistes er gewesen ist, und große Stärke des Glaubens und hitzige Liebe, das uns vorgeschrieben ist, ihm auch nachzufolgen, wie er sich so gar der armen Leute annimmt, als wäre er selbst in der Noth, wie gehört, daß er über die Maßen treffliche Worte gebraucht, ob er sie möchte erretten und den Zorn abwenden, und auch allenthalben sein Gebet erhört wird, wie oft er bittet.

16. Das ist nun nach der Historie hin gesagt. Was ander tief und subtil Ding ist, laß ich anstehen. Sonderlich vorne im Capitel, da Moses also redet, daß er Einen habe angebetet, und doch Drei gesehen. Denn die Worte lauten also: „Der HErr erschien ihm im Hain Mamre, und da er seine Augen aufhub, sahe er drei Männer vor ihm stehen, da lief er hin, und bückte sich vor ihm nieder, und sprach: HErr, habe ich Gnade funden vor dir, so gehe nicht vor deinem Knechte über, man soll euch ein wenig Wassers bringen“ etc. Daher sagen etliche, wie jetzt gesagt, er habe drei Männer gesehen, und mit Einem geredet; damit sei die Dreifaltigkeit in GOtt angezeigt. Das lasse ich die Gelehrten ausfechten. Das ist je wahr, daß Abraham da hat Einen angesehen an GOttes Statt, welcher es auch gewesen ist. Das gibt der Text klarlich, als er spricht: „HErr, habe ich Gnade funden vor deinen Augen“, welches aus der hebräischen Sprache bewiesen wird, daß es der rechte Name GOttes ist.

17. Ob er nun das gethan habe, daß er die Person in GOttes Namen angenommen habe, oder selbst so erleuchtet gewesen ist, daß er\footnote{„er“ fehlt in der Erlanger.} wahrer GOtt wäre, befehle ich abermal den Gelehrten. Es sieht ihm fast gleich, daß er hernach so andächtig mit GOtt selbst redet, als habe er die eine Person für den rechten GOtt angesehen. Denn wir haben gehört, daß er mit dem HErrn selbst redet, und er ihm antwortet, welches doch geschehen ist in einer menschlichen Person. Darum muß er einen hohen Geist gehabt haben, der es ihm gesagt hat, daß er unter den Dreien Einen sieht, der GOtt ist, und die andern Engel, und doch alle Drei annimmt wie Menschen, als sie auch geberden.

18. Summa, so viel ist es, daß Abraham hier ein Bild und Blick gesehen hat, daß GOtt müßte sein, und sollte Mensch werden; denn er sieht das Menschenbild vor sich, hält ihn doch für wahren GOtt, betet ihn auch an wie einen GOtt. Es ist eine große Gnade und merklich Exempel, wie der Mann so angenehm gegen GOtt ist, daß er ihm so erscheint, und solch sein Geheimniß zu erkennen gibt; also, daß es darum geschrieben ist, daß wir je lernen sollen, wie groß Gut es ist um einen gläubigen Menschen. Denn, was darf es Abraham, daß man es schreibe? Es ist uns geschrieben, daß wir sehen, wie GOtt
\quellseite{162}
Lust und Liebe hat an denen, die ihm glauben, daß er sich so nahe zu ihnen macht und umgeht, wie ein Geselle mit dem andern, daß wir uns auch so zu ihm versehen, und mit ihm reden; doch, daß solches geschehe durch den HErrn Christum. Wer den Glauben hat, der mag frei mit GOtt reden, so redet er wieder mit ihm wie ein Bruder, oder wie ein Vater mit einem Kinde. Das hat nun\footnote{„nun“ fehlt in der Erlanger.} die Schrift hier vorgebildet, zu reizen zu dem Glauben durch Christum. Denn Abraham mußte GOtt nicht bloß sehen, sondern im Menschenbilde: so müssen wir ihn auch sehen in dem Bilde, da er sich eingestellt hat, das ist Christus. Kommt der nicht vor, so ist es verloren; wie die Werktheiligen thun, die ohne Christum mit ihrem Thun vor GOtt wollen treten.

19. Zuletzt ist hier angezeigt die Zuversicht, so wir haben sollen zu bitten, daß, wer ein Christ ist und glaubt, mag bitten, was er will, [so] ist er sicher, daß er wird erhört. Siehe, wie Abraham so gewaltiglich redet, und was er sagt, das muß Ja sein. Ja, wenn er weiter gebeten hätte, hätte er auch müssen gewährt sein. Da ist kein Nein, sondern eitel Ja, uns zum Exempel vorgeschrieben, nicht daß wir allein sehen auf die Heiligen im Himmel, sondern die auf Erden sind, wie Abraham dasmal noch auf Erden lebte. Darum lasset uns solch Exempel fassen, und nur getrost bitten. Ist es gebeten, so ist es auch erhört; kannst du nur glauben, so muß es geschehen. Ja, sagst du, hätte ich so starken Glauben wie Abraham. Antwort: Ja, ist es doch nicht um seinetwillen geschehen; darum gilt es nicht, daß du sagst: Wenn ich so heilig wäre. Glaube wie er, so bist du auch so heilig. Woher ist er so heilig? Sagt doch Josua [Cap. 24, 2.], er habe auch Abgötter angebetet, ist auch ein Sünder gewesen; dadurch ward er aber fromm, daß er glaubte; wie gehört. So glaube nur, so bist du auch so fromm. Wir gehen aber in dem Sinne, er habe es erworben mit großen Werken und heiligem Leben; [es] sind aber eitel Träume. Hätte er den Glauben nicht gehabt, so hätte er nichts überall erworben. So müssen wir auch durch den Glauben solches erlangen. Es ist überaus ein groß Ding um das Exempel, daß ein Christ solche Gewalt hat, daß er kann bitten, und GOtt ihm nichts kann abschlagen.

20. Das hat Jacobus in seiner Epistel auch angezogen, und führt es stark; ist auch das beste Stück in derselben Epistel: [Des Gerechten Gebet (sagt er [Cap. 5, 16—18.]) vermag viel, wenn es thätig ist. Elias war ein Mensch gleich wie wir, und er betete ein Gebet, daß es nicht regnen sollte; und es regnete nicht auf Erden drei Jahre und sechs Monden. Und er betete abermal, und der Himmel gab den Regen, und die Erde brachte ihre Frucht“,] setzt eigentlich das Wort hinzu: „Er war ein Mensch wie wir.“ So sollen wir auch dahin deuten alle Historien und Exempel, daß sie Menschen gewesen sind, sowohl als wir. Denn die Larve betrügt uns nur, daß wir mit den Heiligen so ein Geplärr machen. [Wir] sollten so sagen: Sind sie heilig, so sind wir auch heilig; sind wir Sünder, so sind sie es auch gewesen; sind wir doch alle eben von dem Fleisch und Blut geboren, und hat uns GOtt sowohl geschaffen, als sie; es ist je ein Mensch des andern werth, und unterscheidet nichts, denn der Glaube. Hast du den Glauben und GOttes Wort, so bist du ebenso groß, darfst nicht sorgen, daß du geringer seiest denn er, ohne daß der Glaube nicht so stark kann sein.

21. Und sonderlich ist zu merken das, daß solch Exempel anzeigt, wie kräftig das Gebet zur Besserung der Guten ist. Man findet wohl Exempel, wie stark das Gebet gewesen ist zum Schaden der Widersacher, wie Mosis, der doch der sanftmüthigste Mann gewesen ist, ein kurz Gebet that,\footnote{Erlanger: thät er.} daß Kore, Dathan und Abiram die Erde verschlang [4 Mos. 16, 32.]. Desgleichen liest man auch von David und andern mehr. So können die Heiligen auch zuweilen wohl schaden und strafen. Aber dies Exempel ist noch viel edler, daß das Gebet so stark ist, wenn es soll zur Besserung des Nächsten dienen, und GOtt mehr geneigt ist zu erhören das, so nütz und gut, denn das schädlich ist.

22. Elias, der Prophet, trieb es beiderlei stark. Er hat, daß es vierthalb Jahre nicht regnete [1 Kön. 17, 1. 7.], [es] geschah auch also. [Er] that es aber darum, daß das Land so voll falscher Propheten und Abgötterei war, daß sein Geist erzürnt und erbittert ward, und die Strafe über das Land brachte, darum, daß der König so zornig auf ihn ward, daß er ihn
\quellseite{163}
in allen Ländern ließ suchen, ihn zu erwürgen [1 Kön. 18, 10.]. So siehst du, daß sie können erlangen, was sie bitten; doch mehr und stärker zum Guten, denn zum Schaden. Deß sollen wir uns trösten und stärken, und sicher sein, daß so geschieht, wie wir bitten.

23. So sagt Johannes auch in seiner [ersten] Epistel [Cap. 5, 14. 15.]: „Das ist die Freudigkeit, die wir haben zu GOtt, daß, so wir etwas bitten nach seinem Willen, so höret er uns; und so wir wissen, daß er uns höret, was wir bitten, so wissen wir, daß wir die Bitte haben, die wir von ihm gebeten haben.“ Das ist, wenn wir glauben, so sind wir freudig und fest, daß uns GOtt nicht läßt; wie Abraham nicht zweifelte, bis er auf die „zehn“ kommt. Hätte er weiter gebeten, wäre er auch erhört worden. Hätte ich einen Christen, der für mich bäte, [so] wollte ich gutes Muths sein, und mich vor niemand fürchten. Hätte ich auch einen, der wider mich hätte, so wollte ich auch lieber den türkischen Kaiser zum Feinde haben.

\section{Das neunzehnte Capitel}

\begin{schriftzitat}
V. 1—38. Die zween Engel kamen gen Sodom des Abends. Lot aber saß zu Sodom unter dem Thore. Und da er sie sahe, stund er auf ihnen entgegen, und bückte sich mit seinem Angesichte auf die Erde, und sprach: Siehe, mein HErr, kehret doch ein zum Hause eures Knechtes, und bleibet über Nacht, lasset eure Füße waschen, so stehet ihr morgen frühe auf, und ziehet eure Straße. Aber sie sprachen: Nein, sondern wir wollen über Nacht auf der Gasse bleiben. Da nöthigte er sie fast, und sie kehreten zu ihm ein, und kamen in sein Haus. Und er machte ihnen ein Mahl, und buk ungesäuerte Kuchen, und sie aßen. Aber ehe sie sich legten, kamen die Menschen der Stadt Sodom und umgaben das Haus, jung und alt, das ganze Volk aus allen Enden. Und forderten Lot, und sprachen zu ihm: Wo sind die Männer, die zu dir kommen sind diese Nacht? Laß sie heraus gehen zu uns, daß wir sie erkennen. Lot ging heraus zu ihnen vor die Thüre, und schloß die Thüre hinter ihm zu, und sprach: Ach, lieben Brüder, thut nicht so übel. Siehe, ich habe zwei Töchter, die haben noch keinen Mann erkannt, die will ich heraus geben unter euch, und thut mit ihnen, was euch gefällt; allein diesen Männern GOttes thut nichts, denn darum sind sie unter den Schatten meines Daches eingegangen. Sie aber sprachen: Komm hieher. Da sprachen sie: Bist du einiger nicht herein kommen als ein Fremdling, und willst nun Richter sein? Wohlan, wir wollen dir mehr Uebels thun, denn jenen. Und sie drungen hart auf den Mann Lot, und da sie hinzu liefen, und wollten die Thüre aufbrechen, griffen die Männer hinaus und zogen Lot hinein zu ihnen ins Haus, und schlossen die Thüre zu. Und die Männer vor der Thüre am Hause wurden mit Blindheit geschlagen, beide klein und groß, daß sie die Thüre nicht finden konnten. Und die Männer sprachen zu Lot: Hast du noch irgend hier einen Eidam, und Söhne und Töchter, und wer dich angehört in der Stadt, den führe aus dieser Stätte; denn wir werden diese Stätte verderben, darum, daß ihr Geschrei groß ist vor dem HErrn; der hat uns gesandt, sie zu verderben. Da ging Lot hinaus, und redete mit seinen Eidamen, die seine Töchter nehmen sollten: Machet euch auf, und gehet aus diesem Orte, denn der HErr wird diese Stadt verderben. Aber es war ihnen lächerlich. Da nun die Morgenröthe aufging, hießen die Engel den Lot eilen, und sprachen: Mache dich auf, nimm dein Weib, und deine zwei Töchter, die vorhanden sind, daß du nicht umkommest in der Missethat dieser Stadt. Da er aber verzog, ergriffen die Männer ihn, und sein Weib, und seine zwei Töchter bei der Hand, darum, daß der HErr sein verschonte, und führten ihn hinaus, und ließen ihn außen vor der Stadt. Und als sie ihn hatten hinaus gebracht, sprachen sie: Errette deine Seele, und siehe nicht hinter dich, auch stehe nicht in dieser ganzen Gegend; auf dem Berge errette dich, daß du nicht umkommest. Aber Lot sprach zu ihnen: Ach, nein, mein HErr, siehe, dieweil dein Knecht Gnade gefunden hat vor deinen Augen, so woll-


\quellseite{164}
hast du deine Barmherzigkeit groß machen, die du an mir gethan hast, daß du meine Seele bei dem Leben erhieltst; ich kann mich nicht auf dem Berge erretten, es möchte mich ein Unfall ankommen, daß ich stürbe. Siehe, da ist eine Stadt nahe, darein ich fliehen mag, und ist klein, daselbst will ich mich erretten; ist sie doch klein, daß meine Seele lebendig bleibe. Da sprach er zu ihm: Siehe, ich habe auch in diesem Stücke dich angesehen, daß ich die Stadt nicht umkehre, davon du geredet hast. Eile, und errette dich daselbst; denn ich kann nichts thun, bis daß du hinein kommst. Daher ist die Stadt genannt Zoar.\footnote{Randglosse: Zoar heißt klein. — Diese Randglosse fehlt in der Jenaer. In der Wittenberger steht nur der Text der ersten drei Verse.} Und die Sonne war aufgegangen auf Erden, da Lot gen Zoar einkam. Da ließ der HErr Schwefel und Feuer regnen von dem HErrn vom Himmel herab auf Sodom und Gomorra, und kehret die Städte um, die ganze Gegend und alle Einwohner der Städte, und was auf dem Lande gewachsen war. Und sein Weib sahe hinter sich, und ward zur Salzsäule. Abraham aber machte sich des Morgens frühe auf an den Ort, da er gestanden war vor dem HErrn, und wandte sein Angesicht gegen Sodom und Gomorra und alles Land der Gegend, und schaute, und siehe, da ging ein Dampf auf vom Lande, wie ein Dampf vom Ofen. Denn da GOtt die Städte in der Gegend verderbete, gedachte er an Abraham, und geleitete Lot aus den Städten, die er umkehrette, darinnen Lot wohnte. Und Lot zog aus Zoar, und blieb auf dem Berge mit seinen beiden Töchtern; denn er fürchtete sich, zu Zoar zu bleiben, und blieb also in einer Höhle mit seinen beiden Töchtern. Da sprach die Aelteste zu der Jüngsten: Unser Vater ist alt, und ist kein Mann mehr auf Erden, der uns beschlafen möge, nach aller Welt Weise. So komm, laß uns unserm Vater Wein zu trinken geben, und bei ihm schlafen,\footnote{Erlanger: und mit ihm trunken werden.} daß wir Samen von unserm Vater erhalten. Also gaben sie ihrem Vater Wein zu trinken in derselben Nacht. Und die erste ging hinein, und legte sich zu ihrem Vater, und er ward es nicht gewahr, da sie sich legte, noch da sie aufstund. Des Morgens sprach die Aelteste zu der Jüngsten: Siehe, ich habe gestern bei meinem Vater gelegen, laß uns ihm diese Nacht auch Wein zu trinken geben, daß du hinein gehest, und legest dich zu ihm, daß wir Samen von unserm Vater erhalten. Also gaben sie ihrem Vater die Nacht auch Wein zu trinken. Und die Jüngste machte sich auch auf, und legte sich zu ihm, und er ward es nicht gewahr, da sie sich legte, noch da sie aufstund. Also wurden die beiden Töchter Lots schwanger von ihrem Vater, und die Aelteste gebar einen Sohn, den hieß sie Moab; von dem kommen her die Moabiter, bis auf den heutigen Tag. Und die Jüngste gebar auch einen Sohn, den hieß sie das Kind Ammi; von dem kommen die Kinder Ammon, bis auf den heutigen Tag.
\end{schriftzitat}



1. Zum ersten hören wir hier von den Männern oder Engeln, die von Abraham gegangen waren gen Sodom. Wie nun Abraham, sein Vetter, war, so ist auch Lot zu Sodom. Abraham war ein Mann, der gerne herbergte, also thut dieser hier auch, ja, daß er sie auch bringt und nöthigt, daß sie müssen in sein Haus kehren, und seine Gäste sein.

2. Ist aber ein Exempel, darinne man sehe, was die Leute für einen Brauch haben gehabt, daß sie jedermann bereit gewesen sind zur Wohlthat. Denn man soll nicht denken, daß sie allein auf diesmal so willig gewesen sind, arme, fremde Leute zu herbergen. Sie sind zwar auch selbst Fremdlinge gewesen, haben wohl versucht, wie fremden Leuten zu Muthe ist, und wie sich niemand ihrer annimmt; darum haben sie solche gerne aufgenommen. Es wird aber ohne Zweifel also zugegangen sein, daß sie auch manche undankbare Gäste gehabt haben. [Es] gehört sich auch also. Denn was ein christlich gut Werk sein soll, muß also geschickt sein, daß es die Leute, denen es wohlthut, lasse alle Schalckheit dagegen thun. [Sie] haben so gedacht: Um GOttes willen geschehen, um Menschen willen nicht gelassen. Denn man muß immer das ansehen, daß GOtt Gefallen an den Leuten gehabt habe. Wovon er Zeugniß gibt, das muß recht und köstlich sein.

3. Also haben sie nun alle beide, Lot und Abraham, gesessen und gewohnt zwischen eitel Dornen, Heern und Scorpionen, wie denn alle Christen und GOttes Kinder. Da wird nichts anders aus; es muß gehen, wie im Hohenliede [Cap. 2, 2.] steht: „Wie eine Rose mitten unter den Dornen, so ist meine Liebe unter den Töchtern.“ So ist er ein einiger freier Mensch geweisen
\quellseite{165}
, voll guter Werke, [hat] jedermann wohlgethan, und ihm jedermann Leid und Schaden gethan, daß ihm sein Leben bitter und sauer genug geworden ist. Denn er mußte ein ander Wesen führen, denn die Welt; hat müssen leben unter denen, die nicht recht fromm sind gewesen, und die Ungläubigen strafen und auf den Hals laden; daß sein Leben nicht mehr, denn immer Haß und Neid hat verdient. Darum hat ihn GOtt desto mehr getröstet, auch durch die Engel; hat also mit ihm zu schaffen, als trüge er ihn auf den Händen; läßt ihn doch so jämmerlich stecken unter den Ungläubigen, daß man freilich diese Erempel nimmer genug preisen kann.

4. Da wohnt der fromme Lot in der Hauptstadt Sodom von den fünfen. Und rechne selbst, wie er da sitzt mit Weib, Kind und Gesinde, hat sich müssen schicken in die Bürger, die da waren; hat Töchter wollen vergeben, auch zwei Eidame gehabt, mit denen hat er sich müssen vertragen. Wiederum haben sie doch ein solch schändlich Wesen, Fressen und Saufen, und alle Untugend getrieben, daß sie ihn nicht anders gehalten haben, denn für einen Aschenbrödel; solches\footnote{„solches“ fehlt in der Erlanger.} hat er müssen alles leiden. Wie oft hat er auch die Leute strafen müssen. Denn, wie vor gesagt, es gebührt einem Christenmenschen, daß er wider die Bösen predige und lehre, und nicht stille schweige. So hat er nun müssen großen Glauben haben, daß er\footnote{Erlanger: er ihn.} hat leben und sich behelfen können unter so giftigen und\footnote{„und“ fehlt in der Erlanger.} verzweifelten Leuten.

5. Das zeucht der Apostel Petrus für ein Exempel an: „GOtt hat erlöset (spricht er [2. Ep. 2, 7. 8.]) den rechtfertigen Lot, der übertäubet war von dem unzüchtigen Wandel der Greulichen. Denn dieweil er rechtfertig unter ihnen wohnte, daß er es sehen und hören mußte, quälten sie die gerechte Seele von Tag zu Tage mit ihren unrechten Werken.“ Er sahe, daß sie solche greuliche Untugend führten, keinem Menschen dienten noch halfen; sondern, wie Christus [Luc. 17, 28.] von ihnen sagt: „Sie aßen, sie tranken, sie kauften, sie verkauften, sie pflanzten, sie baueten“, jedermann wandelte seiner Lust und Genuß nach, sahe keiner auf den Nächsten. Er war allein in der ganzen Stadt, half und diente den Leuten, so viel er konnte; mußte dennoch von ihnen gehaßt und unterdrückt werden. So hat er müssen sehen, daß ihm im Herzen wehe gethan hat, ja, daß ihm das Herz hätte mögen brechen; dennoch muß er es leiden, kann nicht wehren noch helfen.

6. Darum hat es der Prophet Ezechiel auch angezeigt: „Siehe, das war die Untugend deiner Schwester Sodom (spricht er [Cap. 16, 49.] zur Synagoge), Hoffahrt, Müßiggang und Ueberfluß, daß sie vollauf zu fressen und saufen hatten, sie und ihre Töchter“, das ist, die Flecken und Dörfer umher, „und den Dürftigen reichte niemand die Hand“, darüber fuhren sie heraus und thaten solche Laster, daß ich sie mußte zu Asche machen. So beschreibt es Christus auch (wie oben gesagt [Luc. 17, 28—30.]): „Wie es war zu Zeiten Lots, so wird es auch sein in den Tagen des Menschen-Sohnes; sie aßen, tranken, kauften, verkauften, pflanzten und baueten, bis daß Lot aus Sodom ging; da regnete es Feuer und Schwefel, und brachte sie alle um.“ Wenn du es nun ansiehst, so ist eben das der Gebrechen gewesen zu Sodom, wie noch heutiges Tages, daß jedermann allein auf sich trachtet und zu sich scharrt, mit Recht und Unrecht; es bleibe ein anderer, wo er kann. GOtt hat es gegeben, daß man bauet und pflanzt; aber daß man allein alles, was man kann, zu sich rafft, und den Nächsten läßt Noth leiden, nimmt sich seines andern Menschen an, das ist die Klage und Schuld alle mit einander.

7. Derhalben spricht der Prophet zu der Stadt Jerusalem: „Siehe, so thust du auch, ja, „Sodom und Samaria hat noch nicht die Hälfte gethan, das du thust“, daß sie noch fromm sind, wenn man sie gegen dir hält [Ezech. 16, 51.], wie auch Christus die Juden straft, und fällt ein Urtheil über sie, und spricht [Matth. 11, 23. 24.]: „Und du Capernaum, du bist erhoben bis an den Himmel, du wirst bis in die Hölle hinunter gestoßen werden. Denn so zu Sodoma die Thaten geschehen wären, die bei dir geschehen sind, sie stünde noch heutiges Tages. Doch, ich sage euch, es wird der Sodomer Lande träglicher ergehen am jüngsten Gerichte, denn dir“, also, daß wir uns nicht dürfen rühmen. Hat GOtt die Stadt verfeucht um ihrer Sünde willen, so wird er uns auch nicht schonen; wie St. Petrus [2. Ep. 2, 6.] auch sagt: „Er hat die Städte Sodom und Gomorra zu Asche gemacht, umgekehret, und verdammt, damit ein Erempel gesetzt
\quellseite{166}
denen, die zukünftig Gottlose sein werden.“ Wahr ist es, daß sie die Sünde gethan haben, so man die stummen Sünden nennt; aber sonderlich vergleicht’s die Schrift allen denen, so das Evangelium hören, verstehen, und folgen nicht, und ist eben denen zum Exempel gesetzt, daß sie sich daran stoßen.

8. So soll man die Exempel ansehen, wie die frommen Kinder Abraham und Lot. Als sie Fremdlinge waren unter den bösen Leuten, da sie nichts haben, denn das ihnen das Herz möchte gebrochen haben, schändlich Wesen, Untreue und Unrecht, noch mußten sie bei ihnen bleiben, konnten aus der Welt nicht laufen; thaten nicht wie unsere Narren, die Mönche, die nicht in der Welt wollen leben. GOtt hatte sie in das Land geführt, [sie] mußten darinne sein; wenn sie gleich anderswo gewohnt hätten, hätten sie es eben so gefunden. Wo man hinkommt, ist der Teufel daheim. Darum sind es große Exempel, die uns GOtt vorbildet, daß wir am Glauben und Liebe fest halten, wider die ganze Welt. [Es] ist ein groß Kreuz, wenn wir sehen, daß niemand Freund ist, müssen eitel Neid und Haß leiden; noch läßt GOtt den frommen Mann fallen, wie wir hören werden, wie und warum er es thut.

9. Das ist nun das große, greuliche Gericht GOttes über die fünf Städte. Denn wiewohl Abraham so mit großem Ernst dafür bat, doch, weil er so gar niemand fand, der gut war, ließ er es fort anstehen und den Zorn vor sich gehen. Darum das Exempel nicht in Wind zu schlagen, sondern wohl in die Augen zu bilden ist\footnote{Erlanger: „ist. Dies ist“ etc. — Wittenberger: den großen Ernst.} der große Ernst GOttes, wie gesagt ist, daß man sich nicht deß tröste, daß man sagt: es seien noch viel fromme Leute unter uns, um deren willen Gott den Zorn aufhält. Wie sie da waren, so sind sie noch, und viel ärger; doch sind mehr fromme Leute da gewesen. Also, daß niemand darf einen Deckel und Schutz nehmen, daß er volle Gottes Gericht entlaufen.

10. Es ist endlich beschlossen, daß er haben will, daß man sich vor ihm fürchten soll, und nicht sicher fahren. Darum ist schon der Zorn angegangen, wo er sieht, daß man es in Wind schlägt, und sich gleichwohl auf anderer Leute Fürbitte will verlassen. Daher sagt die Schrift [Ps. 147, 11.]: „Er hat ein Wohlgefallen an denen, so ihn fürchten, und auf ihn trauen“; wie auch im Evangelio Matthäi [Cap. 18, 26. 27.] von dem Knechte, der seinem Herrn zehntausend Pfund schuldig war, daß er Gnade und Barmherzigkeit erlangte, und den Herrn seiner jammerte, als er niederfiel und [um] Gnade bat. So ist durch und durch das Exempel zu führen, daß GOtt niemand gnädig und barmherzig ist, denn die ihn fürchten, daß er zürnen\footnote{Erlanger: zürnet.} und richten werde. So haben die Propheten das Exempel angezeigt, und wohl gehalten und vorgehalten, daß sie die Leute in Gottesfurcht behielten. Es ist zu einem Schrecken und Zittern gesetzt, wie wir aus Petro [2. Ep. 2, 6.] gehört haben, allen Gottlosen, welchen es gewißlich auch also gehen wird, wie es Sodoma und Gomorra gegangen ist, es verziehe sich, wie lange es wolle. GOtt kann ihrem Wesen wohl eine Zeitlang zusehen, aber zuletzt kommt er mit greulicher Strafe, er bleibt gewiß nicht außen.

11. So ist es nun gesagt, was die Sünde sei gewesen, darum sie gestraft sind. Oben ist gesagt im dreizehnten Capitel [§ 4], daß die Städte gelegen seien in einem fetten, köstlichen Lande und rechter Schmalzgrube, da aller Dinge die Fülle und Genüge gewesen ist, daß Mosis [Cap. 13, 10.] sagt, als ein Garten GOttes; so schön von Bäumen und allen edlen Früchten, Korn, Wein, Oel, als hätte es GOtt selbst erbauet, und ihm einen Garten zugerichtet; darum das Volk, so in der Stadt war, hatte einen guten Muth, meinten, sie hätten einen gnädigen GOtt, schlugen alle Zucht aus. Solche Sicherheit brachte sie zu solcher greulichen Sünde; damit GOtt angezeigt hat, was Reichtum und Ueberfluß Gutes bei den Menschen wirkt. Jedermann strebt darnach, daß wir genug und die Fülle haben, und gut Leben; aber die Schrift lehrt, daß uns nichts Fährlicheres begegnen mag [Pred. 2, 1.]. Darum wäre es besser, daß alles auf das theuerste, und nicht so großer Ueberfluß da gewesen wäre, [so] wären sie wohl fromm geblieben; aber weil sie alle Fülle hatten, und im Saufe lebten, wie Christus [Luc. 17, 27.] spricht, wurden sie also gestraft.

12. Derhalben sollten wir auch denken, daß einem Menschen nichts besser ist denn Armuth, daß wir nicht überläng\footnote{„überläng“ = überflüssig, mehr als genug.} haben, sondern nur Hülle und Fülle, um und an, damit davon.
\quellseite{167}
Denn, wie man spricht und wahr ist: Es müssen gar starke Beine sein, die gute Tage ertragen sollten; und: Ein Mensch kann allerlei leiden, ohne gute Tage; wenn er zu viel Futter hat, so geht er eben, wie der Esel, auf das Eis, und bricht ein Bein vor Wohltagen. Sind etliche reich gewesen und fromm, so ist es eine sonderliche Gnade; sonst ist es nicht nütz im gemeinen Haufen, daß einer reich sei; wie Paulus zu Timotheo schreibt [1. Ep. 6, 8.]: „Wenn wir aber Futter und Decke haben, so lasset uns begnügen.“ Wir sehen täglich an\footnote{Jenaer: in.} aller Erfahrung, daß doch jedermann dahin stirbt, bringen nichts davon, denn daß wir den Bauch füllen, decken und hüllen. Darum sollten wir nicht also scharren. Die nach uns kommen, werden auch wohl arbeiten.

13. Weiter haben wir gehört, wie GOtt den frommen Lot mit Wunderzeichen aus der Stadt führt; nicht allein erlöst von der Sünde, sondern auch von der Strafe. Welches Exempel St. Petrus hoch führt in seiner [zweiten] Epistel [Cap. 2, 7.] zum Trost allen, so da glauben. GOtt hat ein Wohlgefallen an ihm, nimmt sein sonderlich wahr; vornehmlich auch um Abrahams willen, denn er ist mit ihm in gleichem Glauben gestanden. Darum zeigt er damit an, wie er Achtung hat auf die ihm trauen. Denn, stelle es vor Augen, so findest du ihn mit seinem Weibe und zwo Töchtern, das sind vier Personen in der ganzen Stadt und Gegend; die andern allzumal haben gelebt, wie sie wollten, gesauft\footnote{Erlanger: gesauft.} und gepraßt. Er müßte ihrer aller Aschenbrödel sein, wie sie zu ihm sagen im Text.

14. Da hat er müssen leiden die zwei Stücke: Zum ersten, wie Petrus sagt, daß sein Herz hätte brechen mögen gegen dem Leben, das er da sahe. Sein Glaube und Geist konnte es nicht leiden, daß er nicht dawider redete; darum er sich hat müssen von ihnen absondern und sie\footnote{„sie“ fehlt in der Erlanger.} auch strafen. So hat er auf sich geladen Haß und Feindschaft. [Es] ist nicht alles geschrieben, was er gelitten hat; St. Petrus rührt es aber ein wenig, als er sagt: „Er hat da müssen sehen und hören, daß seine Seele quälte.“ Wäre er ein einiger Mann gewesen, hätte er mögen davon gehen; so ist er mit Weib und Kind beladen, um deswillen muß er zusehen und leiden. Dar-

um ist es ein mächtiger Preis, daß er unter den Leuten sitzen kann, und allein wider sie redet und thut. Was GOtt lobt, das muß hoch und köstlich sein; er gibt es wohl zu [be]denken, was er hat müssen leiden, mehr denn angezeigt ist, sintemal wir vor Augen sehen, was die noch leiden müssen, die nicht leben, wie der Teufel und die Welt will.

15. Neben solchen Leiden hat er dennoch einen Trost dagegen. Weil er wider viele Leute muß leben, ist GOtt bei ihm, und hat ein Auge auf ihn, straft sie greulich, und errettet ihn, läßt ihn Unglück genug leiden; aber er erhält ihn, beide vor Sünde und Strafe, auf daß wir daran lernen, daß kein Wunder ist, ob er uns die ganze Welt läßt feind sein. Wir müssen uns ermägen, daß Einer müsse allein wider alle Welt stehen.

16. Es kommt niemand so in den Himmel, denn daß er sich lasse dünken, er fahre allein, wie es Christo und den Propheten auch begegnet hat; wie er spricht im Psalter [Ps. 142, 5.]: „Ich habe mich umgesehen, aber da war niemand, der mich kennete.“ Geschieht es nicht vor der Welt im Leben, so muß es am Tode geschehen, daß ich da stehe, und alles verdamme, was wider mich ist, als sei ich klüger denn die ganze Welt; also muß es gehen, sonst ist es nicht recht. Ist alles darum geschrieben, daß wir die Augen zuthun, uns nicht wundern, ob ihrer viel das Evangelium lästern. Es ist noch Wunder, wenn du nur zween siehest, die Christen sind. Wir mögen ihrer wohl viel loben und dafür ansehen, aber wenn es zum Treffen geht, so sei nur gerüstet, daß du könnest sprechen: Hier muß ich allein und verlassen stehen, und hält es niemand mit mir. Darum gibt GOtt diese Exempel, der wir uns trösten, daß nicht allein uns also geschehen, sondern andern heiligen Leuten auch widerfahren ist, und darauf verlassen, daß GOtt ein Auge auf uns hat, und über uns hält, wie er mit ihnen gethan hat.

17. Auch ist noch ein tröstlich Exempel hier, in dem, daß dieser Lot hernach seine Töchter beschläft. Da spricht der Text, wie ihn die Töchter haben trunken gemacht, und sich zu ihm gelegt, daß er ihrer nicht gewahr worden ist, wann sie sich gelegt haben, oder aufgestanden sind. Was ist GOtt daran gelegen, daß er einen solchen keinen Patriarchen, den großen, geistlichen Mann, läßt so schändlich zum Narren werden? Es ist je wunderlich, aber eben darum
\quellseite{168}
gesehen, daß die Vernunft über solchen närrischen Historien zum Narren werde. St. Paulus hat es wohl gesagt: Seine Gerichte sind der Vernunft unergründlich [Röm. 11, 33.]. Er treibt es also mit den Heiligen, daß er alle Klugheit zu Narren macht [1 Cor. 1, 27.].

18. Dieser Mann nähme nicht aller Welt Gut, daß er ein ander Weib berührte, und schläft bei seinen eigenen Töchtern; und sie dürfen mit guter Vernunft und Bedacht den Vater betrügen. Was soll man doch dazu antworten? Wir wollen es unverantwortet lassen, sondern also sagen, daß die zwo Töchter greulich gesündigt haben, ob sie wohl einen Schein vorwenden, und aus Furcht-thun, meinten, es wären keine Leute mehr auf Erden; haben es in der närrischen Andacht gethan, daß sie wollten der Welt helfen. Aber gute Meinung taugt nichts, man muß zuvor GOttes Wort haben, und gewiß sein, daß es wohl gethan sei. So will ich den Lot nicht entschuldigen, wiewohl er nicht so schmerzlich strauchelt; sie aber, die Töchter, will ich viel weniger entschuldigen. Aber dafür halte ich es, daß sie auch einen starken Glauben gehabt haben, sonst wären sie nicht errettet worden, dieweil sein eigen Weib nicht geschonet ward. Er wird es auch ohne Zweifel nicht gelassen haben, seine Kinder zu unterrichten, wie sie glauben sollten, daß des Verstandes halben kein Mangel gewesen ist, daß sie da gefallen sind.

19. Das ist nun auch zum Exempel geschrieben, daß niemand so heilig sei, noch so fest stehe, der nicht wieder fallen könnte.\footnote{Erlanger: könne; Jenaer: künde; Wittenberger: künde.} Hat der Mann können so hoch fallen, was ist es Wunder, ob wir straucheln? Doch ist damit niemand ein Deckel noch Behelf gegeben zu sündigen, sondern ein Trost gestellt denen, so da glauben, daß sie nicht verzagen, ob sie zuweilen fallen. Es ist ein Stück, das da gehört in das Evangelium und Reich Gottes, da eitel Vergebung der Sünden ist. Falle hin oder her, verzweifle nur nicht, und stehe wieder auf.

20. Also werden wir hernach [Cap. 27, 35.] hören, daß Jakob, der Patriarch, auch strauchelte, als er blöde und verzagt war,\footnote{„war“ fehlt in der Erlanger.} und sich kümmerlich stellte. Denn, hätte uns GOtt nicht also vorgebildet, daß die Heiligen also genarrt haben, so könnten wir kein Königreich nicht lernen kennen, wie es nichts Anderes denn Vergebung der Sünden ist. Darum ist es so geschrieben, daß der Preis und Ehre stehe, nicht auf Menschenwerken, sondern auf GOttes Gnade. Lot und Abraham konnten so böse werden als ich, und ich wiederum so fromm als sie. Darum haben wir gehört, wie er die Sodomiter gar nicht verachtet, sondern sich ihrer annimmt, vermahnt sie, und wollte gern, daß sie sich besserten.

21. So muß es zugehen im Reich Christi, daß da eitel aufrichten, immer sich bessern und bessern sei,\footnote{„sei“ fehlt in der Erlanger.} nur daß man nicht sicher und frech werde, und sich in die Schanze schlage, sondern wieder zum Kreuz krieche. Es darf sich keiner des andern schämen. Der mich aus Fleisch und Blut gemacht, hat jene auch also gemacht; und wie er jene heilig gemacht hat, so kann er mich auch machen. Natur ist in allen gleich, Gnade ist auch gleich. Siehe, das sollen solche Exempel der Väter lehren, darum soll man sie nicht so hoch heben, wie unsere Narren, die da meinen, daß sie gar keine Sünde gefühlt haben.

22. Zuletzt, daß wir es beschließen, wollen wir kurz die heimliche Deutung des Capitels überlaufen. Es sind in dieser Historie zwei Stücke, welche deuten beide nichts denn Menschenlehre, wie es der Prophet Ezechiel [Cap. 16, 51.] selbst auslegt, als er Jerusalem und die Sodomiter zusammen vergleichet, und sagt dazu, daß sie viel ärger worden ist, denn Sodom. Das geschieht nun also, daß man den heiligen Samen, das Wort GOttes, mißbraucht und verderbt, daß er nicht kann Frucht schaffen. Das ist das erste.

23. Das andere, daß, wie die zwei Töchter ihrem Vater nachhuren, und von ihm schwanger werden, also thun auch alle, so mit Werken vor GOtt kommen, und vom freien Willen lehren, brauchen den Samen nicht, dahin sich’s gehört. Das ist nun gleich so viel, als wenn die Töchter bei dem Vater schläft, und Hurentinder zeugt. Lot ist lex, das Gesetz; die zwei Töchter populus, des Gesetzes Volk, die machen den Vater trunken; davon Jesaias [Cap. 28, 1.] sagt, von dem trunkenen Ephraim. Das geht also zu: Wenn wir GOttes Wort haben, und geben ihm einen andern Verstand, so machen wir es trunken. Der Wein ist Menschenlehre, so bei uns
\quellseite{169}
möchst und gemacht wird, also wird man denn schwanger von dem trunkenen Gesieß, das ist, salich ausgelegt und verstanden. Wie sie es nun verstehen, so tragen sie Frucht, Ammon und Moab, solche Kinder, die nimmer unter die Christenheit und GOttes Volk kommen dürfen. Wie ihnen auch leiblich verboten war, nicht unter die Kinder Israel zu kommen etc. Denn es gehört ihrer keiner unter Christum, die also gelehrt werden nach dem Verstande, so wir selbst erfunden haben.

24. Soll jemand zu Christo kommen, der muß rein und lauter das Gesetz und Evangelium verstehen, und rechte Frucht tragen. Nun lautet das Gesetz also, daß [es] niemand aus eigenen Kräften thun kann: [es] ist auch nicht darum gegeben, sondern daß es uns zu unserer Erkenntniß bringe. Darnach, daß man GOttes Gnade anrufe, das ist das Evangelium, so bleibt es rechtschaffen und nüchtern, daraus nicht Ammoniten und Moabiten, sondern rechte Kinder Israel werden.

25. Das deuten auch die zween Namen: [\hebraeisch{מוֹאָב}] Moab, der Aelteste, heißt auf deutsch vom Vater, als sollte sie sagen: Den Sohn habe ich vom Vater; [\hebraeisch{עַמּוֹן}] Ammon, filius doloris, Schmerzenssohn, oder [\hebraeisch{בֶּן־עַמִּי}] BenAmmi, das ist, ein Kind meines Volks. Das sind die Werktheiligen, die durch ihre guten Werke wollen Kinder GOttes und die Nächsten bei ihm sein, und doch keinen Glauben haben. Darum sagt auch Jesaias [Cap. 16, b.] von dem Volk: Moab ist ein stolzer Esel, ich kenne seinen Hochmuth wohl, [er] ist fast stolz, und rühmet sich mehr, denn er vermag. So sind sie alle, die aus Werken sind, die ihr Ding allein wollen für recht gehalten haben. Das sind auch die Ammoniten. Vor der Welt scheint ihr Ding köstlich, da ist es Moab; inwendig aber im Gewissen ist kein Friede noch Zuversicht zu GOtt, keine Freude noch Lust. Darum ist es inwendig ein elend Wesen, wie es außen köstlich ist. Vor GOtt sind sie Ammon, und vor der Welt Moab. Denn durch Gesetz und Werke kann dem Gewissen nimmermehr geholfen werden, daß es Ruhe und Frieden habe. Das ist die Deutung dieser Historie; also, daß alles dahin geht und darum zu thun ist, daß man halte am reinen GOttes Wort und nichts Anderes höre denn dasselbige, und Menschenlehre kurzum verdamme.

\section{Das zwanzigste Capitel}

\begin{schriftzitat}
V. 1—12. Abraham aber zog von dannen in das Land gegen Mittag, und wohnete zwischen Kades und Sur, und ward ein Fremdling zu Gerar. Und sprach von seinem Weibe Sara: Es ist meine Schwester. Da sandte Abimelech, der König zu Gerar, nach ihr, und ließ sie holen. Aber GOtt kam zu Abimelech des Nachts im Traum, und sprach zu ihm: Siehe da, du bist des Todes, um des Weibes willen, das du genommen hast; denn sie ist eines Mannes Eheweib. Abimelech aber hatte sie noch nicht berührt, und sprach: HErr, willst du denn auch ein gerecht Volk erwürgen? Hat er nicht zu mir gesagt, sie ist meine Schwester? Und sie hat auch gesagt, er ist mein Bruder. Habe ich doch das gethan mit einfältigem Herzen und unschuldigen Händen. Und GOtt sprach zu ihm im Traum: Ich weiß auch, daß du mit einfältigem Herzen das gethan hast; darum habe ich dich auch aufgehalten, daß du nicht wider mich sündigtest, und habe es dir nicht zugegeben, daß du sie betastest. So gib nun dem Manne sein Weib wieder, denn er ist ein Prophet, und laß ihn für dich bitten, so wirst du lebendig bleiben. Wo du aber sie nicht wieder gibst, so wisse, daß du des Todes sterben mußt, und alles, was dein ist. Da stund Abimelech des Morgens früh auf, und rief allen seinen Knechten, und sagte ihnen dieses alles vor ihren Ohren, und die Leute fürchteten sich sehr. Und Abimelech rief Abraham auch, und sprach zu ihm: Warum hast du das gethan? und was habe ich an dir gesündigt, daß du so eine große Sünde wolltest auf mich und mein Reich bringen? Du hast mit uns gehandelt, nicht wie man handeln soll. Und Abimelech sprach weiter zu Abraham: Was hast du angesehen, daß du solches gethan hast? Abraham sprach: Ich dachte, vielleicht ist keine Gottesfurcht an diesen Orten, und werden
\quellseite{170}
mich um meines Weibes willen erwürgen. Auch ist sie wahrhaftig meine Schwester, denn sie ist meines Vaters Tochter, aber nicht meiner Mutter Tochter, und ist mein Weib worden.
\end{schriftzitat}

1. In diesem Capitel sind etliche Stücke, die vor die Gelehrten gehören, die wollen wir weiter enhinter versparen. Es ist eine schlechte Historie anzusehen, daß GOtt der Allmächtige so seltsam mit Abraham umgeht, erhält sein Weib, und schafft sie ihm wieder, läßt sie ihm doch vorhin nehmen. Wenn es vor die Vernunft kommt, ist es eine schlechte Lection; da schmeckt es nichts, hat weder Saft noch Kraft, daß sie muß sagen: Ist das so nöthig gewesen zu beschreiben? Wie viel ist wohl mehr Jammers in der Welt, daß jemand um Mann, Weib und Kind schändlich kommt, das wohl erbärmlicher zu lesen wäre denn diese Geschichte!

2. Also habe ich aber vor auch gesagt, und sage noch, daß diese Historien darum gar hoch, über alle Historien zu heben und zu setzen sind, so sie geschrieben sind, nicht angesehen, daß es Menschen begangen haben, Abraham oder andere, sondern daß sie sind geschehen durch GOttes Wort und Werk. Darum werden wir noch mehr hernach hören so geringe Werke und närrisch Ding, daß [es] schier verdrießlich ist zu lesen, da er doch so viel Worte von macht, als wäre es das Allerköstlichste; alles darum, daß GOtt der Vernunft feind ist mit ihrem Dünkel, und [sie] meint, was sie groß achtet, das müsse GOtt auch groß achten.

3. Darum soll man darauf sehen, und alle Werke, die beschrieben sind, auf zwei Weisen scheiden. Einmal, ohne Geist gethan, die sind zuweilen hübsch, zuweilen greulich und schändlich, wie der Römer, Persen, Griechen und Anderer Historien, also, daß sie fast lustig zu lesen sind. Darauf fällt denn die Vernunft, und hält’s für köstlich, mißt es der Historie nach, wenn es nur ein groß, schön, tapfer Werk ist, fragt sie nicht weiter, wer es gethan hat; wiederum, ist es gering, so verachtet sie es, ob schon die Person groß ist.

4. Das andere Mal geschehen sie durch und im Geiste. Solche führt allein die Schrift, die sagt auch von großen Werken, richtet aber und urtheilt sie allein nach dem Glauben, darinne sie geschehen sind. Darum hanget darinne das köstliche Kleinod, GOttes Wort, welches man in der andern keinem\footnote{In den Ausgaben: keine.} findet. Derhalben scheinen solche Werke wohl geringe, aber mit geistlichen Augen angesehen sind sie trefflich, allein darum, daß sie gegangen sind in seinem Glauben, Gottes Wort und Heiligem Geist, davon ich oft genug gesagt habe.

5. Es ist aber noch noth, daß man es gar wohl bläue, bis man es fasse (denn es zu tief eingerissen, und den Leuten vorgebläuet ist, daß sie nicht mehr sehen, denn was groß ist), auf daß man recht müßte zu richten, was das Beste an einem Werke ist, nämlich, daß es geht in GOttes Geist, nicht angesehen, wie lang, groß, oder schwer es ist, oder wer es gethan hat. Schlage nur alles in Wind, und siehe hierher, ob es Gott gebeten hat, und im Glauben geht. Denn das ist der Schatz, Maß und Gewicht, darnach GOtt alle Werke mißt und wiegt; wiederum, geht es nicht im Glauben und GOttes Wort, so laß es gleich sein Todte aufwecken, so keusch und rein leben als die Mutter GOttes, so ist es vor GOtt nicht besser, denn öffentliche Sünde und Schande.

6. So wird hier der Mann Abraham groß gepreiset, und hoch gehoben von Abimelech, daß sich GOtt sein so hoch annimmt, und das ganze Land schlägt um Eines Mannes willen. Wie seltsam ist das, daß einem armen Mann Weib und Tochter entführt wird? da doch niemand von schreibt, wird auch nicht also gerochen, deß achtet niemand; und dies will er so groß geachtet haben. Wenn es einem andern widerführe, ohne Glauben und GOttes Wort, so nicht in GOttes Gebot geht, sondern nach der Vernunft fährt, wenn er noch so viel litte als alle Heiligen, achtete er es nicht.

7. Item, daß Abraham da im Lande irre zeucht, hat keine gewisse Stätte, bricht auf von Hebron mitten im Lande, zeucht gegen Mittag, zwischen Egypten; das Irrziehen muß auch beschrieben werden, daß man sich Mosis mit bekümmern. Wie viel sind ihrer aber, die auch irre ziehen, denen es doch herzlich sauer wird, müssen Weib und Kind, Leib und Gut in die Fahr setzen! Warum ist dasselbe nicht auch köstlich? Darum, Abraham ward also getrieben, daß er irre zog, nicht aus seinem Willen, es däuchte ihn nicht so gut, folgte nicht seinem Sinne, sondern GOtt hieß es ihn. Das Göt-
\quellseite{171}
tes Heißen ist die Güte und Adel des Werkes, deß kann GOtt nicht vergessen, [er] läßt es schreiben, so weit die Welt ist, eben darum, daß er nicht leiden will, daß man vor ihm ein Werk thue, das er nicht befohlen hat.

8. So siehe nun hier nicht das Werk, welches ja gering an ihm selbst ist, sondern Abrahams und auch seines Weibes Glauben und Herz an, wie er die Sara dahin gibt, und läßt es GOtt befohlen sein, wie er sie bei Ehren erhalte und ihm wieder schaffe, nachdem er die Verheißung hatte, daß er von ihr einen Sohn sollte haben. Denn solches kann keine Vernunft thun, daß sie solche Jahr frisch auf GOtt wage, und sich darinne zufrieden stelle. Solchen Schatz sieht niemand nicht, der Geist aber sieht es, darum streicht er es auch meisterlich aus. Wenn es ein anderer wäre, der so ein fromm Weib hätte, der letzte sein\footnote{„sein“ fehlt in der Erlanger.} Leib und Leben daran, Gut und Ehre, sie selbst zu retten und zu schützen; könnte er es nicht dahin bringen, so könnte er sich doch nicht brechen, daß er zufrieden wäre, und [es] GOtt heimstellte.

9. Also richte nun von einem jeglichen Werke, das da groß scheint, daß du mögest sagen: Aeußerlich lasse ich das Werk wohl größer sein denn Abrahams Werk; wenn man es aber nach dem Geiste wiegt, so ist jenes schwer, und überwiegt alle Werke auf Erden. Der Teufel mit den Verdammten müssen auch gar viel leiden; item, die Gottlosen noch hier auf Erden müssen zweierlei Unglück haben, daß ihnen viel saur wird, daß sie die Hölle verdienen, denn den Frommen, daß sie den Himmel verdienen. Siehe die Kriegsknechte an, welche Mühe haben sie, daß sie sich lassen würgen, wie sie gegen Schwert, Spieß und Büchsen treten, und weder Ehre noch Gut noch Leben achten. Welcher Heilige macht so viel um GOttes willen, als die um eines Guldens willen? Noch gelten solche Werke und Leiden alle nichts, denn sie haben das Gewicht nicht, das ist, GOttes Wort und Gefallen.

10. So schließen wir nun aus dieser Historie, daß wir uns vorsehen sollen vor allen Dingen, daß niemand in einigem Stande oder Wesen gehe, er wisse denn, daß [es] GOttes Wille ist, und könne sagen, daß gefällt ihm wohl. Im ehelichen Stande soll es auch also gehen. Wer da fühlt, daß er nicht Jungfrau kann sein, der hat seinen Beruf, daß er ehelich werde, und ist sicher, daß [es] GOtt so haben will, sintemal er so geschaffen ist. Was nun dergleichen ist, das sind eitel Werke, die sicher hergehen auf GOttes Wort, und wo sie nicht sicher sind, so taugen sie nichts. Das ist oft genugsam gesagt, ohne daß man sehe, wie die heilige Schrift uns überschüttet, gleichwie mit einem Platzregen, ist immer ein Exempel an andern, daß uns genug vermahne und stärke, nichts zu thun ohne sein Wort und Befehl. Das ist ein Stück von der Historie.

11. Zum andern ist sie auch geschrieben zu Trost allen, die an GOtt glauben und noch glauben sollen. Denn wenn sich GOtt nicht freundlich erzeigte, wer wollte oder könnte auf ihn trauen? Wenn er immer blitzte und donnerte, und Pestilenz ließe kommen, nicht auch Gutes thäte, würde ihm jedermann feind werden. Darum hat er so viel Exempel lassen schreiben, beide derer, die da glauben und nicht glauben, und sind alle Exempel dahin gerichtet, daß wir sehen, wie fleißig er Achtung auf uns hat, daß wir sicher und fest seien, er wolle uns nicht lassen. Das thut er wohl, stellt sich, als wollte er die Hand abziehen, und uns lassen stecken, aber er thut es nur, uns zu versuchen, ob wir fest am Glauben halten.

12. Also thut er hier auch. Was hätte es geschadet, daß er ohne die Fahr Abrahams Weib erhalten hätte, wie er wohl konnte, daß sie ihm nicht wäre genommen? Was gelüstet ihn denn, daß er sie in die Fahr setzt und ihn bekümmert macht, stellt sich eben, als wollte er ihm das Weib nehmen lassen? dazu, daß das Größte ist, daß er ihm verheißen hatte, er wollte ihm von Sara ein Kind geben; nach dem Troste sollte erst der Unfall kommen, daß er zweifeln möchte, ob er sein Wort halten würde? Er wäre ohne Zweifel gerne zu Hebron geblieben, hät nicht aus Vormuth gewandert, hat aber müssen aufbrechen um böser Leute willen; so läßt ihm GOtt in einem fremden Lande das Weib auch nehmen, daß er möchte sagen: Du hast mir zugesagt, mein Weib sollte mir über ein Jahr einen Sohn bringen, und nun lässest du sie hinweg nehmen: wo bleibt denn nun, das du gesagt hast?

13. Also gehen alle GOttes Werke. Er gibt uns die stärksten Verheißungen, so meinen wir denn, es soll gehen, wie wir denken; aber so spricht er: Ich will es über dein Denken und
\quellseite{172}
Vernunft machen, will dazu gebärden, als wollte ich lügen und eben das Widerspiel thun. So wird sich’s auch finden, wenn du sterben sollst, oder auch in Kampf treten über GOttes Wort: da wirst du fühlen, daß du es gerne anders haben wolltest. Warum thut er nun das? Darum, daß wir lernen rechtschaffen glauben; [er] will uns stecken lassen und die Hand abziehen, daß wir dennoch nicht zweifeln, sondern nur sicherer werden, er wolle uns nicht lassen, und uns also trösten: Ich weiß, daß er wird halten, was er gesagt hat.

14. Also hier. Wiewohl der König sein Weib holen läßt und zu sich nimmt, daß er nicht anders konnte denken, denn, es ist verloren, das Weib ist dahin; noch muß er so glauben: Sie bleibe da, oder komme hinweg, ja, wenn sie gleich jetzt vor meinen Augen stürbe, wollte ich dennoch nicht zweifeln, er werde mir ein Kind geben; es gehe mit Sara zu, wie es wolle, so kann er mir nicht lügen, ob sie schon todt dazu wäre. Er kann Ehre zu allen Schanden, und alle Schande zu Ehren machen, denn er ist allmächtig. Darum bleibt er im Glauben, läßt fahren, was da fährt; denkt, er sei wahrhaftig, befiehlt ihm die Weise, wie er es thun soll, welches keine Vernunft nimmermehr thun kann. Die denkt nicht weiter, denn also: Soll ich von ihr ein Kind haben, so müßte sie je bei mir bleiben; nun ist sie dahin, und wird mir zu Schanden. So sagt GOtt: Nicht also, es soll also geschehen, wie ich verheißen habe; wie dies aber geschehen soll, sollt du nicht errathen, ich will es machen, wie mich’s recht dünkt.

15. Siehe, das ist unser Trost, daß wir sehen, wie er erhält, die ihm mit einfältigen, reinen Herzen trauen, unter solchem Scheine und Ansehen, als sei er nicht da, und wir unserer Sachen keinen Rath wissen, und allein da stehen wider Teufel, Tod und alle Dinge, und dennoch sicher sind, es müsse wahr werden, das er geredet hat. Also thut er nun hier, läßt den König das Weib hinnehmen, nimmt sich aber ihrer und ihres Mannes also an, daß er die Sünde sucht und straft, die er nicht gethan hat, so greulich, als Abraham nicht hätte dürfen wünschen noch begehren, daß er nicht zufrieden ist, er habe ihm denn das Weib wieder gegeben.

16. Das sind eitel GOttes That und Rath, wie wir uns in solchen Versuchungen sollen halten. Denn damit hat er angezeigt, was seine Art ist, und wie er gesinnt ist gegen den Gläubigen, daß er nicht leiden kann, was man an ihnen thut; darum er auch (Sach. 2, 8.) sagt: „Wer euch anrühret, der rühret meinen Augapfel an“; allein, daß wir nichts thun, wir fragen ihn denn darum. Trotz, daß er ihm das Weib angerührt hätte; [er] hatte sie bei sich im Hause, noch konnte er nicht zu ihr. Daher abermal im Psalm steht [Ps. 105, 14.]: „Er ließ keinen Menschen sie beleidigen, und strafte Könige um ihretwillen.“

17. Also thue du nun auch. Laß uns fahren, wohin wir fahren; laß kommen und weggehen, was da kommt und geht, so ist er bei uns, und will beweisen, daß er die Seinen halten könne, als ein allmächtiger HErr. Dies ist das Exempel dieses Capitels, uns zu Trost und Unterweisung des Glaubens geschrieben.

18. Nun sind auch etliche Fragen hier zu handeln. Die erste ist davon, daß dieser Abimelech sich rühmt vor GOtt, er habe das gethan mit einfältigem Herzen. Und spricht weiter: „Willst du denn auch ein gerecht Volk erwürgen?“ Was ein einfältig Herz sei, lasse ich gehen. Das unser lateinischer Text heißt cor perfectum, ein vollkommenes Herz, ist nicht zu verstehen nach der Sophisten Weise, wie sie träumen von der Vollkommenheit. Was wollen wir aber dazu sagen, daß GOtt selbst spricht: „Ich weiß auch, daß du mit einfältigem Herzen (das ist, mit rechtschaffenem Herzen und guter Meinung) das gethan hast, darum habe ich dich erhalten, daß du nicht sündigest.“

19. Das ist je ein groß Lob und Argument, das Papstthum zu stärken wider uns, für den freien Willen. Da sagen sie: Der Abimelech war ein heidnischer König, that da Sünde, und war doch nicht Sünde, und wird gestraft, und GOtt gibt ihm dennoch Zeugniß, daß er eine rechtschaffene Meinung habe: so folgt aus diesen Texten, daß außer der Sünde auch die Heiden mögen fromm sein vor GOtt, wie er hier sagt, nicht allein vor den Leuten; so müssen wir lassen stehen, daß er unschuldig in dem Stücke sei, auch vor Gott.

20. Antwort: Wir können nicht gewiß sein aus der Schrift, daß Abimelech sei GOttes Kind und ein rechter Heiliger gewesen; es sieht ihm aber fast ähnlich. Aber antworte also: Der Grund, den wir oft gesagt haben, den siehst du, wie er durch alle Schrift geht, nämlich, daß ohne
\quellseite{173}
Jesum Christum keine Hülfe noch Gnade ist, der hat uns umsonst geholfen; weil wir alle Sünder sind, kostet es sein eigen Blut und Leben. Das ist unser Grund. Was nun für Einreden damit geführt werden, kannst du sie nicht verantworten, so schweig stille; [es] ist besser, denn daß du darum solchen klaren und gewissen Grund leugnen solltest.

21. Denn das soll man nicht leugnen, daß auch oft rechte fromme Christen unter den Heiden gewesen sind, wie wir auch [Dan. 3, 26.] lesen von dem Könige Nabuchodonosor, item, Naeman von Syrien [2 Kön. 5, 17.], welche auch fromme Leute waren, und viel mehr unter den Heiden. Darum ist es glaublich, daß dieser auch gläubig gewesen sei. Denn lasse es nicht einen Schimpf sein, wo GOtt solche Gnade thut, wie diesem Könige, daß er zu ihm kommt und sich offenbart, mit ihm redet, warnt ihn vor Schaden. Es muß je nicht ein geringer Mann sein. Die Heiden läßt er hingehen in ihrer Vernunft; zu welchem er sich aber so nahe zuthut, ist es ein Zeichen, daß er ein Auge auf ihn hat, und vor ihm angesehen sei. Man muß die Heiden nicht also verwerfen, die um die Juden her gewohnt haben, daß etliche auch fromm gewesen sind, als die Königin von Saba [1 Kön. 10, 9.], und Hiram, der König von Tyro [1 Kön. 5, 7.].

22. Darum will ich also antworten: daß Abimelech ein recht frommer Mann ist gewesen, und habe GOtt wohl gekannt und recht gedient. Aber, wie GOtt fährt mit andern Heiligen, daß er sie zuweilen straucheln läßt, als wir von Lot gehört haben, so geht es dem Könige auch. Er ist aber dennoch so redlich, wenn er gewußt hätte, daß [sie] ein Eheweib wäre, würde er sie nicht zu sich genommen haben. Welches Stück auch eben das glaublich anzeigt, daß er ein frommer König sei gewesen, allein damit gesündigt, wie andere Heilige, und GOtt auch ihn gnädiglich straft, wie andere Heilige. Denn die Bösen straft er in die Hölle.

23. Gefällt aber jemand dies nicht, der sage wie St. Augustinus, daß GOtt nicht weiter Zeugniß gibt denn auf das Werk. Denn es möchte wohl sein, daß er in andern Stücken böse gewesen sei, aber in dem nicht. Das aber Gott solches zeugt, damit wäre nicht gesagt, daß er darum ein rechtschaffen rein Herz habe vor ihm, sondern nur in diesem Falle und Werke; also, daß, ob er gleich ein Bube ist, kann er sich dennoch rühmen, daß er dies nicht gethan habe, mit gutem Gewissen. Also legt es St. Augustinus aus. Ich aber wollte bei der ersten Antwort bleiben, auf daß man GOttes Gnade auch unter den Heiden lasse gehen, und halte es für groß Wohlthat und Gnade, daß ihn GOtt warnt und selbst erscheint. Doch, wie gesagt, lasse ich hier einem jeglichen sein Gutdünken.

24. Die andere Frage ist: Wie Abraham sagt von seinem Weibe: „Auch ist sie wahrhaftig meine Schwester, denn sie ist meines Vaters Tochter, aber nicht meiner Mutter Tochter, und ist mein Weib worden“? Das meint er also, daß sie nicht seine leibliche Schwester sei, sondern seine Stiefschwester. Das redet Moses mit umschweifenden Worten, als sollte er sagen: „Mein Vater hat sie nicht gezeuget, ist nicht von der Mutter Leibe, es sei Stiefmutter oder rechte Mutter; darum will er sagen, daß sie nicht der Geburt halben Schwester ist, sondern ein angenommenes Kind, die mein Vater zu sich genommen hat für seine Tochter, also, daß sie Abrahams Vater auch Vater nennt, ohne daß er sie nicht gezeuget hat. Nun haben wir droben [Cap. 11, § 32] gehört, daß sie Abrahams Bruders Tochter gewesen ist.

25. Denn die Schrift hat die Glieder nicht also gespannt, als der Papist, der nicht zuläßt, daß einer ein Weib nehme im vierten Gliede, im dritten viel weniger, doch so gnädig ist, wenn man Geld brächte, ließe er es geschehen. Es heißen aber Buhenstücke, wenn man Frauen-Fleisch feil trägt, und verkauft es um Geld. Aber nach der Schrift mag ich wohl nehmen meine Muhme, herunter zu rechnen, nicht hinaufwärts, das ist, meines Bruders oder Schwester Tochter, aber nicht des Vaters oder Mutter Bruder noch Schwester, wiewohl sie beide in gleichem Gliede sind, Muhmen oder Basen, junge und alte. Das hat GOtt also geordnet, daß er mehr Ehre will haben gegen den Alten denn gegen den Jungen, daß das junge Volk vom alten unterwiesen und gelehrt werde. Das rede ich darum, daß man wisse, wie weit man nach der Schrift freien möge. Es wäre nicht noth, daß man es verwirrte, es verwirrt sich selbst mehr, denn man richten kann. Wenn nun in dem Falle ein Paar zusammen wollten, oder gegeben
\quellseite{174}
wären, möchte es mit GOtt und gutem Gewissen geschehen, wäre auch schier das Beste. Ich halte nicht allzuviel von weiter Freundschaft, die sich so weit weg flieht; [es] ist besser, daß die Freundschaft und das Gut nahe bei einander bleibe; wie es unter den Juden geschah; doch ist uns genug, daß wir aus der Schrift wissen, die Gewissen zu unterrichten. So spricht nun Abraham weiter:

\begin{schriftzitat}
V. 13. Da mich aber GOtt aus meines Vaters Hause wandern hieß, sprach ich zu ihr: Die Barmherzigkeit thue an mir, daß, wo wir hinkommen, du von mir sagest, ich sei dein Bruder.
\end{schriftzitat}

26. Das ist auf hebräisch geredet, welche sehr eine göttliche, freundliche Sprache ist, gibt der Sache feine Namen, daß [es] Lust ist, hat die hübschesten, schönsten Worte der Liebe und Erkenntniß. Also heißt nun ein gut Werk, das wir Freundschaft nennen, auf ihre Sprache „Barmherzigkeit“. So wirst du es durch und durch finden; dafür wir sagen: Thue mir die Freundschaft und Liebe, oder: Thue mir so viel zugute. Daher heißt nun die Gnade, die GOtt in unser Herz gießt und uns fromm macht, auch auf ihre Sprache [\hebraeisch{חֶסֶד}] Chesed, Barmherzigkeit; darum, daß, wie man ein Werk frei umsonst, aus lauter Freundschaft thut, nicht daß man es pflichtig sei. So ist es auch mit GOtt, was er gibt, ist nichts denn Barmherzigkeit, Gabe, Geschenk, lauter Freundschaft, und kommt aus freier, bloßer Güte, Gunst und Liebe, ohne und wider unser Verdienst. Solcher Sprache müssen wir gewohnen im Alten Testamente. Noch eines ist hier, das der Text spricht:

\begin{schriftzitat}
V. 14. 15. Da nahm Abimelech Schafe und Rinder, Knechte und Mägde, und gab sie Abraham, und gab ihm wieder sein Weib Sara, und sprach: Siehe da, mein Land steht dir offen, wohne, wo dir’s wohlgefällt.
\end{schriftzitat}

27. Es muß je ein frommer König gewesen sein und Gott wohl erkannt haben, ohne daß er, wie gesagt, hier gestrauchelt hat. Wiewohl er nicht darauf bleibt, der gute Mann, oder sein Sohn, wie wir hernach sehen werden. Weiter folgt:

\begin{schriftzitat}
V. 16—18. Und sprach zu Sara: Siehe da, ich habe deinem Bruder tausend Silberlinge gegeben. Siehe, der\footnote{Wittenberger und Jenaer: der; Erlanger: er. Vergleich § 33.} soll dir ein Deckel der Augen sein vor allen, die bei dir sind, und allenthalben, und ein Verarbeiter. Abraham aber betete zu GOtt; da heilte GOtt Abimelech, und sein Weib, und seine Mägde, daß sie Kinder gebaren. Denn der HErr hatte zuvor hart verschlossen alle Mütter des Hauses Abimelech, um Sara, Abrahams Weibes willen.
\end{schriftzitat}


28. Da nennt er ihn auch selbst ihren Bruder. Es sollte auch so sein unter den Ehelichen, und viel freundlicher denn Brüder und Schwestern, es ist Ein Leib und zwei Seelen. Die tausend Silberlinge aber, nach dem jüdischen Seckel zu rechnen (welcher fast so viel gewesen ist, als bei uns ein Ort eines Gulbens), derselbigen tausend machen fast drittehalb hundert Gulden; ist ein königliches Geschenk. Das hat er ihr gegeben über die Schafe, Kinder, Knechte und Mägde, die sind auch alles leibeigene Güter, wie ander Vieh, daß sie die verkauften, wie sie wollten. Wie noch schier das Beste wäre, daß es noch wäre, [es] kann doch sonst das Gesinde niemand zwingen noch zähmen.

29. Und hier siehst du, daß dieser Abraham und Abimelech das halten, und das Gesinde bleiben lassen in dem Dienste. Es wäre wohl große Freundschaft und Barmherzigkeit gewesen, daß er sie hätte frei gelassen. Wie konnte die Liebe leiden, daß sie es so gehalten haben? Eben wie sie leiden kann, daß man die Leute an Galgen hängt, oder sonst richtet; denn man muß das leibliche, weltliche Regiment auch halten im Schwange, daß man die Leute zähme und zwinge.

30. Also sind sie auch damit umgegangen. Ihret halben hätten sie es wohl lassen gehen; [es] wäre aber nicht gut, [sie] sollten bald zu stolz worden sein, wenn man ihnen so viel Recht gäbe, oder hielte sie als sich selbst oder ein Kind. Man muß einen jeglichen halten in seinem Stande, wie es GOtt ordnet, Sohn, Tochter, Knecht, Magd, Mann, Weib \&c. Er will es nicht gebrochen haben; [es] kann auch niemand anders das Volk im Zaume halten, denn mit dem Zwange äußerlichen Regimentes.

31. Daher so große Klage ist über Gesinde und Dienstleute in der Welt. Es ist des Teufels und Papsts Schuld, und der Fürsten, daß kein Regiment ist, es thut jedermann, was er will. Wäre aber die Faust und Zwang da, wie die Zeit gewesen ist, daß niemand mucken dürfte,
